% Descriptive inorganic Chemistry: Halogens, Group IV and d-block elements — Chemistry Upper Sixth — Cours signé OnBuch+
% Official MINESEC syllabus. Compiler avec : tectonic CHE03-descriptive-inorganic-chemistry-halogens-group-iv-and-d-bloc.tex
\newcommand{\DOCMATIERE}{Chemistry}
\newcommand{\DOCNIVEAU}{Upper Sixth}
\newcommand{\DOCMODULE}{Upper Sixth — Chemistry}
\newcommand{\DOCLECON}{3}
\newcommand{\DOCTITRE}{Descriptive inorganic Chemistry: Halogens, Group IV and d-block elements}
% OnBuch+ — préambule « cours signés » ENRICHI (compilé avec tectonic / XeLaTeX)
% Système visuel « Pop » d'OnBuch+ : fond crème (jamais blanc), bordures encre
% épaisses, ombres « dures » décalées, titres Archivo Black en capitales.
% Version MATHÉMATIQUES : graphes (pgfplots/tikz), boîtes pédagogiques
% (définition, propriété, méthode, attention, le savais-tu, exercice résolu,
% exercice, TP, bilan), page de garde et signature OnBuch+ ancrée en bas.
%
% Macros attendues dans main.tex AVANT \input{preamble} :
%   \DOCMATIERE  \DOCNIVEAU  \DOCMODULE  \DOCLECON (numéro)  \DOCTITRE
\documentclass[11pt]{article}

\usepackage{fontspec}
\usepackage[english]{babel}
\usepackage[a4paper,margin=2cm,top=3.4cm,bottom=2.8cm,headheight=1.4cm,footskip=1.3cm]{geometry}
\usepackage{amsmath,amssymb,amsfonts}
\usepackage{mathtools} % \xmapsto, \coloneqq... souvent utilisés par les modèles
\usepackage{cancel} % simplifications barrées (bilans de forces)
% \abs{z} (module, valeur absolue) et \norm{u} (norme), utilisés par les modèles
\providecommand{\abs}{}\let\abs\relax\DeclarePairedDelimiter{\abs}{\lvert}{\rvert}
\providecommand{\norm}{}\let\norm\relax\DeclarePairedDelimiter{\norm}{\lVert}{\rVert}
\usepackage[table]{xcolor} % [table] : fournit aussi \rowcolors (alternance de lignes)
\usepackage{pagecolor}
\usepackage{tikz}
\usetikzlibrary{arrows.meta,positioning,calc,shapes.geometric,shapes.misc,shapes.arrows,decorations.pathmorphing,decorations.pathreplacing,decorations.markings,patterns,shadows,mindmap,trees,fit,backgrounds,matrix,intersections,angles,quotes}
\usepackage{pgfplots}
\usepackage[version=4]{mhchem} % équations nucléaires \ce{^{235}_{92}U}
\usepackage[european,siunitx]{circuitikz} % schémas électriques
\pgfplotsset{compat=1.18}
\usepgfplotslibrary{fillbetween,groupplots} % aires entre courbes (calcul intégral)
\usepackage{enumitem}
\usepackage{fancyhdr}
% [most] charge toutes les bibliothèques tcolorbox (listings, minted, raster,
% documentation...) et gonfle inutilement la mémoire TeX sur un document long
% (~300 boîtes) : on ne charge que ce qui est réellement utilisé.
\usepackage[skins,breakable]{tcolorbox}
\usepackage{graphicx}
\usepackage{colortbl}
\usepackage{booktabs}
\usepackage{tabularx}
\usepackage{diagbox} % case d'en-tête diagonale des tableaux à double entrée
% Colonnes extensibles centrée / gauche / droite (C, L, R), souvent utilisées par les modèles
\newcolumntype{C}{>{\centering\arraybackslash}X}
\newcolumntype{L}{>{\raggedright\arraybackslash}X}
\newcolumntype{R}{>{\raggedleft\arraybackslash}X}
\usepackage{array}
\usepackage{multicol}
\usepackage{siunitx}
% parse-numbers=false : accepte aussi \SI{2,5 \times 10^{-3}}{...}
\sisetup{locale=UK,output-decimal-marker={.},group-minimum-digits=4,parse-numbers=false}
\DeclareSIUnit\ohm{\ensuremath{\symup{\Omega}}} % U+2126 absent de la police
\DeclareSIUnit\division{div} % réglages d'oscilloscope (ms/div, V/div)
\DeclareSIUnit\tr{tr} % tours (vitesse de rotation, ex. \SI{1500}{\tr\per\minute})
\usepackage{qrcode}
% \degree : commande fréquemment utilisée par les modèles pour un angle ou une
% température en dehors de \SI{}{} (non fournie par les paquets chargés).
\providecommand{\degree}{^{\circ}}
% \qed (fin de démonstration), utilisé par les modèles sans amsthm
\providecommand{\qed}{\hfill$\square$}
% \barre : double barre des valeurs interdites dans un tableau de variations (tkz-tab)
\providecommand{\barre}{\ensuremath{\|}}
% environnement proof (amsthm non chargé) utilisé par les modèles
\ifdefined\proof\else\newenvironment{proof}[1][Proof]{\par\noindent\textit{#1.}\ }{\qed\par}\fi
\usepackage{lastpage}
\usepackage{hyperref}
\hypersetup{hidelinks}

% ============================================================
% Palette « Pop » OnBuch+ — mêmes hex que la classe P (plus_ui.dart)
% ============================================================
\definecolor{poporange}{HTML}{F59321}
\definecolor{popdark}{HTML}{E07A0C}
\definecolor{popcream}{HTML}{FFF6E8}
\definecolor{popink}{HTML}{1C1714}
\definecolor{poppurple}{HTML}{7A5AE0}
\definecolor{popgreen}{HTML}{2E9E6A}
\definecolor{poppink}{HTML}{E0457B}
\definecolor{popblue}{HTML}{2F7DE1}
\definecolor{popgold}{HTML}{C98A12}
\definecolor{popmuted}{HTML}{8A7F76}
\definecolor{popline}{HTML}{EADFD2}
\definecolor{popgray}{HTML}{8A7F76}
\colorlet{popgrey}{popgray}
% Alias tolérants (le modèle nomme parfois les couleurs différemment)
\colorlet{popwhite}{white}
\colorlet{popblack}{popink}
\colorlet{popred}{poppink}
\colorlet{popyellow}{popgold}
% Teintes claires pour fonds de tableaux / zones
\colorlet{popblueL}{popblue!10!white}
\colorlet{poporangeL}{poporange!14!white}
\colorlet{popgreenL}{popgreen!12!white}
\colorlet{poppurpleL}{poppurple!12!white}
\colorlet{poppinkL}{poppink!12!white}
\colorlet{popgoldL}{popgold!14!white}

\pagecolor{popcream}

% ============================================================
% Typographie
% ============================================================
\defaultfontfeatures{Ligatures=TeX}
\setmainfont{PlusJakartaSans-Medium.ttf}[Path=../../../fonts/,BoldFont=PlusJakartaSans-Bold.ttf]
% Maths en sans-serif (Fira Math) avec chiffres et lettres droites de Plus
% Jakarta, pour que les formules se fondent dans le texte.
\usepackage[math-style=ISO,bold-style=ISO]{unicode-math}
\setmathfont{FiraMath-Regular.otf}[Path=../../../fonts/]
\setmathfont{PlusJakartaSans-Medium.ttf}[Path=../../../fonts/,range={up/num,up/{Latin,latin},\mathup}]
% (icomma retiré : décimales avec point en anglais)
% Caractères mathématiques Unicode absents des polices de texte (Plus Jakarta,
% Archivo Black) : rendus par leur équivalent mathématique, en texte comme en
% formule — sinon ils s'affichent en carré vide (ex. « Arithmétique dans ℤ »).
\usepackage{newunicodechar}
\newunicodechar{ℝ}{\ensuremath{\mathbb{R}}}
\newunicodechar{ℤ}{\ensuremath{\mathbb{Z}}}
\newunicodechar{ℂ}{\ensuremath{\mathbb{C}}}
\newunicodechar{ℕ}{\ensuremath{\mathbb{N}}}
\newunicodechar{ℚ}{\ensuremath{\mathbb{Q}}}
\newunicodechar{𝔻}{\ensuremath{\mathbb{D}}}
\newunicodechar{α}{\ensuremath{\alpha}}
\newunicodechar{β}{\ensuremath{\beta}}
\newunicodechar{γ}{\ensuremath{\gamma}}
\newunicodechar{δ}{\ensuremath{\delta}}
\newunicodechar{ε}{\ensuremath{\varepsilon}}
\newunicodechar{θ}{\ensuremath{\theta}}
\newunicodechar{λ}{\ensuremath{\lambda}}
\newunicodechar{μ}{\ensuremath{\mu}}
\newunicodechar{σ}{\ensuremath{\sigma}}
\newunicodechar{τ}{\ensuremath{\tau}}
\newunicodechar{φ}{\ensuremath{\varphi}}
\newunicodechar{ψ}{\ensuremath{\psi}}
\newunicodechar{ω}{\ensuremath{\omega}}
\newunicodechar{Σ}{\ensuremath{\Sigma}}
% majuscules produites par \MakeUppercase dans les titres (α-aminés -> Α)
\newunicodechar{Α}{\ensuremath{\alpha}}
\newunicodechar{Β}{\ensuremath{\beta}}
\newunicodechar{Γ}{\ensuremath{\Gamma}}
\newunicodechar{Δ}{\ensuremath{\Delta}}
\newunicodechar{Ε}{\ensuremath{\varepsilon}}
\newunicodechar{Θ}{\ensuremath{\Theta}}
\newunicodechar{Λ}{\ensuremath{\Lambda}}
\newunicodechar{Μ}{\ensuremath{\mu}}
\newunicodechar{Τ}{\ensuremath{\tau}}
\newunicodechar{Φ}{\ensuremath{\Phi}}
\newunicodechar{Ψ}{\ensuremath{\Psi}}
\newunicodechar{Ω}{\ensuremath{\Omega}}
\newunicodechar{↦}{\ensuremath{\mapsto}}
\newunicodechar{∈}{\ensuremath{\in}}
\newunicodechar{∉}{\ensuremath{\notin}}
\newunicodechar{∧}{\ensuremath{\wedge}}
\newunicodechar{⇔}{\ensuremath{\Leftrightarrow}}
\newunicodechar{⟺}{\ensuremath{\Longleftrightarrow}}
\newunicodechar{⇒}{\ensuremath{\Rightarrow}}
\newunicodechar{⟹}{\ensuremath{\Longrightarrow}}
\newunicodechar{∘}{\ensuremath{\circ}}
\newunicodechar{∪}{\ensuremath{\cup}}
\newunicodechar{∩}{\ensuremath{\cap}}
\newunicodechar{≡}{\ensuremath{\equiv}}
\newunicodechar{⊂}{\ensuremath{\subset}}
\newunicodechar{⊥}{\ensuremath{\perp}}
\newunicodechar{∃}{\ensuremath{\exists}}
\newunicodechar{∀}{\ensuremath{\forall}}
\newunicodechar{∛}{\ensuremath{\sqrt[3]{\,}}}
\newunicodechar{∜}{\ensuremath{\sqrt[4]{\,}}}
\newunicodechar{⃗}{\ensuremath{{}^{\to}}}
\newunicodechar{ⁿ}{\ifmmode^{n}\else\textsuperscript{\ensuremath{n}}\fi}
\newunicodechar{ˣ}{\ifmmode^{x}\else\textsuperscript{\ensuremath{x}}\fi}
\newunicodechar{ᵇ}{\ifmmode^{b}\else\textsuperscript{\ensuremath{b}}\fi}
\newunicodechar{ʳ}{\ifmmode^{r}\else\textsuperscript{\ensuremath{r}}\fi}
\newunicodechar{ᵘ}{\ifmmode^{u}\else\textsuperscript{\ensuremath{u}}\fi}
\newunicodechar{ʸ}{\ifmmode^{y}\else\textsuperscript{\ensuremath{y}}\fi}
\newunicodechar{ᵃ}{\ifmmode^{a}\else\textsuperscript{\ensuremath{a}}\fi}
\newunicodechar{ᵉ}{\ifmmode^{e}\else\textsuperscript{\ensuremath{e}}\fi}
\newunicodechar{ᵏ}{\ifmmode^{k}\else\textsuperscript{\ensuremath{k}}\fi}
\newunicodechar{ᵖ}{\ifmmode^{p}\else\textsuperscript{\ensuremath{p}}\fi}
\newunicodechar{ᴺ}{\ifmmode^{N}\else\textsuperscript{\ensuremath{N}}\fi}
\newunicodechar{ᶜ}{\ifmmode^{c}\else\textsuperscript{\ensuremath{c}}\fi}
\newunicodechar{ʰ}{\ifmmode^{h}\else\textsuperscript{\ensuremath{h}}\fi}
\newunicodechar{ᵠ}{\ifmmode^{q}\else\textsuperscript{\ensuremath{q}}\fi}
\newunicodechar{ₙ}{\ifmmode_{n}\else\textsubscript{\ensuremath{n}}\fi}
\newunicodechar{ₐ}{\ifmmode_{a}\else\textsubscript{\ensuremath{a}}\fi}
\newunicodechar{ᵢ}{\ifmmode_{i}\else\textsubscript{\ensuremath{i}}\fi}
\newunicodechar{ᵣ}{\ifmmode_{r}\else\textsubscript{\ensuremath{r}}\fi}
\newunicodechar{ₖ}{\ifmmode_{k}\else\textsubscript{\ensuremath{k}}\fi}
\newunicodechar{ᵦ}{\ifmmode_{\beta}\else\textsubscript{\ensuremath{\beta}}\fi}
\newunicodechar{ₓ}{\ifmmode_{x}\else\textsubscript{\ensuremath{x}}\fi}
\newfontfamily\popdisplay{ArchivoBlack-Regular.ttf}[Path=../../../fonts/]
\newfontfamily\poplogo{SpaceGrotesk-Bold.ttf}[Path=../../../fonts/]
\newfontfamily\popsemi{PlusJakartaSans-SemiBold.ttf}[Path=../../../fonts/]
\newfontfamily\popxbold{PlusJakartaSans-ExtraBold.ttf}[Path=../../../fonts/]
\newfontfamily\popxboldls{PlusJakartaSans-ExtraBold.ttf}[Path=../../../fonts/,LetterSpace=3.0]

\setlist[enumerate]{leftmargin=1.7em,itemsep=5pt,topsep=4pt}
\setlist[itemize]{leftmargin=1.7em,itemsep=4pt,topsep=4pt,label={\color{poporange}\textbullet}}
\setlength{\parindent}{0pt}
\setlength{\parskip}{5pt}
\renewcommand{\arraystretch}{1.25}

% pgfplots : style Pop par défaut
\pgfplotsset{
  every axis/.append style={axis line style={popink,line width=1pt},tick style={popink},
    grid style={popline,line width=0.5pt},label style={font=\small},tick label style={font=\footnotesize}},
  popcurve/.style={poporange,line width=1.8pt,no markers,smooth},
}
% Même style disponible dans un \draw TikZ simple (hors pgfplots)
\tikzset{popcurve/.style={poporange,line width=1.8pt}}
\tikzset{popgrid/.style={popline,very thin}}
\colorlet{popcurve}{poporange} % le nom du style est parfois utilisé comme couleur

% ============================================================
% Pill / squiggle
% ============================================================
\newcommand{\poppill}[3]{%
  \tikz[baseline=-0.55ex]\node[draw=popink,line width=1.2pt,rounded corners=2.6pt,%
    fill=#2,inner xsep=6.5pt,inner ysep=3pt]{%
    \popxboldls\fontsize{7}{7}\selectfont\color{#3}\MakeUppercase{#1}};%
}
\newcommand{\popsquiggle}[1]{%
  \tikz[baseline=-0.4ex]{
    \draw[#1,line width=2.6pt,line cap=round]
      plot [smooth] coordinates {(0,0.09) (0.34,0.01) (0.68,0.21) (1.02,0.01) (1.36,0.10)};
  }%
}
% Mot-clé surligné (marqueur orange sous le texte)
\newcommand{\cle}[1]{{\popxbold\color{popdark}#1}}

% ============================================================
% En-tête / pied
% ============================================================
\pagestyle{fancy}
\fancyhf{}
\renewcommand{\headrulewidth}{0pt}
\renewcommand{\footrulewidth}{0pt}
\lhead{\raisebox{-0.15ex}{\poplogo\fontsize{13}{13}\selectfont\color{popink}OnBuch\color{poporange}+}}
\rhead{\poppill{\DOCMATIERE\ \textbullet\ \DOCNIVEAU\ \textbullet\ Chapter \DOCLECON}{white}{popink}}
\lfoot{\popxbold\fontsize{7.6}{7.6}\selectfont\color{popmuted}\MakeUppercase{Signed course OnBuch+}\ \textbullet\ {\popsemi\DOCTITRE}}
\rfoot{\tikz[baseline=-0.6ex]\node[rounded corners=8.5pt,fill=popink,minimum height=17pt,minimum width=17pt,inner xsep=5pt,inner sep=0]{\popxbold\fontsize{8}{8}\selectfont\color{popcream}\thepage\,/\,\pageref*{LastPage}};}

% ============================================================
% Titres
% ============================================================
\newcommand{\chaptitre}[2]{%
  \begin{tcolorbox}[enhanced,colback=poporange,colframe=popink,boxrule=2.6pt,arc=11pt,
      drop shadow=popink,left=15pt,right=15pt,top=12pt,bottom=11pt,before skip=3pt,after skip=18pt]
    \poppill{#2}{popink}{popcream}\par\vspace{9pt}
    {\raggedright\hyphenpenalty=10000\exhyphenpenalty=10000\popdisplay\fontsize{22}{24}\selectfont\color{popcream}\MakeUppercase{#1}\par}
  \end{tcolorbox}
}
% Section numérotée : pastille orange avec le numéro + titre Archivo
\newcounter{popsec}
\newcommand{\coursec}[1]{%
  \stepcounter{popsec}\setcounter{popsub}{0}%
  \par\vspace{16pt}\noindent
  \begin{minipage}{\linewidth}
  \tikz[baseline=-0.7ex]\node[draw=popink,line width=1.6pt,fill=poporange,rounded corners=4pt,minimum size=22pt,inner sep=0]{\popdisplay\fontsize{12}{12}\selectfont\color{popcream}\thepopsec};%
  \hspace{8pt}{\popdisplay\fontsize{15}{17}\selectfont\color{popink}\MakeUppercase{#1}}\par
  \vspace{4pt}\noindent{\color{poporange}\rule{36pt}{3.6pt}}
  \end{minipage}\par\nopagebreak\vspace{8pt}
}
\newcounter{popsub}
\newcommand{\courssub}[1]{%
  \stepcounter{popsub}%
  \par\vspace{9pt}\noindent{\popxbold\fontsize{11.5}{13}\selectfont\color{popink}{\color{poporange}\thepopsec.\thepopsub}\hspace{6pt}#1}\par\nopagebreak\vspace{4pt}
}

% ============================================================
% Boîtes pédagogiques — même recette : fond blanc, bordure épaisse colorée,
% ombre dure encre, badge plein. Titre optionnel [#1].
% ============================================================
\tcbset{popbase/.style={breakable,enhanced,colback=white,boxrule=2.3pt,arc=9pt,
  shadow={3pt}{-3pt}{0pt}{popink},left=14pt,right=14pt,top=11pt,bottom=11pt,before skip=11pt,after skip=15pt,
  fontupper=\popsemi\fontsize{10.6}{14.5}\selectfont\color{popink},
  fontlower=\popsemi\fontsize{10.6}{14.5}\selectfont\color{popink}}}
% Coche dessinée (le glyphe ✓ n'existe pas dans les polices du document)
\newcommand{\popcheck}[1]{\tikz[baseline=-0.5ex]\draw[#1,line width=1.6pt,line cap=round,line join=round](-0.13,0.0)--(-0.04,-0.09)--(0.13,0.1);}
\providecommand{\coche}{\popcheck{popgreen}}
\providecommand{\resultat}[1]{\par\smallskip\noindent\fcolorbox{popgreen}{popgreenL}{\parbox{\dimexpr\linewidth-2\fboxsep-2\fboxrule\relax}{\textbf{Result.} #1}}\par\smallskip}
\newcommand{\popboxhead}[3]{\poppill{#1}{#2}{white}\ifx\relax#3\relax\else\hspace{6pt}{\popxbold\fontsize{10.6}{12}\selectfont\color{popink}#3}\fi\par\vspace{7pt}}

\newtcolorbox{aretenir}[1][]{popbase,colframe=popgreen,before upper={\popboxhead{Key points}{popgreen}{#1}}}
\newtcolorbox{exemplebox}[1][]{popbase,colframe=poporange,before upper={\popboxhead{Exemple}{poporange}{#1}}}
\newtcolorbox{definition}[1][]{popbase,colframe=popblue,before upper={\popboxhead{Definition}{popblue}{#1}}}
\newtcolorbox{propriete}[1][]{popbase,colframe=poppurple,before upper={\popboxhead{Property}{poppurple}{#1}}}
\newtcolorbox{methode}[1][]{popbase,colframe=popgold,before upper={\popboxhead{Method}{popgold}{#1}}}
\newtcolorbox{attention}[1][]{popbase,colframe=poppink,before upper={\popboxhead{Warning}{poppink}{#1}}}
\newtcolorbox{experience}[1][]{popbase,colframe=popblue,colback=popblueL,before upper={\popboxhead{Experiment}{popblue}{#1}}}
% « Le savais-tu ? » : carte violette pleine (contexte, culture, Cameroun)
\newtcolorbox{savaistu}[1][]{popbase,colback=poppurple,colframe=popink,
  fontupper=\popsemi\fontsize{10.4}{14.2}\selectfont\color{white},
  before upper={\poppill{Did you know?}{white}{poppurple}\ifx\relax#1\relax\else\hspace{6pt}{\popxbold\color{white}#1}\fi\par\vspace{7pt}}}
% Exercice résolu : énoncé en haut, \tcblower puis la solution sur fond crème
\newcounter{popexo}\newcounter{popexr}
\newtcolorbox{exoresolu}[1][]{popbase,colframe=poporange,colbacklower=popcream,
  segmentation style={popink,line width=1.2pt,dashed},
  before upper={\stepcounter{popexr}\popboxhead{Worked example \thepopexr}{poporange}{#1}},
  before lower={\poppill{Solution}{popink}{popcream}\par\vspace{6pt}}}
% Exercice (non résolu en place) — la correction est dans la partie Corrigés
\newtcolorbox{exercice}[1][]{popbase,colframe=popink,
  before upper={\stepcounter{popexo}\popboxhead{Exercise \thepopexo}{popink}{#1}}}
% Corrigé
\newtcolorbox{corrige}[1][]{popbase,colframe=popgreen,colback=popgreenL,
  before upper={\popboxhead{Solution}{popgreen}{#1}}}
% Objectifs / prérequis / bilan
\newtcolorbox{objectifs}{popbase,colframe=popink,colback=poporangeL,
  before upper={\popboxhead{Chapter objectives}{popink}{}}}
\newtcolorbox{prerequis}{popbase,colframe=popink,colback=popblueL,
  before upper={\popboxhead{Prerequisites}{popblue}{}}}
\newtcolorbox{bilan}{popbase,colframe=popink,colback=white,boxrule=2.8pt,
  before upper={\popboxhead{Fiche bilan}{poporange}{L'essentiel à connaître pour l'examen}}}
% Figure encadrée avec légende
\newtcolorbox{popfigure}[1][]{enhanced,breakable=false,colback=white,colframe=popline,boxrule=1.4pt,arc=8pt,
  left=10pt,right=10pt,top=10pt,bottom=8pt,before skip=10pt,after skip=12pt,halign=center,
  lower separated=false,halign lower=center,
  fontlower=\popsemi\fontsize{9}{11.5}\selectfont\color{popmuted},
  #1}
\newcommand{\legende}[1]{\par\vspace{6pt}{\popsemi\fontsize{9}{11.5}\selectfont\color{popmuted}#1}}

% ============================================================
% Signature OnBuch+
% ============================================================
\newcommand{\poplogoinline}[1]{{\poplogo\fontsize{#1}{#1}\selectfont\color{popink}OnBuch\color{poporange}+}}
\newcommand{\signatureonbuch}{%
  \begin{tcolorbox}[enhanced,colback=popink,colframe=popink,boxrule=2.2pt,arc=10pt,
      drop shadow=poporange,left=14pt,right=14pt,top=10pt,bottom=10pt,before skip=0pt,after skip=0pt]
    \tikz[baseline=-0.7ex]\node[circle,fill=popgreen,draw=popcream,line width=1.3pt,minimum size=22pt,inner sep=0]{\popcheck{white}};%
    \hspace{9pt}\begin{minipage}[c]{0.62\linewidth}
      {\popxbold\fontsize{12}{13}\selectfont\color{popcream}Signed\ \textbullet\ {\poplogo OnBuch\color{poporange}+}}\par\vspace{2pt}
      {\raggedright\popsemi\fontsize{8}{10.5}\selectfont\color{popline}\MakeUppercase{OnBuch+ teaching team}\\\MakeUppercase{Follows the official MINESEC syllabus}\par}
    \end{minipage}\hfill
    {\poplogo\fontsize{22}{22}\selectfont\color{popcream}OnBuch\color{poporange}+}
  \end{tcolorbox}%
}

% ============================================================
% Page de garde
% #1 titre · #2 matière · #3 niveau · #4 module · #5 n° leçon · #6 durée ·
% #7 description / objectifs (texte ou itemize)
% ============================================================
\newcommand{\pagegarde}[7]{%
  \thispagestyle{empty}
  \newgeometry{a4paper,margin=1.7cm,top=1.6cm,bottom=1.4cm}%
  \begin{tikzpicture}[remember picture,overlay]
    \fill[poporange!18] ([xshift=-3.2cm,yshift=-3.6cm]current page.north east) circle (4.6cm);
    \fill[poppurple!14] ([xshift=2.4cm,yshift=7.6cm]current page.south west) circle (3.8cm);
  \end{tikzpicture}%
  {\poplogo\fontsize{30}{30}\selectfont\color{popink}OnBuch\color{poporange}+}\hfill
  \raisebox{0.6ex}{\poppill{Signed course \textbullet\ Premium}{popink}{popcream}}\par
  \vspace{1pt}\popsquiggle{poppurple}\par
  \vspace{8pt}
  \begin{tcolorbox}[enhanced,colback=poporange,colframe=popink,boxrule=2.8pt,arc=13pt,
      drop shadow=popink,left=18pt,right=18pt,top=12pt,bottom=13pt,after skip=10pt]
    \poppill{#2 \textbullet\ #3}{popink}{popcream}\hspace{5pt}\poppill{#4}{white}{popink}\par\vspace{8pt}
    {\popdisplay\fontsize{14}{14}\selectfont\color{popink}CHAPTER #5}\par\vspace{4pt}
    {\raggedright\hyphenpenalty=10000\exhyphenpenalty=10000\popdisplay\fontsize{25}{27}\selectfont\color{popcream}\MakeUppercase{#1}\par}
    \vspace{8pt}
    {\popxbold\fontsize{9.5}{9.5}\selectfont\color{popink}\MakeUppercase{Suggested time: #6}}
  \end{tcolorbox}
  \begin{tcolorbox}[enhanced,colback=white,colframe=popink,boxrule=2.2pt,arc=10pt,
      drop shadow=popink,left=16pt,right=16pt,top=10pt,bottom=10pt,after skip=8pt]
    \poppill{In this chapter}{poppurple}{white}\par\vspace{5pt}
    \popsemi\fontsize{9.3}{12.6}\selectfont\color{popink}#7
  \end{tcolorbox}
  \noindent\begin{minipage}[t]{0.487\linewidth}
    \begin{tcolorbox}[enhanced,colback=popgreen,colframe=popink,boxrule=2.2pt,arc=10pt,
        drop shadow=popink,left=11pt,right=11pt,top=10pt,bottom=10pt,height=3.1cm,valign=center]
      \tcbox[colback=white,colframe=popink,boxrule=1.3pt,arc=5pt,boxsep=0pt,nobeforeafter,
        left=3pt,right=3pt,top=3pt,bottom=3pt,valign=center]{\qrcode[height=1.9cm]{https://play.google.com/store/apps/details?id=com.luvvix.onbuch&hl=en}}%
      \hspace{8pt}\begin{minipage}[c]{0.5\linewidth}
        {\popdisplay\fontsize{11}{12.5}\selectfont\color{white}\MakeUppercase{Download OnBuch}}\par\vspace{4pt}
        {\raggedright\popsemi\fontsize{8.4}{11}\selectfont\color{white}Free \textbullet\ Google Play\\AI tutor, past papers, results\par}
      \end{minipage}
    \end{tcolorbox}
  \end{minipage}\hfill
  \begin{minipage}[t]{0.487\linewidth}
    \begin{tcolorbox}[enhanced,colback=poppurple,colframe=popink,boxrule=2.2pt,arc=10pt,
        drop shadow=popink,left=11pt,right=11pt,top=10pt,bottom=10pt,height=3.1cm,valign=center]
      \tikz[baseline=-0.7ex]\node[circle,fill=white,draw=popink,line width=1.3pt,minimum size=1.6cm,inner sep=0]{\popdisplay\fontsize{24}{24}\selectfont\color{poppurple}+};%
      \hspace{8pt}\begin{minipage}[c]{0.6\linewidth}
        {\popdisplay\fontsize{11}{12.5}\selectfont\color{white}\MakeUppercase{Go OnBuch+}}\par\vspace{4pt}
        {\raggedright\popsemi\fontsize{8.4}{11}\selectfont\color{white}All signed courses, unlimited AI tutor, PDF sheets — OnBuch+ tab in the app\par}
      \end{minipage}
    \end{tcolorbox}
  \end{minipage}
  \vspace{8pt}
  \popcontenu
  \vfill
  \signatureonbuch
  \restoregeometry
  \newpage
}

% Rangée « Ce cours contient » (page de garde)
\newcommand{\poptile}[3]{% #1 couleur #2 gros texte #3 légende
  \begin{tcolorbox}[enhanced,colback=#1,colframe=popink,boxrule=2pt,arc=9pt,shadow={2.5pt}{-2.5pt}{0pt}{popink},
      left=6pt,right=6pt,top=9pt,bottom=9pt,halign=center,valign=center,height=1.75cm,width=\linewidth,nobeforeafter]
    {\popdisplay\fontsize{12.5}{13}\selectfont\color{white}\MakeUppercase{#2}}\par\vspace{3pt}
    {\centering\popsemi\fontsize{7.8}{9.5}\selectfont\color{white}#3\par}
  \end{tcolorbox}}
\newcommand{\popcontenu}{%
  \par\noindent{\popxboldls\fontsize{7.5}{7.5}\selectfont\color{popmuted}THIS SIGNED COURSE CONTAINS}\par\vspace{7pt}
  \noindent\begin{minipage}[t]{0.235\linewidth}\poptile{popblue}{Course}{complete, illustrated, step by step}\end{minipage}\hfill
  \begin{minipage}[t]{0.235\linewidth}\poptile{poporange}{Methods}{and worked examples}\end{minipage}\hfill
  \begin{minipage}[t]{0.235\linewidth}\poptile{poppink}{Exercises}{exam style + solutions}\end{minipage}\hfill
  \begin{minipage}[t]{0.235\linewidth}\poptile{popgreen}{Summary sheet}{the essentials to remember}\end{minipage}\par}

% Sommaire
\newtcolorbox{sommaire}{enhanced,colback=white,colframe=popink,boxrule=2.2pt,arc=10pt,
  drop shadow=popink,left=18pt,right=18pt,top=13pt,bottom=13pt,before skip=6pt,after skip=20pt,
  before upper={\poppill{Contents}{poppurple}{white}\par\vspace{9pt}\popsemi\fontsize{10.6}{14.5}\selectfont\color{popink}}}

% Signature de fin de cours — poussée tout en bas de la dernière page
\newcommand{\findecours}{%
  \par\vfill\nopagebreak
  \begin{center}\popsquiggle{poporange}\end{center}\vspace{4pt}
  \signatureonbuch
}
\AtBeginDocument{\def\llbracket{\mathopen{[\mkern-3.5mu[}}\def\rrbracket{\mathclose{]\mkern-3.5mu]}}} % intervalles d'entiers (glyphe ⟦ absent de la police)
% Glyphes absents de Fira Math : redéfinis à partir de glyphes disponibles
\AtBeginDocument{%
  \def\setminus{\mathbin{\backslash}}\let\smallsetminus\setminus
  \def\vdots{\mathord{\vbox{\baselineskip3.2pt\lineskiplimit0pt\kern4pt\hbox{.}\hbox{.}\hbox{.}}}}%
  \def\ddots{\mathinner{\mkern1mu\raise6.5pt\hbox{.}\mkern2mu\raise3.6pt\hbox{.}\mkern2mu\raise0.7pt\hbox{.}\mkern1mu}}%
  \def\triangle{\mathord{\scalebox{0.9}{$\symup{\Delta}$}}}%
  \def\nabla{\mathord{\raisebox{\depth}{\scalebox{1}[-1]{$\symup{\Delta}$}}}}%
  \def\checkmark{\popcheck{popdark}}%
  \def\star{\mathbin{\ast}}\def\bigstar{\mathord{\ast}}%
  \def\diamond{\mathbin{\scalebox{0.6}{$\Diamond$}}}%
  \def\bigcup{\mathop{\mathchoice{\raisebox{-0.3em}{\scalebox{1.7}{$\cup$}}}{\scalebox{1.25}{$\cup$}}{\cup}{\cup}}\displaylimits}%
  \def\bigcap{\mathop{\mathchoice{\raisebox{-0.3em}{\scalebox{1.7}{$\cap$}}}{\scalebox{1.25}{$\cap$}}{\cap}{\cap}}\displaylimits}%
  \def\lesseqqgtr{\mathrel{\lessgtr}}\def\bigoplus{\mathop{\oplus}\displaylimits}%
}
\providecommand{\ssi}{if and only if} % abréviation fréquente
\AtBeginDocument{\providecommand{\wideparen}{\overparen}} % arc de cercle

\begin{document}

\pagegarde{Descriptive inorganic Chemistry: Halogens, Group IV and d-block elements}{Chemistry}{Upper Sixth}{Upper Sixth — Chemistry}{3}{20 periods}{This chapter explores the descriptive inorganic chemistry of three important families of the periodic table: the halogens (Group VII), Group IV elements from carbon to lead, and the d-block transition metals of the first series. Learners will master trends in physical and chemical properties, key compounds and their reactions, and the characteristic features of transition metals — variable oxidation states, coloured complexes, catalysis and magnetism — all essential for GCE Advanced Level success.
\vspace{4pt}
\begin{multicols}{2}\small
\begin{itemize}
\item By the end of this chapter I can describe the physical states, colours and trends in physical properties of the halogens and relate them to their uses.
\item By the end of this chapter I can compare the chemical reactivity of the halogens and explain it using oxidizing power and electronegativity trends.
\item By the end of this chapter I can write balanced equations for the reactions of halogens and describe the preparation, properties and tests for hydrogen halides and halide ions.
\item By the end of this chapter I can name the oxoacids of the halogens and compare their acid strength, oxidizing strength and thermal stability.
\item By the end of this chapter I can describe the trends in physical properties of Group IV elements and explain the unique behaviour and allotropy of carbon.
\item By the end of this chapter I can write equations for the reactions of Group IV elements with oxygen, acids and alkalis, and describe the formulae, structures and properties of their hydrides, oxides and chlorides.
\end{itemize}\end{multicols}}

\begin{sommaire}
\begin{multicols}{2}
\begin{enumerate}
\item The Halogens: Introduction, Physical Properties and Comparative Chemistry
\item Variable Oxidation States and Reactions of the Halogens
\item Hydrogen Halides, Halide Tests and Oxoacids of the Halogens
\item Group IV Elements: Introduction, Allotropy and Trends in Physical Properties
\item Reactions and Compounds of Group IV: Hydrides, Oxides and Chlorides
\item Introduction to the d-Block: Electronic Configurations and Physical Trends
\item Variable Oxidation States and Complex Formation in Transition Metals
\item Naming, Shapes and Isomerism of Complexes
\item Colour, Catalysis, Magnetism and Uses of d-Block Elements
\item Integration activity
\item Methods and worked examples
\item Exercises
\item Solutions to the exercises
\item Summary sheet
\end{enumerate}
\end{multicols}
\end{sommaire}

\begin{objectifs}
\begin{itemize}
\item By the end of this chapter I can describe the physical states, colours and trends in physical properties of the halogens and relate them to their uses.
\item By the end of this chapter I can compare the chemical reactivity of the halogens and explain it using oxidizing power and electronegativity trends.
\item By the end of this chapter I can write balanced equations for the reactions of halogens and describe the preparation, properties and tests for hydrogen halides and halide ions.
\item By the end of this chapter I can name the oxoacids of the halogens and compare their acid strength, oxidizing strength and thermal stability.
\item By the end of this chapter I can describe the trends in physical properties of Group IV elements and explain the unique behaviour and allotropy of carbon.
\item By the end of this chapter I can write equations for the reactions of Group IV elements with oxygen, acids and alkalis, and describe the formulae, structures and properties of their hydrides, oxides and chlorides.
\item By the end of this chapter I can write the electronic configurations (spdf and electron-in-boxes) of the first transition series elements and their ions.
\item By the end of this chapter I can explain the characteristic properties of transition metals: variable oxidation states, complex formation, colour, catalytic activity and magnetic behaviour.
\item By the end of this chapter I can name coordination complexes, predict their shapes and identify structural and stereoisomerism in complexes.
\end{itemize}
\end{objectifs}

\begin{prerequis}
\begin{itemize}
\item Atomic structure and electronic configuration (s, p, d orbitals) from Lower Sixth
\item Periodic table organization and periodic trends (atomic radius, ionization energy, electronegativity)
\item Chemical bonding: ionic, covalent, dative (coordinate) bonding and molecular shapes (VSEPR)
\item Oxidation numbers and redox reactions (oxidizing and reducing agents)
\item Acids, bases and qualitative tests for ions from Ordinary Level
\end{itemize}
\end{prerequis}

\newpage

\coursec{The Halogens: Introduction, Physical Properties and Comparative Chemistry}

At the municipal water treatment station in Douala, chlorine is added to kill germs, yet the same element that disinfects our drinking water was once used as a deadly war gas. In the markets of Kribi and Bamenda, iodised table salt is sold to prevent goitre. Why do elements in the same family behave so similarly yet differ so dramatically --- fluorine a pale yellow gas, iodine a shiny grey solid? This chapter decodes the chemistry of the halogens, the Group IV elements and the d-block metals that surround us daily. We begin with the halogens, the most reactive family of non-metals in the periodic table.

\courssub{The Halogens: Position in the Periodic Table and Electronic Configuration}

Look at the right-hand side of the periodic table, one column before the noble gases. There you find five elements that chemists group together because they behave like members of one family: fluorine (F), chlorine (Cl), bromine (Br), iodine (I) and astatine (At). The name \emph{halogen} comes from Greek words meaning ``salt-former'', because these elements combine directly with metals to give salts --- sodium chloride, our common table salt, being the most famous.

What makes them a family? Write out their outer electronic configurations:

\begin{center}
\begin{tabularx}{\linewidth}{|l|l|X|}
\hline
\rowcolor{popblueL} \textbf{Element} & \textbf{Symbol} & \textbf{Outer electronic configuration} \\
\hline
Fluorine & F & $2s^2\,2p^5$ \\
\hline
Chlorine & Cl & $3s^2\,3p^5$ \\
\hline
Bromine & Br & $4s^2\,4p^5$ \\
\hline
Iodine & I & $5s^2\,5p^5$ \\
\hline
Astatine & At & $6s^2\,6p^5$ \\
\hline
\end{tabularx}
\end{center}

Every halogen is \emph{one electron short} of a stable noble gas configuration. This single fact explains almost everything about their chemistry: each halogen atom readily gains one electron (forming the ion \ce{X-}) or shares one electron (forming a single covalent bond). It also explains why the free elements exist as \cle{diatomic molecules} \ce{X2}: two atoms share a pair of electrons in a single covalent bond, and each atom then achieves a complete octet.

\begin{definition}[The halogens]
A \cle{halogen} is an element of Group VII (Group 17) of the periodic table, with outer electronic configuration $ns^2\,np^5$. The free elements exist as diatomic molecules \ce{X2} ($X = F, Cl, Br, I$) held together by a single covalent bond. The halogens are fluorine, chlorine, bromine, iodine and astatine.
\end{definition}

\begin{attention}[Astatine]
Astatine is radioactive and extremely rare; its chemistry is studied only in tracer amounts. In GCE questions, ``the halogens'' almost always means F, Cl, Br and I. Do not waste time discussing At unless the question demands it.
\end{attention}

\courssub{Colours and Physical States at Room Temperature}

If you could line up samples of the four common halogens on a bench (do not try this --- they are all toxic!), you would see a beautiful gradation: a pale yellow gas, a greenish-yellow gas, a dark red-brown liquid that fumes, and a grey-black crystalline solid which, when gently warmed, sublimes into a striking purple vapour.

\begin{center}
\begin{tabularx}{\linewidth}{|l|X|X|}
\hline
\rowcolor{popblueL} \textbf{Halogen} & \textbf{Appearance at room temperature} & \textbf{State} \\
\hline
\ce{F2} & Pale yellow & Gas \\
\hline
\ce{Cl2} & Greenish-yellow (yellow-green) & Gas \\
\hline
\ce{Br2} & Red-brown (dark red liquid, orange-brown vapour) & Liquid \\
\hline
\ce{I2} & Grey-black shiny solid; purple vapour on warming & Solid (sublimes) \\
\hline
\end{tabularx}
\end{center}

\begin{attention}[Colour of element vs colour of solution]
Do not confuse the colour of the halogen vapour with the colour of its solution. Iodine gives a \textbf{brown} solution in water (or in aqueous KI, where it dissolves by forming the tri-iodide ion \ce{I3-}) but a \textbf{purple/violet} solution in organic solvents such as cyclohexane or tetrachloromethane. Similarly, bromine water is orange-yellow, while bromine in cyclohexane is orange-red. Examiners love this trap.
\end{attention}

\begin{popfigure}
\begin{tikzpicture}[scale=1.0]
\draw[->, very thick, popdark] (6.6,0.4) -- (6.6,-8.6) node[midway, right, align=left, text width=3.4cm]{\small Increasing melting point, boiling point and molecular size};
\draw[fill=popgold!40, draw=popdark] (0,0) rectangle (4.5,-1.6);
\node[align=center] at (2.25,-0.8) {\textbf{Fluorine \ce{F2}}\\ \small pale yellow gas};
\draw[fill=popgreen!30, draw=popdark] (0,-2.1) rectangle (4.5,-3.7);
\node[align=center] at (2.25,-2.9) {\textbf{Chlorine \ce{Cl2}}\\ \small greenish-yellow gas};
\draw[fill=poporange!45, draw=popdark] (0,-4.2) rectangle (4.5,-5.8);
\node[align=center] at (2.25,-5.0) {\textbf{Bromine \ce{Br2}}\\ \small red-brown liquid};
\draw[fill=poppurple!35, draw=popdark] (0,-6.3) rectangle (4.5,-7.9);
\node[align=center] at (2.25,-7.1) {\textbf{Iodine \ce{I2}}\\ \small grey-black solid, purple vapour};
\end{tikzpicture}
\legende{Group VII at room temperature: state and colour change progressively down the group, while melting point, boiling point and molecular size increase.}
\end{popfigure}

\courssub{Trends in Physical Properties Down Group VII}

\textbf{Melting and boiling points.} The change of state from gas (\ce{F2}, \ce{Cl2}) to liquid (\ce{Br2}) to solid (\ce{I2}) tells us immediately that melting and boiling points \emph{increase} down the group. Why? The molecules \ce{X2} are non-polar, so the only attractions between them are \cle{van der Waals forces} (instantaneous dipole--induced dipole forces). Down the group, the number of electrons per molecule increases, the electron cloud becomes larger and more easily polarised, so the temporary dipoles become stronger and the van der Waals attractions increase. More energy is therefore needed to separate the molecules, and the melting and boiling points rise.

\begin{propriete}[Trend in melting and boiling points]
Down Group VII, melting and boiling points \textbf{increase} because the strength of the van der Waals forces between the \ce{X2} molecules increases with the size (and number of electrons) of the molecule. \emph{(Explanation examinable --- learn the three steps: more electrons $\rightarrow$ more polarisable cloud $\rightarrow$ stronger van der Waals forces.)}
\end{propriete}

\textbf{Atomic radius.} Down the group each element has one more filled electron shell, so the atomic radius increases steadily: F $<$ Cl $<$ Br $<$ I.

\textbf{Electronegativity and first ionization energy.} Electronegativity (the power of an atom in a molecule to attract the bonding electrons) \emph{decreases} down the group: fluorine is the most electronegative element of the whole periodic table. The first ionization energy also \emph{decreases} down the group, because the outermost electron is further from the nucleus and better shielded by inner shells, so it is held less tightly and is easier to remove.

\begin{popfigure}
\begin{tikzpicture}
\begin{axis}[
    ybar, bar width=10pt,
    width=12cm, height=7cm,
    ymin=0, ymax=4.5,
    symbolic x coords={F, Cl, Br, I},
    xtick=data,
    ylabel={Value},
    legend style={at={(0.98,0.98)}, anchor=north east, font=\small},
    ymajorgrids=true, grid style={popline!40},
    axis line style={popdark},
    tick label style={font=\small},
    ylabel style={font=\small},
]
\addplot[fill=popblue, draw=popdark] coordinates {(F,4.0) (Cl,3.0) (Br,2.8) (I,2.5)};
\addlegendentry{Electronegativity (Pauling)}
\addplot[fill=poporange, draw=popdark] coordinates {(F,4.2) (Cl,3.1) (Br,2.85) (I,2.5)};
\addlegendentry{First I.E. (scaled: $E_i/400$ in $\mathrm{kJ\,mol^{-1}}$)}
\end{axis}
\end{tikzpicture}
\legende{Electronegativity and first ionization energy both decrease down Group VII (first ionization energies, about 1681, 1251, 1140 and 1008 $\mathrm{kJ\,mol^{-1}}$ for F, Cl, Br, I, are scaled by $1/400$ so both quantities fit on one axis).}
\end{popfigure}

\begin{aretenir}[Summary of trends down Group VII]
\begin{itemize}
\item Atomic radius: \textbf{increases} (extra electron shells).
\item Melting and boiling points: \textbf{increase} (stronger van der Waals forces).
\item Electronegativity: \textbf{decreases}.
\item First ionization energy: \textbf{decreases}.
\item Oxidizing power (reactivity): \textbf{decreases}.
\end{itemize}
\end{aretenir}

\courssub{Comparative Chemistry of the Halogens}

\textbf{Trend in oxidizing power.} A halogen reacts by \emph{gaining} an electron: \ce{X2 + 2e- -> 2X-}. It is therefore an \cle{oxidizing agent}. The more easily the atom gains the electron, the stronger the oxidizing agent. Down the group the atomic radius increases, so the nucleus attracts the incoming electron less strongly (more shielding, greater distance), and the oxidizing power \emph{decreases}: \ce{F2} $>$ \ce{Cl2} $>$ \ce{Br2} $>$ \ce{I2}.

\begin{propriete}[Oxidizing power down the group]
Oxidizing power decreases down Group VII because the atom's ability to attract and gain an electron weakens as the atomic radius increases and shielding by inner electron shells grows. Fluorine is the strongest oxidizing agent of all the elements.
\end{propriete}

\textbf{Evidence 1: displacement reactions.} A more reactive (more oxidizing) halogen displaces a less reactive one from a solution of its halide. Bubble chlorine through potassium bromide solution and the colourless solution turns orange-brown --- bromine has been set free:
\[
\ce{Cl2(aq) + 2KBr(aq) -> 2KCl(aq) + Br2(aq)}
\qquad\text{ionic:}\qquad
\ce{Cl2 + 2Br- -> 2Cl- + Br2}
\]
Similarly, chlorine displaces iodine from potassium iodide (solution turns brown, and gives a blue-black colour with starch):
\[
\ce{Cl2(aq) + 2KI(aq) -> 2KCl(aq) + I2(aq)}
\]
and bromine displaces iodine:
\[
\ce{Br2(aq) + 2KI(aq) -> 2KBr(aq) + I2(aq)}
\]
But iodine cannot displace bromine or chlorine, and bromine cannot displace chlorine. The order of displacement exactly matches the order of oxidizing power.

\begin{methode}[Comparing halogen reactivity by displacement]
\begin{enumerate}
\item Take aqueous solutions of the halides (e.g.\ \ce{KCl}, \ce{KBr}, \ce{KI}) in separate test tubes.
\item Add a few drops of the halogen water (e.g.\ chlorine water or bromine water) to each.
\item Add a little cyclohexane, shake, and observe the colour of the organic layer: purple indicates \ce{I2}, orange-red indicates \ce{Br2}.
\item Conclusion: a colour change means displacement has occurred, so the added halogen is \emph{more} reactive than the halogen of the halide. No change means it is \emph{less} reactive.
\end{enumerate}
\end{methode}

\textbf{Evidence 2: reactions with hydrogen.} All halogens combine with hydrogen to give hydrogen halides, \ce{H2 + X2 -> 2HX}, but the vigour decreases dramatically down the group:

\begin{center}
\begin{tabularx}{\linewidth}{|l|X|}
\hline
\rowcolor{popblueL} \textbf{Halogen} & \textbf{Reaction with \ce{H2}} \\
\hline
\ce{F2} & Explodes violently, even in the dark and in the cold. \\
\hline
\ce{Cl2} & Explodes in sunlight (or on ignition); slow in the dark. \\
\hline
\ce{Br2} & Combines on heating; moderate speed. \\
\hline
\ce{I2} & Reacts only on strong heating, and the reaction is \emph{reversible} (partial): \ce{H2 + I2 <=> 2HI}. \\
\hline
\end{tabularx}
\end{center}

The decreasing vigour, and the decreasing thermal stability of the hydrogen halides formed (\ce{HF} is very stable; \ce{HI} decomposes easily on heating), confirm that reactivity falls down the group.

\textbf{Evidence 3: reactions with metals and with phosphorus.} Heated metals react with halogens to give halides, the vigour again decreasing down the group. Sodium burns brightly in chlorine to give white fumes of sodium chloride, \ce{2Na(s) + Cl2(g) -> 2NaCl(s)}. With iron, the comparison is especially instructive: chlorine oxidises iron to iron(III), \ce{2Fe(s) + 3Cl2(g) -> 2FeCl3(s)}, whereas the weaker oxidant iodine gives only iron(II), \ce{Fe(s) + I2(s) -> FeI2(s)}. With white phosphorus, excess chlorine gives \ce{PCl5} (\ce{P4 + 10Cl2 -> 4PCl5}) while iodine can only manage \ce{PI3} --- again showing chlorine's greater oxidizing strength.

\textbf{Reaction with water and with alkalis: disproportionation.} Chlorine dissolves in water to give ``chlorine water'', a mixture of two acids:
\[
\ce{Cl2(aq) + H2O(l) <=> HCl(aq) + HOCl(aq)}
\]
Here chlorine is simultaneously reduced (to \ce{Cl-} in \ce{HCl}, oxidation state $-1$) and oxidised (to \ce{HOCl}, oxidation state $+1$). A reaction in which the same element is both oxidised and reduced is called a \cle{disproportionation}. The hypochlorous acid \ce{HOCl} is the active germ-killer in water treatment.

With cold dilute sodium hydroxide, chlorine gives chloride and hypochlorite --- this is the basis of household bleach:
\[
\ce{Cl2(g) + 2NaOH(aq) -> NaCl(aq) + NaOCl(aq) + H2O(l)} \quad (\text{cold, dilute})
\]
With hot concentrated sodium hydroxide, the products are chloride and chlorate(V):
\[
\ce{3Cl2(g) + 6NaOH(aq) -> 5NaCl(aq) + NaClO3(aq) + 3H2O(l)} \quad (\text{hot, concentrated})
\]

\begin{exoresolu}[Displacement and oxidizing power]
\textbf{Statement.} A student adds chlorine water to a solution of potassium iodide, then shakes the mixture with cyclohexane. (a) Describe what is observed. (b) Write the ionic equation for the reaction. (c) Explain, in terms of electron transfer, what this shows about the relative oxidizing powers of chlorine and iodine. (d) Would a reaction occur if iodine water were added to potassium chloride solution? Justify.
\tcblower
\textbf{Solution.}
(a) The aqueous layer turns brown (formation of \ce{I2}); after shaking with cyclohexane, the upper organic layer becomes purple/violet.

(b) \ce{Cl2(aq) + 2I-(aq) -> 2Cl-(aq) + I2(aq)}.

(c) Each iodide ion loses one electron (oxidation): \ce{2I- -> I2 + 2e-}. Each chlorine molecule gains two electrons (reduction): \ce{Cl2 + 2e- -> 2Cl-}. Since chlorine \emph{takes} electrons from iodide, chlorine is the stronger oxidizing agent; iodine is the weaker one. This agrees with the decrease of oxidizing power down Group VII.

(d) No reaction. Iodine is a weaker oxidizing agent than chlorine, so it cannot remove electrons from \ce{Cl-} ions. The mixture would simply show the brown colour of iodine, unchanged.
\end{exoresolu}

\courssub{Uses of the Halogens}

\begin{itemize}
\item \textbf{Chlorine:} sterilisation of drinking water and swimming pools (as at the Douala treatment station); manufacture of bleach (\ce{NaOCl}), of hydrochloric acid, of plastics such as PVC, and of solvents.
\item \textbf{Fluorine compounds:} fluorides (e.g.\ \ce{NaF}) in toothpaste to prevent tooth decay; CFCs formerly used as refrigerants and propellants (now phased out because they destroy the ozone layer); the non-stick polymer PTFE (Teflon).
\item \textbf{Bromine:} silver bromide \ce{AgBr} is the light-sensitive salt in photographic film; bromine compounds are used as flame retardants.
\item \textbf{Iodine:} tincture of iodine (iodine in ethanol) as an antiseptic for wounds; iodine added to table salt (as \ce{KIO3} or \ce{KI}) to prevent goitre --- a public-health measure visible in every Cameroonian market.
\end{itemize}

\begin{savaistu}[Chlorine: saviour and killer]
Chlorination of drinking water is one of the greatest public-health advances in history, saving millions from cholera and typhoid. Yet chlorine was also the first chemical weapon used on a large scale (Ypres, 1915). The same oxidizing power that destroys bacteria in your water attacks lung tissue when inhaled --- a sobering reminder that in chemistry, dose and context decide whether a substance is medicine or poison.
\end{savaistu}

\begin{attention}[Exam tip]
GCE examiners very often ask: ``Explain the trend in melting points of the halogens.'' Always answer in terms of the \textbf{increasing strength of van der Waals forces as the molecules become larger} (more electrons, more polarisable). Never say ``the covalent bond gets stronger'' --- melting \ce{I2} does not break any covalent bond; it only separates the molecules from one another.
\end{attention}

\begin{exercice}[Trends and displacement]
\begin{enumerate}
\item Arrange \ce{F2}, \ce{Cl2}, \ce{Br2}, \ce{I2} in order of increasing boiling point and justify your answer.
\item Predict and explain what happens when bromine water is added to (a) potassium chloride solution, (b) potassium iodide solution. Write ionic equations where reaction occurs.
\item Explain why the first ionization energy decreases from fluorine to iodine.
\end{enumerate}
\end{exercice}

\begin{corrige}[Exercise 1]
\begin{enumerate}
\item \ce{F2} $<$ \ce{Cl2} $<$ \ce{Br2} $<$ \ce{I2}. Down the group the \ce{X2} molecules contain more electrons and have larger, more polarisable electron clouds, so the van der Waals forces between molecules strengthen and more energy is needed to separate them.
\item (a) No reaction with \ce{KCl}: bromine is a weaker oxidizing agent than chlorine and cannot displace it. (b) With \ce{KI} the solution turns brown: \ce{Br2 + 2I- -> 2Br- + I2}; bromine, being the stronger oxidant, takes electrons from the iodide ions.
\item Down the group the outer electron occupies a shell further from the nucleus and is more strongly shielded by inner shells; the effective attraction on it is weaker, so less energy is needed to remove it.
\end{enumerate}
\end{corrige}

\coursec{Variable Oxidation States and Reactions of the Halogens}

In the previous section we saw that the halogens are powerful oxidising agents whose strength \emph{decreases} down the group: fluorine is the most ferocious electron-grabber known, while iodine is comparatively gentle. We now ask two deeper questions. First, if halogens love electrons so much, how can compounds such as \ce{HClO4} exist, in which chlorine appears to have \emph{lost} electrons? Second, how do the halogens actually behave when they meet hydrogen, metals, water or alkali? These two questions — the \cle{variable oxidation states} of the halogens and their \cle{reactions} — form the heart of this section and are examined almost every year at the GCE Advanced Level.

\courssub{Oxidation States Shown by the Halogens}

\textbf{The puzzle.} In sodium chloride, \ce{NaCl}, chlorine has clearly gained one electron: its oxidation state is $-1$. Yet in perchloric acid, \ce{HClO4}, chlorine is bonded to four oxygens, each more electronegative than itself, and its oxidation state is $+7$. The same element swings from $-1$ to $+7$ — a span of eight units. Understanding \emph{why} is the key to all halogen oxo-chemistry.

\begin{methode}[Assigning oxidation numbers in oxo-compounds]
To find the oxidation state of a halogen in an oxo-compound:
\begin{enumerate}
\item Assign oxygen its usual value $-2$ and hydrogen $+1$.
\item Write that the sum of all oxidation states equals the overall charge of the species ($0$ for a neutral molecule, the ionic charge for an ion).
\item Solve for the halogen.
\end{enumerate}
Example with \ce{HClO3}: $+1 + x + 3(-2) = 0$, so $x = +5$. Chlorine is in oxidation state $+5$.
Example with \ce{ClO2-}: $x + 2(-2) = -1$, so $x = +3$.
\end{methode}

Applying this method to the chlorine oxoacids and their anhydrides gives the full ladder of oxidation states:

\begin{center}
\begin{tabularx}{\linewidth}{|l|c|X|}
\hline
\rowcolor{popblueL} \textbf{Compound} & \textbf{Oxidation state of Cl} & \textbf{Name (Stock notation)} \\
\hline
\ce{NaCl} & $-1$ & Sodium chloride \\
\hline
\ce{HClO} (or \ce{Cl2O}) & $+1$ & Chloric(I) acid (hypochlorous acid) \\
\hline
\ce{HClO2} & $+3$ & Chloric(III) acid (chlorous acid) \\
\hline
\ce{HClO3} & $+5$ & Chloric(V) acid (chloric acid) \\
\hline
\ce{HClO4} & $+7$ & Chloric(VII) acid (perchloric acid) \\
\hline
\end{tabularx}
\end{center}

\begin{popfigure}
\begin{center}
\begin{tikzpicture}[scale=1.0]
\draw[->, thick, popmuted] (-0.5,0) -- (10.5,0) node[right, popdark]{\small oxidation state of Cl};
\foreach \x/\os in {0.5/{-1}, 2.8/{+1}, 5.1/{+3}, 7.4/{+5}, 9.7/{+7}}{
  \draw[thick, popdark] (\x,0.08) -- (\x,-0.08) node[below, popdark]{\small $\os$};
}
\node[draw=popblue, fill=popblueL, rounded corners, align=center, text width=1.7cm] at (0.5,1.6) {\small \ce{NaCl}\\ \scriptsize sodium chloride};
\node[draw=popgreen, fill=popgreenL, rounded corners, align=center, text width=1.7cm] at (2.8,1.6) {\small \ce{HClO}\\ \scriptsize chloric(I) acid};
\node[draw=poporange, fill=white, rounded corners, align=center, text width=1.7cm] at (5.1,1.6) {\small \ce{HClO2}\\ \scriptsize chloric(III) acid};
\node[draw=poppurple, fill=white, rounded corners, align=center, text width=1.7cm] at (7.4,1.6) {\small \ce{HClO3}\\ \scriptsize chloric(V) acid};
\node[draw=popink, fill=white, rounded corners, align=center, text width=1.7cm] at (9.7,1.6) {\small \ce{HClO4}\\ \scriptsize chloric(VII) acid};
\foreach \x in {0.5,2.8,5.1,7.4,9.7}{ \draw[thin, popmuted, dashed] (\x,0.1) -- (\x,0.9); }
\end{tikzpicture}
\end{center}
\legende{The ladder of chlorine oxidation states, from $-1$ in \ce{NaCl} to $+7$ in \ce{HClO4}.}
\end{popfigure}

\courssub{Why Fluorine Shows Only $-1$}

\begin{propriete}[Fluorine: only $-1$]
Fluorine \textbf{never} shows positive oxidation states. Two reasons combine:
\begin{enumerate}
\item It is the \cle{most electronegative element} in the Periodic Table ($\chi = 4.0$), so in \emph{every} compound the bonding electrons are assigned to fluorine.
\item Being in Period 2, it has \textbf{no accessible d-orbitals}; its octet cannot be expanded, so it can never form more than one covalent bond and can never be oxidised beyond $-1$ (its only states are $0$ in \ce{F2} and $-1$ in compounds).
\end{enumerate}
\end{propriete}

\textbf{Why Cl, Br and I can be positive.} Chlorine ($[\ce{Ne}]3s^2 3p^5$), bromine and iodine possess \cle{empty d-orbitals} of the valence shell ($3d$ for Cl, $4d$ for Br, $5d$ for I). When bonded to a more electronegative element such as oxygen or fluorine, paired $ns$ and $np$ electrons can be \emph{promoted} into these empty $nd$ orbitals, unpairing up to seven electrons and allowing 1, 3, 5 or 7 covalent bonds. Hence the oxidation states $+1, +3, +5, +7$ appear. The energy cost of promotion is repaid by the extra bonds formed — but only with very electronegative partners, which is why the positive states occur in \emph{oxides, oxoacids and interhalogens}, never with metals.

\begin{attention}[Common mistake]
Do not write that fluorine shows $+7$ ``like chlorine''. Fluorine has no $d$-orbitals (Period 2) and is more electronegative than oxygen, so \ce{FO} bonds still assign electrons to F. Its maximum oxidation state in any compound is $-1$.
\end{attention}

\courssub{Reaction of the Halogens with Hydrogen}

All halogens combine with hydrogen to give the hydrogen halides \ce{HX}, but the vigour of the reaction falls dramatically down the group — a beautiful illustration of decreasing oxidising power:

\begin{align*}
\ce{H2(g) + F2(g) & -> 2HF(g)} && \text{explosive, even in the dark and cold}\\
\ce{H2(g) + Cl2(g) & -> 2HCl(g)} && \text{explosive in sunlight (photochemical), slow in the dark}\\
\ce{H2(g) + Br2(g) & <=> 2HBr(g)} && \text{needs heat ($\sim$200\ensuremath{{}^{\circ}\mathrm{C}}) and a Pt catalyst; equilibrium}\\
\ce{H2(g) + I2(g) & <=> 2HI(g)} && \text{slow, reversible, even on heating; equilibrium lies far left}
\end{align*}

Because the \ce{H-X} bond weakens as the halogen atom grows larger (bond enthalpy: \ce{H-F} 562, \ce{H-Cl} 431, \ce{H-Br} 366, \ce{H-I} 299 kJ mol$^{-1}$), the \cle{thermal stability} of the hydrogen halides decreases down the group: \ce{HF} and \ce{HCl} resist heat, \ce{HBr} decomposes partially on strong heating, and \ce{HI} decomposes appreciably into \ce{H2} and \ce{I2} (notice the violet fumes) when a hot glass rod is held in the gas.

\courssub{Reaction with Metals and Non-metals}

\textbf{Metals.} Halogens oxidise metals to \cle{ionic halides}; the more vigorous the halogen, the higher the oxidation state of the metal product:

\begin{align*}
\ce{2Na(s) + Cl2(g) & -> 2NaCl(s)}\\
\ce{2Fe(s) + 3Cl2(g) & -> 2FeCl3(s)} && \text{excess chlorine}\\
\ce{2Fe(s) + I2(g) & -> 2FeI2(s)} && \text{iodine is weaker, stops at Fe(II)}\\
\ce{2Al(s) + 3Cl2(g) & -> 2AlCl3(s)} && \text{sublimes as dimer \ce{Al2Cl6}}
\end{align*}

Note that chlorine, a strong oxidant, drives iron to its \emph{higher} state, iron(III). With the weaker iodine, iron gives only \ce{FeI2} — direct evidence of the oxidising-power trend.

\textbf{Non-metals.} With phosphorus, chlorine gives \ce{PCl3} (limited chlorine) or \ce{PCl5} (excess chlorine); sulphur gives \ce{SCl2}:

\begin{align*}
\ce{2P(s) + 3Cl2(g) & -> 2PCl3(l)} &
\ce{2P(s) + 5Cl2(g) & -> 2PCl5(s)} &
\ce{S(s) + Cl2(g) & -> SCl2(l)}
\end{align*}

\courssub{Chlorine and Water: the Bleaching Reaction}

Chlorine dissolves slightly in water ($\sim$0.09 mol dm$^{-3}$ at 298 K, 1 atm) and reacts reversibly:

\[ \ce{Cl2(g) + H2O(l) <=> HCl(aq) + HClO(aq)} \quad K_c \approx 4 \times 10^{-4} \]

Look at the oxidation states: chlorine goes from $0$ in \ce{Cl2} to $-1$ in \ce{HCl} \emph{and} to $+1$ in \ce{HClO}. One and the same element is reduced and oxidised in the same reaction.

\begin{definition}[Disproportionation]
A \cle{disproportionation} is a redox reaction in which atoms of the \textbf{same element} are simultaneously \textbf{oxidised} and \textbf{reduced}.
\end{definition}

The chloric(I) acid \ce{HClO} is the active \cle{bleaching agent}: it oxidises coloured organic dyes to colourless products, itself being reduced:

\[ \ce{HClO(aq) + dye -> HCl(aq) + (dye + O) (colourless)} \]

This is why moist litmus paper is bleached by chlorine (first red, then white), and why chlorine water disinfects the drinking water of many Cameroonian towns.

\courssub{Reaction with Alkali: Cold Dilute versus Hot Concentrated}

Chlorine disproportionates in sodium hydroxide, but the product depends on conditions:

\begin{align*}
\ce{Cl2(g) + 2NaOH(aq) & -> NaCl(aq) + NaClO(aq) + H2O(l)} && \text{cold, dilute}\\
\ce{3Cl2(g) + 6NaOH(aq) & -> 5NaCl(aq) + NaClO3(aq) + 3H2O(l)} && \text{hot, concentrated}
\end{align*}

In cold alkali, chlorine goes from $0$ to $-1$ (\ce{NaCl}) and $+1$ (\ce{NaClO}); in hot alkali the chlorate(I) itself disproportionates further ($3\ce{ClO-} \to 2\ce{Cl-} + \ce{ClO3-}$), giving chlorate(V), \ce{NaClO3} ($+5$). 

Bromine and iodine behave similarly, but the intermediate hypohalites (\ce{BrO-}, \ce{IO-}) are unstable even at room temperature and disproportionate rapidly to halide and halate(V):

\[ \ce{3Br2 + 6OH- -> 5Br- + BrO3- + 3H2O} \quad \text{(even cold)} \]

\begin{attention}[Hot alkali trap]
A very common mistake is to write \ce{Cl2 + 2NaOH -> NaCl + NaClO + H2O} for \emph{hot concentrated} alkali. Hot conditions give the chlorate(V): \ce{3Cl2 + 6NaOH -> 5NaCl + NaClO3 + 3H2O}. Read the question's conditions before writing the equation!
\end{attention}

\begin{aretenir}[Exam tip]
Always balance disproportionation equations carefully and show the oxidation-state changes. GCE examiners award separate marks for correct stoichiometry and for identifying the species oxidised and reduced.
\end{aretenir}

\courssub{Displacement of Halides and the Thiosulfate Reaction}

A more reactive halogen displaces a less reactive one from its halide solution — the classic test for halide ions and the basis of the redox series:

\begin{align*}
\ce{Cl2(aq) + 2KBr(aq) & -> 2KCl(aq) + Br2(aq)} && \text{orange/brown colour appears in organic layer}\\
\ce{Cl2(aq) + 2KI(aq) & -> 2KCl(aq) + I2(aq)} && \text{brown solution / black solid}\\
\ce{Br2(aq) + 2KI(aq) & -> 2KBr(aq) + I2(aq)} &&
\end{align*}

The liberated iodine can be estimated quantitatively by titration with sodium thiosulfate — a classic practical reaction:

\[ \ce{I2(aq) + 2Na2S2O3(aq) -> 2NaI(aq) + Na2S4O6(aq)} \]

Here iodine is reduced to iodide ($0 \to -1$) while thiosulfate is oxidised to tetrathionate ($+2 \to +2.5$ average). Starch indicator (deep blue-black with \ce{I2}) is added near the end-point, which is the sharp disappearance of the blue colour.

\begin{exoresolu}[Chlorine and hot alkali]
\textbf{Statement.} Chlorine is bubbled into hot concentrated sodium hydroxide. (a) Write the balanced equation. (b) State the oxidation number of chlorine in each chlorine-containing species. (c) Identify the type of redox reaction. \tcblower
\textbf{Solution.}
(a) \ce{3Cl2(g) + 6NaOH(aq) -> 5NaCl(aq) + NaClO3(aq) + 3H2O(l)}.

(b) In \ce{Cl2}: $0$. In \ce{NaCl}: $-1$. In \ce{NaClO3}: $+1 + x + 3(-2) = 0$, so $x = +5$.

(c) Since chlorine is simultaneously reduced ($0 \to -1$) and oxidised ($0 \to +5$), this is a \textbf{disproportionation}. Check the electron balance: five Cl atoms gain one electron each ($5e^-$), one Cl atom loses five electrons ($5e^-$) — consistent.
\end{exoresolu}

\begin{popfigure}
\begin{center}
\begin{tikzpicture}[node distance=1.1cm]
\node[draw=popink, fill=popblueL, rounded corners, align=center, text width=2.4cm] (cl2) {\textbf{\ce{Cl2}}\\ \scriptsize chlorine};
\node[draw=popblue, rounded corners, align=center, text width=2.6cm, above left=1.0cm and 1.6cm of cl2] (h2) {\ce{H2}, sunlight\\ \scriptsize $\to$ \ce{2HCl(g)}};
\node[draw=popblue, rounded corners, align=center, text width=2.6cm, below left=1.0cm and 1.6cm of cl2] (met) {Fe or Al, heat\\ \scriptsize $\to$ \ce{FeCl3}, \ce{AlCl3}};
\node[draw=popgreen, rounded corners, align=center, text width=2.6cm, above right=1.0cm and 1.6cm of cl2] (wat) {\ce{H2O}\\ \scriptsize $\to$ \ce{HCl} + \ce{HClO}};
\node[draw=poporange, rounded corners, align=center, text width=2.8cm, right=1.6cm of cl2] (cold) {cold dilute \ce{NaOH}\\ \scriptsize $\to$ \ce{NaCl} + \ce{NaClO}};
\node[draw=poppurple, rounded corners, align=center, text width=2.8cm, below right=1.0cm and 1.6cm of cl2] (hot) {hot conc.\ \ce{NaOH}\\ \scriptsize $\to$ \ce{5NaCl} + \ce{NaClO3}};
\draw[->, thick, popblue] (cl2) -- (h2);
\draw[->, thick, popblue] (cl2) -- (met);
\draw[->, thick, popgreen] (cl2) -- (wat);
\draw[->, thick, poporange] (cl2) -- (cold);
\draw[->, thick, poppurple] (cl2) -- (hot);
\end{tikzpicture}
\end{center}
\legende{Summary flow diagram of the principal reactions of chlorine.}
\end{popfigure}

\begin{exercice}[Oxidation states and displacement]
\begin{enumerate}
\item Determine the oxidation state of the halogen in: \ce{KBr}, \ce{Br2O}, \ce{KIO3}, \ce{HIO4}.
\item Explain, in terms of electronic structure, why iodine can show $+7$ but fluorine cannot.
\item Predict and explain what is observed when chlorine water is added to potassium iodide solution, and write the ionic equation.
\item Balance: \ce{Br2 + NaOH -> NaBr + NaBrO3 + H2O} (hot conditions).
\end{enumerate}
\end{exercice}

\begin{corrige}[Exercise 1]
\begin{enumerate}
\item \ce{KBr}: $-1$; \ce{Br2O}: $+1$; \ce{KIO3}: $+5$; \ce{HIO4}: $+7$.
\item Iodine ($[\ce{Kr}]5s^2 5p^5$) has empty $5d$ orbitals into which paired electrons can be promoted, unpairing up to seven electrons for bonding with oxygen. Fluorine (Period 2) has no $d$-orbitals and is the most electronegative element, so it can never be assigned a positive oxidation state.
\item The solution turns brown (black solid on standing): chlorine oxidises iodide to iodine. \ce{Cl2(aq) + 2I^{-}(aq) -> 2Cl^{-}(aq) + I2(aq)}.
\item \ce{3Br2 + 6NaOH -> 5NaBr + NaBrO3 + 3H2O}.
\end{enumerate}
\end{corrige}

\coursec{Hydrogen Halides, Halide Tests and Oxoacids of the Halogens}

In the previous section we saw that the halogens are powerful oxidising agents whose strength decreases down the group, and that they display variable oxidation states in their compounds. We now turn to two important families of halogen compounds: the \cle{hydrogen halides} \ce{HX}, in which the halogen has oxidation state $-1$, and the \cle{oxoacids}, in which the halogen shows positive oxidation states up to $+7$. These compounds are not only examination favourites; they are also the basis of the qualitative analysis tests for halide ions that you will meet in the practical paper.

\courssub{Preparation of the Hydrogen Halides}

\textbf{Action of concentrated sulphuric acid on metal halides}\par

The standard laboratory preparation of a hydrogen halide is the action of concentrated sulphuric acid on a solid metal halide. Concentrated \ce{H2SO4} is chosen because it is a \emph{non-volatile} acid (boiling point about $338\ ^\circ\mathrm{C}$), so it displaces the more volatile hydrogen halide from its salt:

\[
\ce{NaCl(s) + H2SO4(l) -> NaHSO4(s) + HCl(g)}
\]
On heating with excess \ce{NaCl}:
\[
\ce{NaCl(s) + NaHSO4(s) -> Na2SO4(s) + HCl(g)}
\]

\[
\ce{CaF2(s) + H2SO4(l) -> CaSO4(s) + 2HF(g)}
\]

The misty fumes of \ce{HCl} or \ce{HF} are collected by upward delivery (they are denser than air) or dissolved in water to give the aqueous acid. Note that the preparation of \ce{HF} is carried out in a lead or platinum apparatus, because \ce{HF} attacks glass (it reacts with \ce{SiO2} to form \ce{SiF4} and \ce{H2SiF6}).

This method, however, \textbf{fails for \ce{HBr} and \ce{HI}}. The reason is chemical, not practical: concentrated sulphuric acid is itself an oxidising agent, and bromide and iodide ions are good enough reducing agents to be oxidised by it. With sodium bromide, the \ce{HBr} first formed is partly oxidised:

\[
\ce{NaBr(s) + H2SO4(l) -> NaHSO4(s) + HBr(g)}
\]

\[
\ce{2HBr(g) + H2SO4(l) -> Br2(g) + SO2(g) + 2H2O(l)}
\]

The product is therefore contaminated with brown \ce{Br2} vapour and colourless, choking \ce{SO2}. With sodium iodide the oxidation goes even further, because \ce{I-} is the strongest reducing agent of the halides; a mixture of reduction products (\ce{SO2}, \ce{S} and even \ce{H2S}) is obtained:

\[
\ce{8HI(g) + H2SO4(l) -> 4I2(s) + H2S(g) + 4H2O(l)}
\]

\begin{attention}[A classic trap]
Concentrated \ce{H2SO4} cannot be used to prepare pure \ce{HBr} or \ce{HI}, because it oxidises \ce{Br-} and \ce{I-} to \ce{Br2} and \ce{I2}. In the examination, if you are asked to ``prepare pure hydrogen bromide'', do \emph{not} write the sulphuric acid method.
\end{attention}

\textbf{Alternative preparations of \ce{HBr} and \ce{HI}}\par

To avoid oxidation we replace concentrated sulphuric acid by \cle{phosphoric(V) acid}, \ce{H3PO4}, which is non-volatile but \emph{not} an oxidising agent:

\[
\ce{NaBr(s) + H3PO4(l) -> NaH2PO4(s) + HBr(g)}
\]

\[
\ce{NaI(s) + H3PO4(l) -> NaH2PO4(s) + HI(g)}
\]

Alternatively, \ce{HBr} and \ce{HI} can be prepared by the \emph{hydrolysis of phosphorus trihalides}, made \emph{in situ} from red phosphorus and the halogen:

\[
\ce{2P(s) + 3Br2(l) -> 2PBr3(l)}
\]
\[
\ce{PBr3(l) + 3H2O(l) -> H3PO3(aq) + 3HBr(g)}
\]

\[
\ce{2P(s) + 3I2(s) -> 2PI3(s)}
\]
\[
\ce{PI3(s) + 3H2O(l) -> H3PO3(aq) + 3HI(g)}
\]

\courssub{Physical Properties and Acidic Strength of the Hydrogen Halides}

All the hydrogen halides are \textbf{colourless gases} at room temperature which fume in moist air (the gas dissolves in atmospheric moisture to form tiny droplets of acid). They are \textbf{extremely soluble in water}; this is beautifully demonstrated by the \cle{fountain experiment}.

\begin{experience}[The fountain experiment]
A dry round-bottomed flask is filled with \ce{HCl} gas and fitted with a jet tube dipping into water containing blue litmus. A few drops of water are introduced into the flask; the \ce{HCl} dissolves so readily that the pressure inside drops sharply, and water is driven up through the jet as a spectacular fountain. The litmus turns red, showing that an acidic solution (hydrochloric acid) has formed. The same experiment works with \ce{HBr} and \ce{HI}.
\end{experience}

In water the hydrogen halides behave as acids: $\ce{HX(aq) + H2O(l) -> H3O+(aq) + X-(aq)}$. Their strengths differ strikingly:

\begin{propriete}[Acid strength of the hydrogen halides]
Aqueous acid strength increases down the group:
\[
\ce{HF} \ll \ce{HCl} < \ce{HBr} < \ce{HI}
\]
\ce{HF} is a \emph{weak} acid ($K_{\mathrm{a}} \approx 6.8 \times 10^{-4}$); \ce{HCl}, \ce{HBr} and \ce{HI} are strong (fully ionised). The explanation is the \textbf{strength of the \ce{H-X} bond}: as the halogen atom becomes larger down the group, the bonding pair is further from the hydrogen nucleus, the bond becomes longer and weaker, and \ce{H+} is released more easily. The \ce{H-F} bond is short and very strong (bond dissociation energy $\approx 565\ \mathrm{kJ\,mol^{-1}}$), so \ce{HF} ionises only slightly. (This proof is qualitative and is frequently demanded in ``explain'' questions.)
\end{propriete}

\begin{attention}[Do not confuse the two trends]
Down Group VII, the \emph{oxidising power of the halogens} \ce{X2} decreases, but the \emph{acid strength} and the \emph{reducing power} of the hydrogen halides \ce{HX} increase. Writing ``\ce{HF} is the strongest acid because fluorine is the most reactive halogen'' is a serious error that costs marks every year.
\end{attention}

\textbf{Reducing power of the hydrogen halides}\par

The ease with which \ce{X-} loses electrons increases down the group, so the reducing power increases: $\ce{HF} < \ce{HCl} < \ce{HBr} < \ce{HI}$. This is exactly what we exploited above: \ce{HI} reduces concentrated \ce{H2SO4} all the way to \ce{H2S} (sulphur reduced from $+6$ to $-2$), \ce{HBr} reduces it only to \ce{SO2} ($+6$ to $+4$), while \ce{HCl} and \ce{HF} do not reduce it at all.

\courssub{Qualitative Tests for Halide Ions}

Identifying \ce{Cl-}, \ce{Br-} and \ce{I-} in an unknown solution is a compulsory practical skill. The key reagent is \cle{silver nitrate solution}, \ce{AgNO3(aq)}, because the three silver halides have characteristic colours and different solubilities in ammonia solution.

\begin{methode}[Testing for halide ions]
\begin{enumerate}
\item To about $2\ \mathrm{cm^3}$ of the unknown solution, add dilute nitric acid, \ce{HNO3}, to destroy interfering ions such as \ce{CO3^{2-}} (which would also give a precipitate with \ce{Ag+}).
\item Add a few drops of \ce{AgNO3(aq)}. A precipitate indicates a halide:
\begin{itemize}
\item \textbf{white} precipitate $\Rightarrow$ \ce{Cl-}: \quad $\ce{Ag+(aq) + Cl-(aq) -> AgCl(s)}$
\item \textbf{cream} precipitate $\Rightarrow$ \ce{Br-}: \quad $\ce{Ag+(aq) + Br-(aq) -> AgBr(s)}$
\item \textbf{yellow} precipitate $\Rightarrow$ \ce{I-}: \quad $\ce{Ag+(aq) + I-(aq) -> AgI(s)}$
\end{itemize}
\item Confirm by adding ammonia solution: \ce{AgCl} dissolves in \textbf{dilute} \ce{NH3(aq)}, \ce{AgBr} dissolves only in \textbf{concentrated} \ce{NH3(aq)}, and \ce{AgI} is \textbf{insoluble} even in concentrated ammonia. Dissolution is due to formation of the soluble complex $\ce{[Ag(NH3)2]+}$.
\end{enumerate}
\end{methode}

\begin{popfigure}
\begin{center}
\begin{tikzpicture}[
  font=\small,
  box/.style={draw=popblue, thick, rounded corners=2pt, fill=popblueL, text width=5.2cm, align=center, inner sep=4pt},
  res/.style={draw=popdark, thick, rounded corners=2pt, fill=white, text width=3.2cm, align=center, inner sep=4pt},
  arr/.style={-{Stealth[length=2.5mm]}, thick, popdark}]
\node[box, align=center] (start) at (0,0){Unknown solution\\ $+$ dilute \ce{HNO3}, then \ce{AgNO3(aq)}};
\node[box] (ppt) at (0,-2.2) {Precipitate formed?};
\node[res, align=center] (no) at (5.2,-2.2){No precipitate:\\ halide absent};
\node[res, fill=poppink!15, draw=poppink, align=center] (cl) at (-5.2,-4.6){White ppt \ce{AgCl}\\ dissolves in \textbf{dilute} \ce{NH3}};
\node[res, fill=popgold!20, draw=popgold, align=center] (br) at (0,-4.6){Cream ppt \ce{AgBr}\\ dissolves in \textbf{conc.} \ce{NH3}};
\node[res, fill=poporange!20, draw=poporange, align=center] (i) at (5.2,-4.6){Yellow ppt \ce{AgI}\\ \textbf{insoluble} in \ce{NH3}};
\node[res, draw=popgreen, fill=popgreenL] (c1) at (-5.2,-6.6) {\ce{Cl-} confirmed};
\node[res, draw=popgreen, fill=popgreenL] (c2) at (0,-6.6) {\ce{Br-} confirmed};
\node[res, draw=popgreen, fill=popgreenL] (c3) at (5.2,-6.6) {\ce{I-} confirmed};
\draw[arr] (start) -- (ppt);
\draw[arr] (ppt) -- (no) node[midway, above]{no};
\draw[arr] (ppt) -- (cl) node[midway, above left]{white};
\draw[arr] (ppt) -- (br) node[midway, right]{cream};
\draw[arr] (ppt) -- (i) node[midway, above right]{yellow};
\draw[arr] (cl) -- (c1);
\draw[arr] (br) -- (c2);
\draw[arr] (i) -- (c3);
\end{tikzpicture}
\end{center}
\legende{Flowchart for the identification of \ce{Cl-}, \ce{Br-} and \ce{I-} using silver nitrate and ammonia solutions.}
\end{popfigure}

\textbf{Confirmatory tests}\par

\begin{itemize}
\item \textbf{Concentrated \ce{H2SO4} on the solid salt.} A chloride gives only misty fumes of \ce{HCl}; a bromide gives misty fumes plus brown \ce{Br2} vapour; an iodide gives violet \ce{I2} vapour and the rotten-egg smell of \ce{H2S}. The deeper the oxidation, the heavier the halide.
\item \textbf{Test for chlorine gas.} Chlorine \emph{bleaches damp litmus paper} (it turns red first, then white), because \ce{Cl2} forms the bleaching agent \ce{HOCl} with the moisture: $\ce{Cl2(g) + H2O(l) -> HCl(aq) + HOCl(aq)}$. Note: the paper must be \emph{damp}; dry chlorine does not bleach.
\end{itemize}

\begin{exoresolu}[Identifying an unknown halide]
A white solid P dissolves in water. Addition of dilute \ce{HNO3} followed by \ce{AgNO3(aq)} gives a cream precipitate which does not dissolve in dilute ammonia but dissolves in concentrated ammonia. When concentrated \ce{H2SO4} is added to solid P, brown fumes are evolved. Identify the anion in P and write equations for the reactions observed.
\tcblower
\textbf{Solution.} The cream precipitate with \ce{AgNO3} suggests \ce{AgBr}, hence the anion is the bromide ion, \ce{Br-}. Confirmation: \ce{AgBr} is insoluble in dilute \ce{NH3} but soluble in concentrated \ce{NH3} (formation of $\ce{[Ag(NH3)2]+}$), exactly as observed. The brown fumes with concentrated \ce{H2SO4} are \ce{Br2}, formed by oxidation of \ce{Br-}.\\[2pt]
Equations:
\begin{align*}
\ce{Ag+(aq) + Br-(aq) &-> AgBr(s)}\\
\ce{AgBr(s) + 2NH3(aq) &-> [Ag(NH3)2]+(aq) + Br-(aq)}\\
\ce{2Br- + 2H2SO4 &-> Br2 + SO2 + SO4^{2-} + 2H2O}
\end{align*}
Conclusion: P contains the \textbf{bromide ion}, e.g.\ \ce{NaBr} or \ce{KBr}.
\end{exoresolu}

\courssub{Oxoacids of the Halogens}

\begin{definition}[Oxoacid]
An \cle{oxoacid} is an acid containing hydrogen, oxygen and another element. The halogen oxoacids have the general formula \ensuremath{\mathrm{HXO_{n}}} with $n = 1$ to $4$, in which the halogen has a \emph{positive} oxidation state.
\end{definition}

For chlorine the four oxoacids are:

\begin{center}
\begin{tabular}{|l|l|c|c|}
\hline
\rowcolor{popblueL} \textbf{Formula} & \textbf{Name} & \textbf{Oxidation state of Cl} & \textbf{Salt name} \\
\hline
\ce{HOCl} (or \ce{HClO}) & hypochlorous acid & $+1$ & hypochlorite, \ce{ClO-} \\
\hline
\ce{HClO2} & chlorous acid & $+3$ & chlorite, \ce{ClO2-} \\
\hline
\ce{HClO3} & chloric acid & $+5$ & chlorate, \ce{ClO3-} \\
\hline
\ce{HClO4} & perchloric acid & $+7$ & perchlorate, \ce{ClO4-} \\
\hline
\end{tabular}
\end{center}

The naming pattern is systematic: the \emph{lowest} oxoacid takes the prefix \emph{hypo-} and suffix \emph{-ous}; the next is simply \emph{-ous}; then \emph{-ic}; and the highest takes \emph{per-} and \emph{-ic}. Bromine and iodine form analogous series (e.g.\ \ce{HBrO3} bromic acid, \ce{HIO4} periodic acid), though not all members are stable.

\begin{popfigure}
\begin{center}
\begin{tikzpicture}[font=\small,
  atom/.style={circle, draw=popdark, thick, minimum size=7mm, inner sep=1pt},
  cl/.style={atom, fill=popgreenL, draw=popgreen},
  o/.style={atom, fill=poppink!20, draw=poppink},
  h/.style={atom, fill=popblueL, draw=popblue, minimum size=5.5mm},
  bond/.style={thick, popdark}]
% HOCl
\begin{scope}[shift={(0,0)}]
\node[h] (h1) at (0,0) {H};
\node[o] (o1) at (1.1,0) {O};
\node[cl] (c1) at (2.2,0) {Cl};
\draw[bond] (h1)--(o1)--(c1);
\node[below=4mm of c1, align=center] {\ce{HOCl}\\ hypochlorous};
\end{scope}
% HClO2
\begin{scope}[shift={(4.2,0)}]
\node[h] (h2) at (0,0) {H};
\node[o] (o2) at (1.1,0) {O};
\node[cl] (c2) at (2.2,0) {Cl};
\node[o] (o2b) at (2.2,1.1) {O};
\draw[bond] (h2)--(o2)--(c2)--(o2b);
\node[below=4mm of c2, align=center] {\ce{HClO2}\\ chlorous};
\end{scope}
% HClO3
\begin{scope}[shift={(8.4,0)}]
\node[h] (h3) at (0,0) {H};
\node[o] (o3) at (1.1,0) {O};
\node[cl] (c3) at (2.2,0) {Cl};
\node[o] (o3b) at (2.2,1.1) {O};
\node[o] (o3c) at (3.3,0) {O};
\draw[bond] (h3)--(o3)--(c3)--(o3b);
\draw[bond] (c3)--(o3c);
\node[below=4mm of c3, align=center] {\ce{HClO3}\\ chloric};
\end{scope}
% HClO4
\begin{scope}[shift={(12.6,0)}]
\node[h] (h4) at (0,0) {H};
\node[o] (o4) at (1.1,0) {O};
\node[cl] (c4) at (2.2,0) {Cl};
\node[o] (o4b) at (2.2,1.1) {O};
\node[o] (o4c) at (3.3,0) {O};
\node[o] (o4d) at (2.2,-1.1) {O};
\draw[bond] (h4)--(o4)--(c4)--(o4b);
\draw[bond] (c4)--(o4c);
\draw[bond] (c4)--(o4d);
\node[below=8mm of c4, align=center] {\ce{HClO4}\\ perchloric};
\end{scope}
\end{tikzpicture}
\end{center}
\legende{Structures of the four oxoacids of chlorine: the acidic hydrogen is always bonded to an oxygen atom, and the number of oxygen atoms coordinated to chlorine increases from one to four. (Double bonds to oxygen are omitted for clarity; the actual geometry around Cl is tetrahedral.)}
\end{popfigure}

\begin{attention}[Where is the acidic hydrogen?]
In \emph{all} halogen oxoacids the ionisable hydrogen is attached to \textbf{oxygen} (\ce{H-O-X}), never directly to the halogen. Writing \ce{HClO} with the H bonded to Cl is a structural error.
\end{attention}

\courssub{Trends in the Oxoacids: Acid Strength, Oxidising Strength and Thermal Stability}

\begin{propriete}[Trends in the halogen oxoacids]
For oxoacids of the \emph{same} halogen:
\begin{itemize}
\item \textbf{Acid strength increases} with the number of oxygen atoms:
\[
\ce{HOCl} < \ce{HClO2} < \ce{HClO3} < \ce{HClO4}
\]
\ce{HOCl} is a weak acid ($K_{\mathrm{a}} \approx 3.0 \times 10^{-8}$) while \ce{HClO4} is one of the strongest acids known. Each additional oxygen atom withdraws electron density from the \ce{O-H} bond (through the electronegative chlorine), weakening it and stabilising the conjugate anion by delocalising the negative charge over more oxygen atoms.
\item \textbf{Oxidising strength decreases} as the number of oxygen atoms increases: \ce{HOCl} is the most powerful oxidising agent of the series, \ce{HClO4} (when cold and dilute) the least. The lower oxoacids contain chlorine in a low positive oxidation state with few stabilising oxygens, so they gain electrons readily.
\item \textbf{Thermal stability increases} with the number of oxygen atoms: \ce{HOCl} decomposes easily ($\ce{2HOCl -> 2HCl + O2}$, accelerated by light), whereas \ce{HClO4} is the most stable, though it becomes a dangerously powerful oxidant when hot and concentrated.
\end{itemize}
\end{propriete}

\begin{aretenir}[Exam tip: the oxoacid mnemonic]
``\textbf{More oxygens, stronger acid.}'' Acid strength and thermal stability rise from \ce{HOX} to \ce{HXO4}, while oxidising strength falls. The lower oxoacids (\ce{HOCl}, \ce{HClO2}) are the strongest oxidising agents; the higher ones (\ce{HClO4}) are the strongest acids. Learn the two opposite trends together so you never swap them.
\end{aretenir}

\textbf{Uses of the oxoacids and their salts}\par

\begin{itemize}
\item \textbf{Sodium hypochlorite}, \ce{NaClO}, is the active ingredient of household \emph{bleach} and of many water-treatment products; it kills germs precisely because \ce{ClO-} is a strong oxidising agent. In Cameroon, dilute hypochlorite solutions are routinely used to disinfect drinking water.
\item \textbf{Chlorates}, e.g.\ \ce{KClO3}, are used in \emph{match heads} and \emph{fireworks}, where they supply oxygen for rapid combustion.
\item \textbf{Perchlorates}, e.g.\ \ce{KClO4} and \ce{NH4ClO4}, are used in \emph{explosives}, pyrotechnics and solid rocket propellants.
\end{itemize}

\begin{savaistu}[Bleach and the chemistry of clean water]
The bleaching power of ``eau de Javel'' sold in every Cameroonian market comes from the hypochlorite ion \ce{ClO-}, the conjugate base of hypochlorous acid \ce{HOCl} — the very species formed when chlorine dissolves in water. The same chemistry that bleaches damp litmus paper in your test for chlorine disinfects wells and hospital surfaces. Never mix bleach with acids: the reaction $\ce{ClO- + Cl- + 2H+ -> Cl2 + H2O}$ releases toxic chlorine gas!
\end{savaistu}

\begin{exoresolu}[Oxidation states and acid strength]
(a) Determine the oxidation state of chlorine in \ce{HClO}, \ce{HClO2}, \ce{HClO3} and \ce{HClO4}. (b) Arrange the acids in order of increasing acid strength and justify your answer. (c) Which of them is the strongest oxidising agent?
\tcblower
\textbf{Solution.} (a) Using $x + 1 + n(-2) = 0$ for \ensuremath{\mathrm{HClO_{n}}}:
\begin{align*}
\ce{HClO}:&\quad x = +1 & \ce{HClO2}:&\quad x = +3\\
\ce{HClO3}:&\quad x = +5 & \ce{HClO4}:&\quad x = +7
\end{align*}
(b) Increasing acid strength: $\ce{HClO} < \ce{HClO2} < \ce{HClO3} < \ce{HClO4}$. Justification: each extra oxygen atom withdraws electron density from the \ce{O-H} bond, weakening it, and stabilises the anion \ensuremath{\mathrm{ClO_{n}-}} by spreading the negative charge over more oxygen atoms.\\[2pt]
(c) \ce{HClO} (hypochlorous acid) is the strongest oxidising agent: oxidising strength \emph{decreases} as the number of oxygen atoms increases.
\end{exoresolu}

\begin{exercice}[Preparation and identification]
(a) Explain, with equations, why pure \ce{HI} cannot be prepared by the action of concentrated \ce{H2SO4} on \ce{NaI}, and describe a suitable alternative preparation.\\
(b) A solution gives a yellow precipitate with \ce{AgNO3(aq)} after acidification with dilute \ce{HNO3}; the precipitate is insoluble in concentrated ammonia. Identify the ion present and write the precipitation equation.\\
(c) Arrange \ce{HBrO}, \ce{HBrO3} and \ce{HBrO4} in order of (i) increasing acid strength, (ii) increasing oxidising strength.
\end{exercice}

\begin{corrige}[Exercise 1]
(a) \ce{I-} is a strong reducing agent; concentrated \ce{H2SO4} oxidises it: $\ce{8HI + H2SO4 -> 4I2 + H2S + 4H2O}$, so the \ce{HI} is contaminated by \ce{I2}, \ce{SO2}, \ce{S} and \ce{H2S}. Alternative: use non-oxidising phosphoric(V) acid, $\ce{NaI + H3PO4 -> NaH2PO4 + HI}$, or hydrolyse \ce{PI3}: $\ce{PI3 + 3H2O -> H3PO3 + 3HI}$.\\
(b) The ion is \ce{I-}: $\ce{Ag+(aq) + I-(aq) -> AgI(s)}$ (yellow, insoluble in \ce{NH3}).\\
(c) (i) Acid strength: $\ce{HBrO} < \ce{HBrO3} < \ce{HBrO4}$. (ii) Oxidising strength: $\ce{HBrO4} < \ce{HBrO3} < \ce{HBrO}$ (the reverse order).
\end{corrige}

\begin{aretenir}[Key points of this section]
\begin{itemize}
\item \ce{HCl} and \ce{HF} are prepared from metal halides with concentrated \ce{H2SO4}; \ce{HBr} and \ce{HI} require \ce{H3PO4} or hydrolysis of \ce{PX3}, because \ce{H2SO4} oxidises \ce{Br-} and \ce{I-}.
\item Acid strength: $\ce{HF} < \ce{HCl} < \ce{HBr} < \ce{HI}$ (weakening \ce{H-X} bond); reducing power increases in the same order.
\item Halide test: dilute \ce{HNO3}, then \ce{AgNO3}; \ce{AgCl} white (soluble in dilute \ce{NH3}), \ce{AgBr} cream (soluble in conc.\ \ce{NH3}), \ce{AgI} yellow (insoluble). Chlorine bleaches damp litmus paper.
\item Oxoacids \ensuremath{\mathrm{HXO_{n}}}: hypochlorous ($+1$), chlorous ($+3$), chloric ($+5$), perchloric ($+7$). More oxygens $\Rightarrow$ stronger acid, more thermally stable, but weaker oxidising agent.
\item Uses: \ce{NaClO} in bleach, chlorates in matches and fireworks, perchlorates in explosives.
\end{itemize}
\end{aretenir}

\coursec{Group IV Elements: Introduction, Allotropy and Trends in Physical Properties}

In the previous section we completed our study of the halogens by examining the hydrogen halides, the tests for halide ions and the oxoacids of chlorine, bromine and iodine. We now leave Group VII and move two steps to the left of the periodic table, to \cle{Group IV} (also called Group 14). This group is remarkable for one reason above all: it contains \cle{carbon}, the element on which all life and all organic chemistry are built. Yet as we descend the group we meet silicon, the element of computer chips and glass, then germanium, tin and lead — and the character of the elements changes completely from non-metal to metal. Understanding \emph{why} this happens is the heart of this section.

\courssub{The Group IV Elements and Their Electronic Configuration}

Group IV occupies the central p-block column of the periodic table, between Group III (B, Al, \dots) and Group V (N, P, \dots). Its five elements are carbon (C), silicon (Si), germanium (Ge), tin (Sn) and lead (Pb).

\begin{definition}[Group IV elements]
The Group IV elements are the elements of the fourth main group of the periodic table: \ce{C}, \ce{Si}, \ce{Ge}, \ce{Sn} and \ce{Pb}. Every Group IV atom has the outer electronic configuration \cle{$ns^2np^2$}, i.e.\ four valence electrons.
\end{definition}

\begin{center}
\begin{tabularx}{\linewidth}{|l|c|X|}
\hline
\rowcolor{popblueL} \textbf{Element} & \textbf{Symbol} & \textbf{Outer electronic configuration} \\
\hline
Carbon & \ce{C} & $2s^2 2p^2$ \\
\hline
Silicon & \ce{Si} & $3s^2 3p^2$ \\
\hline
Germanium & \ce{Ge} & $4s^2 4p^2$ \\
\hline
Tin & \ce{Sn} & $5s^2 5p^2$ \\
\hline
Lead & \ce{Pb} & $6s^2 6p^2$ \\
\hline
\end{tabularx}
\end{center}

With four valence electrons, each atom can form \textbf{four covalent bonds} (as in \ce{CH4} or \ce{SiCl4}), giving the maximum oxidation state of $+4$. The alternative oxidation state, $+2$, appears when only the two $np$ electrons are used in bonding — we shall return to this crucial idea when we discuss the \cle{inert pair effect}.

\courssub{The Unique Properties of Carbon}

Carbon behaves very differently from the rest of its group. Four features explain this uniqueness:

\begin{itemize}
\item \textbf{Small atomic size.} Carbon is the smallest member of the group. Its atoms can approach each other closely, so its $2p$ orbitals overlap effectively \emph{sideways} to form strong $\pi$ bonds. Hence carbon readily forms \cle{multiple bonds}: \ce{C=C}, \ce{C#C}, \ce{C=O}. Silicon and the heavier elements cannot do this efficiently because their larger $3p$, $4p$ \dots\ orbitals overlap poorly sideways.
\item \textbf{Catenation.} Carbon shows \cle{catenation} — the ability to form long chains and rings of identical atoms — to a unique extent. The \ce{C-C} single bond is very strong (about \SI{348}{kJ.mol^{-1}}), stronger than \ce{Si-Si} (about \SI{226}{kJ.mol^{-1}}), so carbon chains are stable. This is why millions of organic compounds exist.
\item \textbf{High electronegativity.} Carbon ($\chi \approx 2.5$ on the Pauling scale) is the most electronegative member of the group, so it forms essentially covalent compounds and never forms a simple \ce{C^4+} ion.
\item \textbf{Maximum covalency of four.} Carbon has no available $d$ orbitals in its valence shell ($n=2$), so it can never expand its octet: \ce{CCl4} does not hydrolyse, whereas \ce{SiCl4} hydrolyses violently because silicon can use its empty $3d$ orbitals to accept a water molecule.
\end{itemize}

\begin{attention}[Common mistake]
Do not write that carbon forms \ce{C^4+} ions. The energy required to remove four electrons is far too large; carbon is \emph{always covalent} in its compounds.
\end{attention}

\courssub{Allotropy of Carbon: Diamond, Graphite and Fullerenes}

\begin{definition}[Allotropy]
\cle{Allotropy} is the existence of an element in two or more different structural forms, called \textbf{allotropes}, in the same physical state.
\end{definition}

Carbon provides the classic example of allotropy: \textbf{diamond} and \textbf{graphite} are both pure solid carbon, yet diamond is the hardest natural substance known while graphite is soft enough to write with. The difference lies entirely in the bonding arrangement.

\textbf{Diamond.} Each carbon atom is $sp^3$ hybridised and bonded \emph{tetrahedrally} to four neighbours, building a \cle{giant covalent} three-dimensional lattice. All four valence electrons are locked in $\sigma$ bonds. Consequently diamond is extremely hard, has a very high melting point (sublimes above \SI{3600}{\ensuremath{{}^{\circ}\mathrm{C}}}) and is an electrical \textbf{insulator}.

\textbf{Graphite.} Each carbon atom is $sp^2$ hybridised and bonded to three neighbours in flat hexagonal layers. The fourth electron of each atom occupies a $p$ orbital perpendicular to the layer, and these electrons merge into a \cle{delocalised} $\pi$ system \emph{between} the layers. The layers are held only by weak van der Waals forces, so they slide over one another: graphite is soft, slippery (a lubricant and pencil ``lead'') and conducts electricity along the layers.

\begin{propriete}[Electrical conductivity of graphite]
Graphite conducts electricity because each carbon atom contributes one \cle{delocalised electron} that moves freely within and between the layers; diamond, in which every valence electron is localised in a $\sigma$ bond, cannot conduct. (Proof not examinable; the explanation must be quoted.)
\end{propriete}

\begin{popfigure}
\begin{center}
\begin{tikzpicture}[scale=0.85]
% ---- Diamond ----
\begin{scope}
\coordinate (A) at (0,0); \coordinate (B) at (1.3,0.75);
\coordinate (C) at (2.6,0); \coordinate (D) at (1.3,-0.9);
\coordinate (E) at (1.3,1.7);
\foreach \p in {A,B,C,D,E} \fill[popink] (\p) circle (0.09);
\draw[popdark,thick] (A)--(B) (B)--(C) (A)--(D) (D)--(C) (B)--(E);
\draw[popdark,thick] (A)--(1.3,0.1);
\node[popdark,font=\small\bfseries] at (1.3,-1.5) {Diamond};
\node[popmuted,font=\footnotesize,align=center,text width=3.4cm] at (1.3,-2.3) {tetrahedral 3D lattice,\\ 4 $\sigma$ bonds per C};
\end{scope}
% ---- Graphite ----
\begin{scope}[xshift=6.2cm]
\foreach \y in {0,1.1}{
  \foreach \x in {0,0.9,1.8,2.7}{
    \fill[popblue] (\x,\y) circle (0.09);
    \fill[popblue] (\x+0.45,\y+0.55) circle (0.09);
  }
  \foreach \x in {0,0.9,1.8}{
    \draw[popblue,thick] (\x,\y)--(\x+0.45,\y+0.55)--(\x+0.9,\y);
  }
  \draw[popblue,thick] (2.7,\y)--(3.15,\y+0.55);
  \draw[popblue,thick] (0.45,\y+0.55)--(0.9,\y+1.1*0);
}
\foreach \x in {0.45,1.35,2.25,3.15}{
  \draw[popblue,thick] (\x,0.55)--(\x,1.65);
}
\foreach \x in {0,0.9,1.8,2.7}{
  \fill[popblue] (\x,1.1) circle (0.09);
  \fill[popblue] (\x+0.45,1.65) circle (0.09);
}
\foreach \x in {0,0.9,1.8}{
  \draw[popblue,thick] (\x,1.1)--(\x+0.45,1.65)--(\x+0.9,1.1);
}
\draw[popblue,thick] (2.7,1.1)--(3.15,1.65);
\draw[poporange,dashed,thick] (3.6,0.2)--(3.6,1.5);
\node[poporange,font=\footnotesize,align=center] at (4.35,0.85) {weak\\ forces};
\node[popblue,font=\small\bfseries] at (1.6,-1.5) {Graphite};
\node[popmuted,font=\footnotesize,align=center,text width=3.6cm] at (1.6,-2.3) {hexagonal layers, 3 bonds per C,\\ one delocalised electron};
\end{scope}
\end{tikzpicture}
\end{center}
\legende{Structures of diamond (left) and graphite (right). Diamond: every C forms four $\sigma$ bonds in a rigid tetrahedral network. Graphite: every C forms three $\sigma$ bonds in layers; the fourth electron is delocalised, giving conductivity and softness.}
\end{popfigure}

\begin{savaistu}[Fullerenes — an exam-useful extra]
In 1985 a third allotrope family was discovered: the \cle{fullerenes}, molecules such as \ce{C60} (``buckyballs'') in which carbon atoms form closed cages of pentagons and hexagons, like a football. Fullerenes and the related \textbf{carbon nanotubes} are studied for drug delivery and super-strong materials. Mentioning them in an essay on allotropy earns credit for breadth of knowledge.
\end{savaistu}

\courssub{Allotropy in Tin and the Metallic Character Trend}

Carbon is not the only allotropic member. \textbf{Tin} exists as \cle{grey tin} ($\alpha$-tin, a brittle non-metallic form with a diamond-type structure) stable below \SI{13.2}{\ensuremath{{}^{\circ}\mathrm{C}}}, and \cle{white tin} ($\beta$-tin, the familiar metallic form) stable above that temperature. The slow crumbling of tin objects in severe cold (``tin pest'') is a real consequence of this transition.

This dual personality of tin illustrates the general trend: \textbf{metallic character increases down the group}. Carbon is a non-metal; silicon and germanium are \cle{metalloids} (semiconductors); tin and lead are metals, with positive electrode potentials and basic oxides in their lower oxidation state.

\courssub{Trend in Structure: From Giant Covalent to Metallic}

The change of character down the group is rooted in a change of \textbf{structure}:

\begin{center}
\begin{tabularx}{\linewidth}{|l|X|X|}
\hline
\rowcolor{popblueL} \textbf{Element} & \textbf{Structure of the element} & \textbf{Character} \\
\hline
\ce{C} (diamond) & Giant covalent, tetrahedral & Non-metal \\
\hline
\ce{Si}, \ce{Ge} & Giant covalent, diamond-type & Metalloids (semiconductors) \\
\hline
\ce{Sn} & Giant covalent (grey) or metallic (white) & Metal \\
\hline
\ce{Pb} & Metallic lattice & Metal \\
\hline
\end{tabularx}
\end{center}

As the atoms grow larger, the \ce{E-E} covalent bonds become longer and weaker, and the giant covalent lattice loses out to a metallic lattice in which the valence electrons are delocalised over positive ions.

\courssub{Trends in Atomic Radius, Ionization Energy, Electronegativity and Melting Point}

\begin{propriete}[Trends down Group IV]
Descending Group IV from C to Pb:
\begin{itemize}
\item \textbf{Atomic radius increases} — each element has one more filled electron shell.
\item \textbf{First ionization energy decreases} — the outer electrons are further from the nucleus and more shielded.
\item \textbf{Electronegativity decreases} (C 2.5 to Pb about 1.8) — bonding electrons are attracted less strongly.
\item \textbf{Melting point decreases overall} — but the explanation must be structural (see warning below).
\end{itemize}
\end{propriete}

\begin{popfigure}
\begin{center}
\begin{tikzpicture}
\begin{axis}[
  width=12cm, height=7cm,
  xlabel={Element}, ylabel={First ionization energy / $\mathrm{kJ\,mol^{-1}}$},
  symbolic x coords={C,Si,Ge,Sn,Pb}, xtick=data,
  ymin=600, ymax=1150, grid=major, grid style={popline!40},
  legend style={at={(0.97,0.97)},anchor=north east,font=\small},
  axis y line*=left, axis x line*=bottom]
\addplot[mark=*,popblue,thick] coordinates {(C,1086) (Si,786) (Ge,762) (Sn,709) (Pb,716)};
\addlegendentry{First ionization energy}
\end{axis}
\begin{axis}[
  width=12cm, height=7cm,
  symbolic x coords={C,Si,Ge,Sn,Pb}, xtick=data,
  ymin=60, ymax=140, hide x axis,
  axis y line*=right, ylabel={Atomic radius / pm},
  legend style={at={(0.97,0.75)},anchor=north east,font=\small}]
\addplot[mark=square*,poporange,thick,dashed] coordinates {(C,77) (Si,118) (Ge,122) (Sn,140) (Pb,135)};
\addlegendentry{Atomic radius}
\end{axis}
\end{tikzpicture}
\end{center}
\legende{First ionization energy (blue, left axis) falls while atomic radius (orange, right axis) rises from C to Pb. Note the slight irregularity at Pb — a reminder that real trends are smooth only to first approximation.}
\end{popfigure}

\begin{attention}[The melting point trap]
Do \textbf{not} write simply that ``melting points decrease down Group IV'' without explanation. Carbon (as diamond or graphite, subliming above \SI{3600}{\ensuremath{{}^{\circ}\mathrm{C}}}) has an \emph{exceptionally high} melting point because melting requires breaking strong covalent bonds throughout a giant lattice. Silicon (\SI{1410}{\ensuremath{{}^{\circ}\mathrm{C}}}) and germanium (\SI{937}{\ensuremath{{}^{\circ}\mathrm{C}}}) are lower because their \ce{E-E} bonds are longer and weaker. Tin (\SI{232}{\ensuremath{{}^{\circ}\mathrm{C}}}) and lead (\SI{327}{\ensuremath{{}^{\circ}\mathrm{C}}}) are metallic, and their metallic bonding is weak because the atoms are large. Always explain the trend \emph{via the change of structure}, from giant covalent to metallic.
\end{attention}

\courssub{Oxidation States $+2$ and $+4$: The Inert Pair Effect}

Group IV elements show oxidation states $+2$ and $+4$. At the top of the group, $+4$ dominates: \ce{CO2}, \ce{CCl4}, \ce{SiO2} are all stable, while \ce{CO} is easily oxidised. At the bottom, the picture reverses: \ce{PbO} and \ce{PbCl2} are stable, whereas \ce{PbO2} and \ce{Pb(IV)} compounds are strong \textbf{oxidizing agents}, eager to be reduced to \ce{Pb(II)}.

\begin{definition}[The inert pair effect]
The \cle{inert pair effect} is the reluctance of the $ns^2$ pair of valence electrons to participate in bonding. Its importance increases down Groups III, IV and V, so the lower oxidation state ($+2$ for Group IV) becomes progressively more stable relative to the higher one ($+4$).
\end{definition}

The explanation is energetic: down the group the \ce{E-E} and \ce{E-X} bonds become weaker, so the energy released by forming two extra bonds no longer compensates the energy needed to involve the $ns^2$ pair. Hence:

\begin{itemize}
\item Stability of $+4$ decreases: $\ce{C} > \ce{Si} > \ce{Ge} > \ce{Sn} > \ce{Pb}$.
\item Stability of $+2$ increases: \ce{Pb^2+} is more stable than \ce{Pb^4+}; \ce{Sn^2+} is a mild reducing agent, but \ce{Pb(IV)} is a strong oxidizing agent.
\end{itemize}

\begin{exemplebox}[The inert pair effect in action]
Lead(IV) oxide oxidises concentrated hydrochloric acid to chlorine, being itself reduced to lead(II):
\[ \ce{PbO2(s) + 4HCl(aq) -> PbCl2(aq) + Cl2(g) + 2H2O(l)} \]
This is the reaction used practically to distinguish \ce{PbO2} from \ce{PbO}: only the Pb(IV) oxide liberates chlorine.
\end{exemplebox}

\begin{exoresolu}[Explaining the oxidizing power of Pb(IV)]
Explain why \ce{PbO2} is a good oxidizing agent whereas \ce{CO2} is not.
\tcblower
\textbf{Step 1 — Identify the oxidation states.} In \ce{PbO2}, lead is in the $+4$ state; in \ce{CO2}, carbon is in the $+4$ state.

\textbf{Step 2 — Apply the inert pair effect.} For lead, at the bottom of Group IV, the $6s^2$ pair is reluctant to bond (inert pair effect), so the $+2$ state is more stable than $+4$. \ce{PbO2} therefore readily gains electrons and is reduced to \ce{Pb(II)}:
\[ \ce{PbO2 + 4H+ + 2e- -> Pb^{2+} + 2H2O} \]
This makes \ce{PbO2} a strong oxidizing agent.

\textbf{Step 3 — Contrast with carbon.} Carbon, at the top of the group, shows no inert pair effect: $+4$ is its most stable state, so \ce{CO2} has no tendency to be reduced and is not an oxidizing agent under ordinary conditions.

\textbf{Conclusion.} The oxidizing power of Pb(IV) arises from the stability of Pb(II), a direct consequence of the inert pair effect.
\end{exoresolu}

\begin{methode}[Exam tip: answering ``why is Pb(IV) oxidizing?'']
When asked why Pb(IV) compounds are oxidizing agents:
\begin{enumerate}
\item State that the stability of $+2$ relative to $+4$ increases down Group IV.
\item Name the cause: the \cle{inert pair effect} — the $6s^2$ electrons of lead are reluctant to participate in bonding.
\item Conclude that Pb(IV) readily accepts electrons to form the more stable Pb(II), hence it oxidizes other species. Quote an example such as \ce{PbO2} with concentrated \ce{HCl}.
\end{enumerate}
\end{methode}

\courssub{Uses of the Group IV Elements}

The physical trends we have just studied translate directly into everyday applications, several of them familiar in Cameroon:

\begin{itemize}
\item \textbf{Carbon:} diamond in cutting and drilling tools; graphite in pencil ``leads'', electrodes and lubricants; coke as a reducing agent in metallurgy.
\item \textbf{Silicon:} purified silicon is the \cle{semiconductor} at the heart of computer chips and solar panels; silicon(IV) oxide, \ce{SiO2}, is the raw material of \textbf{glass}.
\item \textbf{Germanium:} used in semiconductors and in infrared optics.
\item \textbf{Tin:} an unreactive, non-toxic metal used to \textbf{plate} steel food cans (``tin cans'') and, alloyed with lead, as \textbf{solder} in electronics.
\item \textbf{Lead:} dense and corrosion-resistant, used in \textbf{car batteries} (lead--acid accumulators), in roofing sheets, and as shielding against X-rays; its compounds were formerly used in paints and petrol, uses now abandoned because of lead's toxicity.
\end{itemize}

\begin{aretenir}[Key points]
\begin{itemize}
\item Group IV elements (\ce{C}, \ce{Si}, \ce{Ge}, \ce{Sn}, \ce{Pb}) all have the outer configuration $ns^2np^2$.
\item Carbon is unique: small size, high electronegativity, strong catenation and multiple bonding, maximum covalency of four.
\item Carbon shows allotropy: diamond (giant tetrahedral, hard, insulator) and graphite (layered, soft, conductor thanks to delocalised electrons); fullerenes are a modern third form.
\item Structure changes from giant covalent (C, Si, Ge) to metallic (Sn, Pb); metallic character increases down the group.
\item Atomic radius increases; first ionization energy, electronegativity and melting point decrease — the melting point trend must be explained by the change of structure.
\item Oxidation states $+2$ and $+4$: the inert pair effect makes $+2$ increasingly stable down the group, so Pb(IV) compounds are oxidizing agents.
\end{itemize}
\end{aretenir}

\begin{exercice}[Checking your understanding]
\begin{enumerate}
\item Explain why \ce{CCl4} is not hydrolysed by water whereas \ce{SiCl4} is hydrolysed rapidly.
\item Account for the fact that graphite is used as an electrode material while diamond is not.
\item Arrange \ce{CO}, \ce{SnO}, \ce{PbO2} and \ce{PbO} in order of increasing stability and justify your answer using the inert pair effect.
\item Explain why the melting point of lead (\SI{327}{\ensuremath{{}^{\circ}\mathrm{C}}}) is far lower than that of carbon, even though lead is much lower in the group.
\end{enumerate}
\end{exercice}

\begin{corrige}[Exercise 1]
\textbf{1.} Carbon has no $d$ orbitals in its $n=2$ valence shell, so it cannot expand its octet and cannot accept a lone pair from water; \ce{CCl4} is inert. Silicon has empty $3d$ orbitals, accepts a water molecule as a ligand, and \ce{SiCl4} hydrolyses: $\ce{SiCl4(l) + 2H2O(l) -> SiO2(s) + 4HCl(aq)}$.

\textbf{2.} In graphite each carbon uses only three electrons in $\sigma$ bonds; the fourth is delocalised between the layers and carries current. In diamond all four electrons are localised in $\sigma$ bonds, so no charge carriers exist and diamond is an insulator.

\textbf{3.} Increasing stability: $\ce{PbO2} < \ce{CO} < \ce{SnO} < \ce{PbO}$. The inert pair effect makes the $+2$ state more stable down the group: Pb(II) oxide is the most stable, while Pb(IV) oxide is the least stable (a strong oxidizing agent); for carbon the $+4$ state is favoured, so \ce{CO} (C in $+2$) is readily oxidised.

\textbf{4.} Carbon is a giant covalent structure: melting requires breaking very strong \ce{C-C} bonds throughout the lattice. Lead has a metallic lattice with large atoms and weak metallic bonding, so far less energy is needed; hence its much lower melting point.
\end{corrige}

\coursec{Reactions and Compounds of Group IV: Hydrides, Oxides and Chlorides}

In the previous section we surveyed the \emph{physical} trends across Group~IV — atomic radii, ionisation energies, allotropy and the change from non‑metallic carbon to metallic lead. We now turn to their \emph{chemical} behaviour: how the elements react with oxygen, acids and alkalis, and the nature of the hydrides, oxides and chlorides they form. Understanding these reactions is essential for predicting the outcome of qualitative analysis and for explaining the contrasting properties of, for example, gaseous \ce{CO2} versus the refractory solid \ce{SiO2}.

\courssub{Reactions of the Group IV Elements}

\begin{experience}[Reactions with oxygen]
When the Group~IV elements are heated in air or oxygen they form oxides. The conditions and products vary down the group:
\begin{itemize}
\item \textbf{Carbon} burns in excess oxygen to give carbon dioxide; in a limited supply it forms carbon monoxide:
\[
\ensuremath{\mathrm{C(s) + O_{2}(g) \rightarrow [\Delta] CO_{2}(g)}} \qquad
\ensuremath{\mathrm{2C(s) + O_{2}(g) \rightarrow [\Delta] 2CO(g)}}
\]
\item \textbf{Silicon} requires strong heating (≈ 1000 °C) to react, yielding silicon dioxide:
\[
\ensuremath{\mathrm{Si(s) + O_{2}(g) \rightarrow [\Delta] SiO_{2}(s)}}
\]
\item \textbf{Germanium} behaves similarly to silicon, forming \ce{GeO2}.
\item \textbf{Tin} and \textbf{lead} are more reactive. Tin gives tin(IV) oxide \ce{SnO2} on strong heating; lead forms a mixture of lead(II) oxide \ce{PbO} and lead(II,IV) oxide \ce{Pb3O4} (red lead) depending on temperature and oxygen pressure:
\[
\ensuremath{\mathrm{2Pb(s) + O_{2}(g) \rightarrow [\Delta] 2PbO(s)}} \qquad
\ensuremath{\mathrm{3Pb(s) + 2O_{2}(g) \rightarrow [\Delta] Pb_{3}O_{4}(s)}}
\]
\end{itemize}
\end{experience}

\begin{experience}[Reactions with acids]
\begin{itemize}
\item \textbf{Carbon} and \textbf{silicon} are resistant to dilute non‑oxidising acids (e.g. \ce{HCl}, dilute \ce{H2SO4}) because of their strong covalent network structures.
\item \textbf{Tin} and \textbf{lead} react with \emph{oxidising} acids. Concentrated nitric acid oxidises Sn to \ce{SnO2} (or \ce{Sn^{4+}}) and Pb to \ce{PbO2} (or \ce{Pb^{4+}}), with evolution of \ce{NO2}. Concentrated sulfuric acid (hot) also acts as an oxidising agent:
\[
\ce{Sn(s) + 4HNO3(conc) -> SnO2(s) + 4NO2(g) + 2H2O(l)}
\]
\[
\ce{Pb(s) + 4HNO3(conc) -> PbO2(s) + 4NO2(g) + 2H2O(l)}
\]
\[
\ensuremath{\mathrm{Sn(s) + 2H_{2}SO_{4}(conc, \Delta) \rightarrow  SnO_{2}(s) + 2SO_{2}(g) + 2H_{2}O(l)}}
\]
\[
\ensuremath{\mathrm{Pb(s) + 2H_{2}SO_{4}(conc, \Delta) \rightarrow  PbO_{2}(s) + 2SO_{2}(g) + 2H_{2}O(l)}}
\]
\item Dilute \ce{HCl} reacts slowly with Sn and Pb, producing hydrogen and the corresponding \ce{M^{2+}} chlorides. The reaction with lead is slow due to the protective layer of insoluble \ce{PbCl2}:
\[
\ce{Sn(s) + 2HCl(aq) -> SnCl2(aq) + H2(g)}
\]
\[
\ce{Pb(s) + 2HCl(aq) -> PbCl2(s) + H2(g)} \quad\text{(slow)}
\]
\end{itemize}
\end{experience}

\begin{experience}[Reactions with alkalis]
\begin{itemize}
\item \textbf{Silicon} dissolves in hot concentrated \ce{NaOH} to give sodium silicate and hydrogen:
\[
\ce{Si(s) + 2NaOH(aq) + H2O(l) -> Na2SiO3(aq) + 2H2(g)}
\]
\item \textbf{Tin} and \textbf{lead} show \emph{amphoteric} behaviour: they dissolve in both acids and hot concentrated alkalis, forming hexahydroxostannate(IV) \ensuremath{\mathrm{[Sn(OH)_{6}]^{2-}}} and hexahydroxoplumbate(IV) \ensuremath{\mathrm{[Pb(OH)_{6}]^{2-}}} respectively (often written as stannates \ce{Na2SnO3} and plumbates \ce{Na2PbO3}).
\item Carbon does not react with alkalis under normal conditions.
\end{itemize}
\end{experience}

\begin{propriete}[Amphoteric oxides of Sn and Pb]
The oxides \ce{SnO}, \ce{SnO2}, \ce{PbO} and \ce{PbO2} react with both acids and bases. For example:
\[
\ce{SnO2(s) + 2H2SO4(aq) -> Sn(SO4)2(aq) + 2H2O(l)} \quad\text{(acidic medium)}
\]
\[
\ce{SnO2(s) + 2NaOH(aq) + 2H2O(l) -> Na2[Sn(OH)6](aq)} \quad\text{(basic medium)}
\]
This dual reactivity is a hallmark of the metallic side of Group~IV.
\end{propriete}

\courssub{Hydrides of Group IV}

All Group~IV elements form tetrahedral hydrides of general formula \ce{XH4} (X = C, Si, Ge, Sn, Pb). The structures are based on \emph{sp$^{3}$} hybridisation at the central atom.

\begin{definition}[Group IV hydrides]
The binary hydrides are:
\begin{itemize}
\item \ce{CH4} (methane) — gas, very stable
\item \ce{SiH4} (silane) — gas, pyrophoric
\item \ce{GeH4} (germane) — gas, less stable
\item \ce{SnH4} (stannane) — unstable, decomposes at room temperature
\item \ce{PbH4} (plumbane) — only observed at low temperature, extremely unstable
\end{itemize}
All adopt a tetrahedral geometry (bond angles ≈ 109.5°).
\end{definition}

\begin{propriete}[Thermal stability and reducing power of Group IV hydrides]
The thermal stability of the hydrides \emph{decreases} down the group as the \ce{X-H} bond becomes longer and weaker. Consequently, the reducing power \emph{increases} down the group: \ce{SnH4} and \ce{PbH4} are strong reducing agents, whereas \ce{CH4} is essentially inert.
\end{propriete}

\begin{aretenir}[Key trends in Group IV hydrides]
\begin{itemize}
\item Formula: \ce{XH4} (tetrahedral)
\item Thermal stability: \ce{CH4} > \ce{SiH4} > \ce{GeH4} > \ce{SnH4} > \ce{PbH4}
\item Reducing power: opposite order
\item Catenation: Carbon forms countless \ce{C-C} chains (hydrocarbons); silicon forms limited silanes (\ensuremath{\mathrm{Si_nH_{2n+2}}} up to n≈8); heavier analogues show negligible catenation.
\end{itemize}
\end{aretenir}

\begin{exemplebox}[Comparing \ce{CH4} and \ce{SiH4}]
Methane (\ce{CH4}) is the main component of natural gas; it burns cleanly to \ce{CO2} and \ce{H2O} and does not hydrolyse. Silane (\ce{SiH4}) ignites spontaneously in air (pyrophoric) and hydrolyses rapidly:
\[
\ce{SiH4(g) + 2H2O(l) -> SiO2(s) + 4H2(g)}
\]
The readiness of \ce{SiH4} to hydrolyse reflects the polarity of the \ce{Si-H} bond and the accessibility of silicon's empty 3d orbitals.
\end{exemplebox}

\courssub{Oxides of Group IV}

Group~IV oxides exist in two oxidation states: +2 (\ce{XO}) and +4 (\ce{XO2}). Their structures and acid–base characters change dramatically down the group.

\begin{propriete}[Structures of the dioxides]
\ce{CO2} consists of discrete linear \ce{O=C=O} molecules (gas at room temperature). \ce{SiO2}, \ce{GeO2} and \ce{SnO2} adopt giant covalent network structures based on \ce{[XO4]} tetrahedra sharing corners. \ce{PbO2} has a layered structure (rutile type) with \ce{Pb} in octahedral coordination.
\end{propriete}

\begin{popfigure}
\begin{tikzpicture}[scale=0.9]
% CO2 molecule
\node[align=center] at (0,3) {\textbf{\ce{CO2} (g)}};
\draw[thick, popdark] (-1.2,2.5) -- (0,2.8) -- (1.2,2.5);
\node at (-1.2,2.5) {O};
\node at (0,2.8) {C};
\node at (1.2,2.5) {O};
% SiO2 network fragment
\node[align=center] at (6,3) {\textbf{\ce{SiO2} (s)}};
% central Si
\node[circle,draw,fill=popblue!20,minimum size=6mm] (Si1) at (6,2) {Si};
% four O around
\foreach \ang/\pos in {45/NE,135/NW,225/SW,315/SE} {
  \node[circle,draw,fill=poporange!20,minimum size=5mm] (O\ang) at ($(Si1)+(\ang:1.2)$) {O};
  \draw[thick] (Si1) -- (O\ang);
}
% bridging O to second Si
\node[circle,draw,fill=popblue!20,minimum size=6mm] (Si2) at ($(O45)+(0:1.5)$) {Si};
\draw[thick] (O45) -- (Si2);
% third O on second Si
\node[circle,draw,fill=poporange!20,minimum size=5mm] (O2) at ($(Si2)+(0:1.2)$) {O};
\draw[thick] (Si2) -- (O2);
\end{tikzpicture}
\legende{Comparison of discrete molecular \ce{CO2} (left) and extended tetrahedral network of \ce{SiO2} (right).}
\end{popfigure}

\begin{propriete}[Acid–base character of Group IV oxides]
\begin{itemize}
\item \ce{CO2} and \ce{SiO2} are \emph{acidic} oxides. They react with bases to form carbonates/silicates.
\item \ce{CO} is a \emph{neutral} oxide — it does not react with either acids or bases under normal conditions.
\item \ce{SnO}, \ce{SnO2}, \ce{PbO}, \ce{PbO2} are \emph{amphoteric}.
\item \ce{GeO2} is weakly acidic; \ce{GeO} is amphoteric.
\end{itemize}
\end{propriete}

\begin{popfigure}
\begin{center}
\begin{tabular}{|c|c|c|}
\hline
\rowcolor{popblueL}
\textbf{Oxide} & \textbf{Formula} & \textbf{Acid–base character} \\
\hline
Carbon dioxide & \ce{CO2} & \cellcolor{popgreenL}Acidic \\
Silicon dioxide & \ce{SiO2} & \cellcolor{popgreenL}Acidic \\
Germanium dioxide & \ce{GeO2} & \cellcolor{popgreenL}Acidic (weak) \\
Tin(II) oxide & \ce{SnO} & \cellcolor{poporange}Amphoteric \\
Tin(IV) oxide & \ce{SnO2} & \cellcolor{poporange}Amphoteric \\
Lead(II) oxide & \ce{PbO} & \cellcolor{poporange}Amphoteric \\
Lead(IV) oxide & \ce{PbO2} & \cellcolor{poporange}Amphoteric \\
Carbon monoxide & \ce{CO} & \cellcolor{popblueL}Neutral \\
\hline
\end{tabular}
\end{center}
\legende{Acid–base classification of Group IV oxides (colour‑coded).}
\end{popfigure}

\begin{attention}[CO is neutral, not acidic]
A common error is to classify \ce{CO} as acidic because it is an oxide of carbon. \ce{CO} does \emph{not} react with water to form an acid, nor does it react with \ce{NaOH} or \ce{HCl}. Only \ce{CO2} is the acidic oxide of carbon.
\end{attention}

\begin{methode}[Classifying an oxide as acidic, basic or amphoteric]
\begin{enumerate}
\item Write the reaction of the oxide with a strong acid (e.g. \ce{HCl}). If a salt forms, the oxide shows basic behaviour.
\item Write the reaction with a strong base (e.g. \ce{NaOH}). If a salt forms, the oxide shows acidic behaviour.
\item If both reactions occur, the oxide is amphoteric. If neither occurs, it is neutral.
\end{enumerate}
\end{methode}

\begin{exemplebox}[Amphoteric behaviour of \ce{SnO2}]
\[
\ce{SnO2(s) + 4HCl(aq) -> SnCl4(aq) + 2H2O(l)} \quad\text{(basic behaviour)}
\]
\[
\ce{SnO2(s) + 2NaOH(aq) + 2H2O(l) -> Na2[Sn(OH)6](aq)} \quad\text{(acidic behaviour)}
\]
Thus \ce{SnO2} is amphoteric.
\end{exemplebox}

\courssub{Chlorides of Group IV}

The chlorides follow the oxidation states +2 and +4: \ce{XCl2} and \ce{XCl4}. Their stability and reactivity with water show a clear trend.

\begin{propriete}[Hydrolysis of tetrachlorides]
\ce{CCl4} is inert to water at room temperature. \ce{SiCl4}, \ce{GeCl4} and \ce{SnCl4} hydrolyse violently:
\[
\ce{SiCl4(l) + 2H2O(l) -> SiO2(s) + 4HCl(aq)}
\]
The reaction proceeds because silicon (and Ge, Sn) possess low‑lying empty \emph{d‑orbitals} (3d for Si) that can accept a lone pair from a water molecule, forming a pentacoordinate intermediate. Carbon lacks accessible d‑orbitals, so \ce{CCl4} cannot expand its coordination sphere and does not hydrolyse under normal conditions.
\end{propriete}

\begin{propriete}[Stability of the +4 chlorides down the group]
The +4 oxidation state becomes less stable relative to +2 due to the \emph{inert pair effect}. Consequently:
\begin{itemize}
\item \ce{CCl4}, \ce{SiCl4}, \ce{GeCl4}, \ce{SnCl4} are stable liquids (or solids).
\item \ce{PbCl4} is unstable; it decomposes readily to \ce{PbCl2} and chlorine:
\[
\ensuremath{\mathrm{PbCl_{4}(s) \rightarrow [\Delta] PbCl_{2}(s) + Cl_{2}(g)}}
\]
\end{itemize}
\end{propriete}

\begin{aretenir}[Trends in Group IV chlorides]
\begin{itemize}
\item Formulae: \ce{XCl2} and \ce{XCl4} (tetrahedral for \ce{XCl4}, bent/angular for \ce{XCl2} due to lone pair).
\item Hydrolysis: \ce{CCl4} inert; \ce{SiCl4}, \ce{GeCl4}, \ce{SnCl4} rapid; \ce{PbCl4} hydrolyses but decomposes first.
\item Thermal stability of \ce{XCl4}: decreases down the group; \ce{PbCl4} decomposes to \ce{PbCl2} + \ce{Cl2}.
\end{itemize}
\end{aretenir}

\begin{exemplebox}[Contrasting \ce{CCl4} and \ce{SiCl4} hydrolysis]
When water is added to \ce{CCl4}, two immiscible layers form — no reaction. With \ce{SiCl4}, white fumes of \ce{HCl} appear instantly and a white precipitate of \ce{SiO2} forms. The difference is exploited in the laboratory: \ce{SiCl4} is used to produce pure silica smoke for smoke screens.
\end{exemplebox}

\begin{exoresolu}[Hydrolysis and acid–base behaviour]
\textbf{Statement:} A student adds a few drops of water to separate test tubes containing \ce{CCl4}, \ce{SiCl4} and \ce{SnCl4}. She then tests the resulting mixtures with blue and red litmus. Predict and explain the observations.

\textbf{Solution:}
\begin{enumerate}
\item \ce{CCl4}: No reaction. Two layers remain. Litmus unchanged (neutral).
\item \ce{SiCl4}: Violent hydrolysis: \ce{SiCl4 + 2H2O -> SiO2(s) + 4HCl(aq)}. The solution becomes strongly acidic — blue litmus turns red, red litmus stays red.
\item \ce{SnCl4}: Similarly hydrolyses: \ce{SnCl4 + 2H2O -> SnO2(s) + 4HCl(aq)}. Acidic solution — same litmus result as \ce{SiCl4}.
\end{enumerate}
\textbf{Key point:} The hydrolysis of \ce{SiCl4} and \ce{SnCl4} produces \ce{HCl}, giving an acidic solution, whereas \ce{CCl4} does not hydrolyse because carbon cannot expand its octet.
\end{exoresolu}

\begin{aretenir}[GCE favourite: \ce{CO2} vs \ce{SiO2}]
Examiners frequently ask: ``Explain why \ce{CO2} is a gas at room temperature while \ce{SiO2} is a high‑melting solid, even though both are Group~IV dioxides.'' The answer must mention:
\begin{itemize}
\item \ce{CO2}: discrete linear molecules, weak van der Waals forces.
\item \ce{SiO2}: giant covalent network of \ce{[SiO4]} tetrahedra, strong \ce{Si-O} bonds throughout the lattice.
\end{itemize}
\end{aretenir}

\begin{savaistu}[Silicon in Cameroon]
Cameroon's volcanic regions (e.g. Mount Cameroon) are rich in silica (\ce{SiO2}) deposits. Local industries use quartz sand to produce glass and ceramics. The high melting point of \ce{SiO2} (≈ 1710 °C) makes it ideal for refractory bricks in metallurgical furnaces.
\end{savaistu}

\begin{exercice}[Practice questions]
\begin{enumerate}
\item Write balanced equations for the reaction of (a) \ce{Sn} with concentrated \ce{HNO3}, (b) \ce{Si} with hot concentrated \ce{NaOH}.
\item Classify each of the following oxides as acidic, basic, amphoteric or neutral: \ce{CO}, \ce{CO2}, \ce{SnO}, \ce{PbO2}, \ce{SiO2}.
\item Explain why \ce{PbCl4} decomposes on heating while \ce{SnCl4} does not.
\item Describe the trend in thermal stability of the Group~IV hydrides and relate it to bond strength.
\end{enumerate}
\end{exercice}

\begin{corrige}[Exercise 1]
\begin{enumerate}
\item (a) \ce{Sn(s) + 4HNO3(conc) -> SnO2(s) + 4NO2(g) + 2H2O(l)} \\
      (b) \ce{Si(s) + 2NaOH(aq) + H2O(l) -> Na2SiO3(aq) + 2H2(g)}
\item \ce{CO} — neutral; \ce{CO2} — acidic; \ce{SnO} — amphoteric; \ce{PbO2} — amphoteric; \ce{SiO2} — acidic.
\item \ce{PbCl4} decomposes because the +4 oxidation state of lead is destabilised by the inert pair effect (reluctance of the 6s² electrons to participate in bonding). Tin's 5s² pair is less inert, so \ce{SnCl4} remains stable.
\item Thermal stability decreases down the group: \ce{CH4} > \ce{SiH4} > \ce{GeH4} > \ce{SnH4} > \ce{PbH4}. The \ce{X-H} bond length increases and bond dissociation energy decreases, making the heavier hydrides easier to decompose.
\end{enumerate}
\end{corrige}

\coursec{Introduction to the d-Block: Electronic Configurations and Physical Trends}

Having surveyed the chemistry of Group~IV elements — their hydrides, oxides and chlorides — we now turn to the \cle{d-block} of the periodic table.  These elements, often called the \cle{transition metals}, are the workhorses of modern industry: iron in construction, copper in electrical wiring, nickel in batteries, and many serve as catalysts in chemical manufacture.  In this section we define the d-block, distinguish true transition elements from the borderline cases, master the electronic configurations of atoms and ions, and rationalise the observed trends in atomic radius, ionisation energy, melting point and related properties.

\courssub{What are d-Block Elements?}

\begin{definition}[d-Block Elements]
The \textbf{d-block} comprises the ten groups (3–12) of the periodic table in which the \cle{differentiating electron} enters a \cle{(n-1)d} subshell.  For the first transition series this means the 3d subshell is filled after the 4s subshell.  The block lies between the s-block (Groups 1–2) and the p-block (Groups 13–18).
\end{definition}

The first transition series runs from \ce{Sc} (Z=21) to \ce{Zn} (Z=30).  All are metals with typical metallic properties: high electrical and thermal conductivity, malleability, ductility and lustre.

\courssub{Transition Elements: The Strict Definition}

Not every d-block element qualifies as a \cle{transition element} in the strict IUPAC sense.

\begin{definition}[Transition Element]
A \textbf{transition element} is a d-block element that forms at least one stable ion with a \textbf{partially filled d-subshell}.
\end{definition}

\begin{attention}[Borderline Cases: Scandium and Zinc]
\ce{Sc} ([\ce{Ar}] 3d$^{1}$ 4s$^{2}$) forms only \ce{Sc^{3+}} ([\ce{Ar}] 3d$^{0}$) — an empty d-subshell.  
\ce{Zn} ([\ce{Ar}] 3d$^{10}$ 4s$^{2}$) forms only \ce{Zn^{2+}} ([\ce{Ar}] 3d$^{10}$) — a completely filled d-subshell.  
Hence \textbf{Sc and Zn are d-block elements but are NOT transition elements}.  This distinction is a favourite GCE examination point.
\end{attention}

\courssub{Electronic Configurations of the First Transition Series}

\courssub{Building Up: The Aufbau Principle and the 4s/3d Order}

According to the Aufbau principle, the 4s orbital is lower in energy than 3d for the neutral atoms of the first series, so it fills first.  The general configuration is \ce{[Ar]} 3d$^{n}$ 4s$^{2}$ (n = 1–10).  However, two elements adopt anomalous configurations because a half-filled (d$^{5}$) or fully-filled (d$^{10}$) d-subshell confers extra stability through exchange energy and symmetry.

\begin{propriete}[Anomalous Configurations of Cr and Cu]
\begin{itemize}
\item \ce{Cr} (Z=24): expected \ce{[Ar]} 3d$^{4}$ 4s$^{2}$ → actual \ce{[Ar]} \textbf{3d$^{5}$ 4s$^{1}$} (half-filled d-subshell).
\item \ce{Cu} (Z=29): expected \ce{[Ar]} 3d$^{9}$ 4s$^{2}$ → actual \ce{[Ar]} \textbf{3d$^{10}$ 4s$^{1}$} (fully-filled d-subshell).
\end{itemize}
\end{propriete}

\begin{methode}[Writing the Configuration of a Transition Metal Ion]
When a transition metal forms a cation, electrons are removed \textbf{from the highest principal quantum number first}, i.e. from the 4s orbital before the 3d orbital.
\begin{enumerate}
\item Write the ground-state configuration of the neutral atom (including anomalies for Cr and Cu).
\item Remove the required number of electrons from the 4s subshell.
\item If more electrons must be removed, take them from the 3d subshell.
\end{enumerate}
\end{methode}

\paragraph{Example:} \ce{Fe} (Z=26) is \ce{[Ar]} 3d$^{6}$ 4s$^{2}$.  
\ce{Fe^{2+}}: remove two 4s electrons → \ce{[Ar]} 3d$^{6}$.  
\ce{Fe^{3+}}: remove two 4s then one 3d → \ce{[Ar]} 3d$^{5}$.

\begin{exemplebox}[Electron-in-Boxes Diagrams]
The orbital (electron-in-boxes) notation makes the pairing and unpaired electrons explicit.  Below are the diagrams for \ce{Cr}, \ce{Cu} and \ce{Fe^{2+}}.  Each box represents one orbital; arrows represent electrons with opposite spins.
\end{exemplebox}

\begin{popfigure}
\centering
% Chromium
\begin{tikzpicture}[scale=0.7, baseline=(Cr.base)]
\node (Cr) at (0,0.8) {\textbf{Cr: [Ar] 3d$^{5}$ 4s$^{1}$}};
% 3d boxes
\foreach \i in {1,...,5}{
  \draw (\i*1.2-3, -1) rectangle ++(1,1);
  \draw (\i*1.2-3+0.5, -0.5) node {$\uparrow$};
}
\node at (3, -1.5) {3d};
% 4s box
\draw (6.5, -1) rectangle ++(1,1);
\draw (7, -0.5) node {$\uparrow$};
\node at (7, -1.5) {4s};
\end{tikzpicture}
\quad
% Copper
\begin{tikzpicture}[scale=0.7, baseline=(Cu.base)]
\node (Cu) at (0,0.8) {\textbf{Cu: [Ar] 3d$^{10}$ 4s$^{1}$}};
\foreach \i in {1,...,5}{
  \draw (\i*1.2-3, -1) rectangle ++(1,1);
  \draw (\i*1.2-3+0.5, -0.5) node {$\uparrow\downarrow$};
}
\node at (3, -1.5) {3d};
\draw (6.5, -1) rectangle ++(1,1);
\draw (7, -0.5) node {$\uparrow$};
\node at (7, -1.5) {4s};
\end{tikzpicture}
\quad
% Fe2+
\begin{tikzpicture}[scale=0.7, baseline=(Fe.base)]
\node (Fe) at (0,0.8) {\textbf{Fe$^{2+}$: [Ar] 3d$^{6}$}};
\foreach \i in {1,...,5}{
  \draw (\i*1.2-3, -1) rectangle ++(1,1);
  \ifnum\i<4
    \draw (\i*1.2-3+0.5, -0.5) node {$\uparrow\downarrow$};
  \else
    \draw (\i*1.2-3+0.5, -0.5) node {$\uparrow$};
  \fi
}
\node at (3, -1.5) {3d};
\node at (7, -1.5) {4s (empty)};
\end{tikzpicture}
\legende{Electron-in-boxes (orbital) diagrams for Cr, Cu and Fe$^{2+}$.  Each box is one orbital; arrows show electron spin.  Hund's rule: electrons occupy degenerate orbitals singly before pairing.}
\end{popfigure}

\begin{exoresolu}[Worked Example: Configurations of Co$^{2+}$ and Ni$^{3+}$]
\textbf{Statement:} Write the full electronic configuration (in spdf notation) and the electron-in-boxes diagram for the 3d subshell of \ce{Co^{2+}} and \ce{Ni^{3+}}.

\textbf{Solution:}
\begin{enumerate}
\item \ce{Co} (Z=27): neutral = \ce{[Ar]} 3d$^{7}$ 4s$^{2}$.  
      \ce{Co^{2+}}: remove two 4s electrons → \ce{[Ar]} \textbf{3d$^{7}$}.  
      3d boxes (5 orbitals, 7 electrons): three orbitals doubly occupied, two singly occupied.  
      Diagram: $\uparrow\downarrow$ $\uparrow\downarrow$ $\uparrow\downarrow$ $\uparrow$ $\uparrow$.
\item \ce{Ni} (Z=28): neutral = \ce{[Ar]} 3d$^{8}$ 4s$^{2}$.  
      \ce{Ni^{3+}}: remove two 4s then one 3d → \ce{[Ar]} \textbf{3d$^{7}$}.  
      Same d-electron count as \ce{Co^{2+}}, hence identical 3d diagram.
\end{enumerate}
\textbf{Key point:} Ions with the same d$^{n}$ configuration have similar magnetic and colour properties.
\end{exoresolu}

\begin{aretenir}[Complete Configurations of the First Transition Series (Neutral Atoms and Common Ions)]
\begin{center}
\begin{tabular}{|c|c|c|c|c|}\hline
\textbf{Element (Z)} & \textbf{Neutral Atom} & \textbf{M$^{2+}$} & \textbf{M$^{3+}$} & \textbf{Transition Element?} \\ \hline
Sc (21) & [Ar] 3d$^{1}$ 4s$^{2}$ & [Ar] 3d$^{1}$ & [Ar] 3d$^{0}$ & No (Sc$^{3+}$ is d$^{0}$) \\ \hline
Ti (22) & [Ar] 3d$^{2}$ 4s$^{2}$ & [Ar] 3d$^{2}$ & [Ar] 3d$^{1}$ & Yes \\ \hline
V (23) & [Ar] 3d$^{3}$ 4s$^{2}$ & [Ar] 3d$^{3}$ & [Ar] 3d$^{2}$ & Yes \\ \hline
Cr (24) & [Ar] 3d$^{5}$ 4s$^{1}$ & [Ar] 3d$^{4}$ & [Ar] 3d$^{3}$ & Yes \\ \hline
Mn (25) & [Ar] 3d$^{5}$ 4s$^{2}$ & [Ar] 3d$^{5}$ & [Ar] 3d$^{4}$ & Yes \\ \hline
Fe (26) & [Ar] 3d$^{6}$ 4s$^{2}$ & [Ar] 3d$^{6}$ & [Ar] 3d$^{5}$ & Yes \\ \hline
Co (27) & [Ar] 3d$^{7}$ 4s$^{2}$ & [Ar] 3d$^{7}$ & [Ar] 3d$^{6}$ & Yes \\ \hline
Ni (28) & [Ar] 3d$^{8}$ 4s$^{2}$ & [Ar] 3d$^{8}$ & [Ar] 3d$^{7}$ & Yes \\ \hline
Cu (29) & [Ar] 3d$^{10}$ 4s$^{1}$ & [Ar] 3d$^{9}$ & [Ar] 3d$^{8}$ & Yes \\ \hline
Zn (30) & [Ar] 3d$^{10}$ 4s$^{2}$ & [Ar] 3d$^{10}$ & — & No (Zn$^{2+}$ is d$^{10}$) \\ \hline
\end{tabular}
\end{center}
\end{aretenir}

\courssub{Trends in Physical Properties Across the First Transition Series}

\courssub{Atomic Radius}

Across the series the atomic radius shows a \textbf{small decrease from Sc to Cr}, then remains \textbf{nearly constant} from Mn to Cu, with a slight increase at Zn.

\begin{aretenir}[Why the Radius is Nearly Constant]
\begin{itemize}
\item Nuclear charge increases by +1 each step.
\item The added electrons enter the \textbf{inner 3d subshell}, which shields the outer 4s electrons effectively.
\item The increased attraction is largely offset by increased shielding → 4s electron cloud experiences similar effective nuclear charge.
\end{itemize}
\end{aretenir}

\begin{attention}[Exam Tip]
When asked to explain the near-constancy of atomic radii, \textbf{always mention}: “The additional 3d electrons shield the 4s electrons from the increasing nuclear charge, so the effective nuclear charge felt by the 4s electrons changes very little.”  Do \textbf{not} mention lanthanide contraction — that applies to the 4d/5d series.
\end{attention}

\courssub{First Ionisation Energy}

The first ionisation energy (\ce{I_1}) generally increases across the series but with \textbf{irregularities} due to exchange energy and subshell stability.  The overall increase reflects the rising nuclear charge.  The small dips at Cr and Cu arise because the electron removed comes from a singly occupied 4s orbital, leaving a stable half-filled (d$^{5}$) or fully-filled (d$^{10}$) d-subshell in the cation.  The sharp rise at Zn occurs because the electron is removed from a stable 4s$^{2}$ orbital into a stable 3d$^{10}$ core.

\begin{popfigure}
\centering
\begin{tikzpicture}
\begin{axis}[
  width=0.95\linewidth,
  height=7cm,
  xlabel={Element (Sc=1 … Zn=10)},
  ylabel={First Ionisation Energy / kJ\,mol$^{-1}$},
  xtick={1,2,3,4,5,6,7,8,9,10},
  xticklabels={Sc,Ti,V,Cr,Mn,Fe,Co,Ni,Cu,Zn},
  ymin=600, ymax=950,
  grid=major,
  legend pos=north west,
  mark=*,
  mark size=2,
]
\addplot[popblue, thick] coordinates {
  (1,633) (2,658) (3,650) (4,653) (5,717) (6,762) (7,760) (8,737) (9,745) (10,906)
};
\addlegendentry{IE$_1$}
\end{axis}
\end{tikzpicture}
\legende{First ionisation energy across the first transition series.  Note the general increase, small irregularities at V, Cr, Ni, Cu, and the sharp rise at Zn.}
\end{popfigure}

\courssub{Melting and Boiling Points}

Transition metals have high melting and boiling points due to strong metallic bonding involving both 4s and 3d electrons.  The melting point peaks around the middle of the series (V, Cr) where the number of unpaired d-electrons available for metallic bonding is maximal.  Mn has a lower melting point due to its complex crystal structure.  Zn, with a full 3d$^{10}$ subshell, has a much lower melting point because only the 4s electrons contribute to metallic bonding.

\begin{popfigure}
\centering
\begin{tikzpicture}
\begin{axis}[
  width=0.95\linewidth,
  height=7cm,
  xlabel={Element (Sc=1 … Zn=10)},
  ylabel={Melting Point / \ensuremath{{}^{\circ}} C},
  xtick={1,2,3,4,5,6,7,8,9,10},
  xticklabels={Sc,Ti,V,Cr,Mn,Fe,Co,Ni,Cu,Zn},
  ymin=300, ymax=2000,
  grid=major,
  legend pos=north west,
  mark=*,
  mark size=2,
]
\addplot[poporange, thick] coordinates {
  (1,1541) (2,1668) (3,1910) (4,1907) (5,1246) (6,1538) (7,1495) (8,1455) (9,1085) (10,419)
};
\addlegendentry{M.P.}
\end{axis}
\end{tikzpicture}
\legende{Melting points across the first transition series.  The peak at V/Cr reflects maximum unpaired d-electrons; Mn dips due to structure; Zn is anomalously low.}
\end{popfigure}

\courssub{Density and Hardness Compared with s-Block Metals}

Transition metals are \textbf{much denser and harder} than s-block metals (e.g. Na, Mg, Ca).  This is because:
\begin{itemize}
\item Smaller atomic radii → closer packing.
\item Greater number of valence electrons (4s + 3d) participating in metallic bonding → stronger electrostatic attraction between cations and electron sea.
\end{itemize}
For example, \ce{Fe} density = 7.87 g cm$^{-3}$ vs \ce{Mg} = 1.74 g cm$^{-3}$; \ce{Fe} is hard, \ce{Mg} is relatively soft.

\begin{propriete}[Summary of Physical Trends (Sc → Zn)]
\begin{center}
\begin{tabularx}{\linewidth}{|l|X|X|}\hline
\textbf{Property} & \textbf{Trend} & \textbf{Reason} \\ \hline
Atomic radius & Small decrease, then nearly constant & 3d electrons shield 4s from increasing nuclear charge \\ \hline
First ionisation energy & General increase with irregularities & Rising nuclear charge; dips at Cr, Cu due to stable d$^{5}$/d$^{10}$ in M$^{+}$; sharp rise at Zn \\ \hline
Melting point & High, peaks at V/Cr, dips at Mn, drops at Zn & Number of unpaired d-electrons for metallic bonding; Mn structure; Zn only 4s electrons \\ \hline
Density / Hardness & Much higher than s-block & Stronger metallic bonding (more valence electrons, smaller size) \\ \hline
\end{tabularx}
\end{center}
\end{propriete}

\courssub{Common Mistakes and Exam Tips}

\begin{attention}[Frequent Errors]
\begin{itemize}
\item Writing ion configurations as \ce{3d^4 4s^2} for \ce{Fe^{2+}} instead of \ce{3d^6}.  \textbf{Always remove 4s electrons first.}
\item Calling Sc and Zn “transition elements”.  Remember the definition: \textbf{partially filled d-subshell in at least one common ion}.
\item Explaining constant atomic radius by “lanthanide contraction” — that applies to the 4d/5d series, not the first series.
\item Forgetting that \ce{Cr} and \ce{Cu} have \ce{4s^1} configurations, which affects their ionisation energies and chemistry.
\item Drawing electron-in-boxes diagrams without applying Hund's rule (maximise unpaired electrons before pairing).
\end{itemize}
\end{attention}

\courssub{Did You Know?}

\begin{savaistu}[Transition Metals in Cameroon]
\begin{itemize}
\item \textbf{Iron ore} (magnetite, hematite) deposits in the \emph{Adamawa} and \emph{East} regions feed the local steel industry.
\item \textbf{Copper} is mined in the \emph{Nkam} area; its excellent conductivity makes it vital for the national electricity grid.
\item \textbf{Nickel} and \textbf{cobalt} lateritic ores in the \emph{South-East} are prospective for battery-grade metals.
\item Many \textbf{homogeneous catalysts} (e.g. \ce{Fe}, \ce{Co} complexes) are used in the \emph{SONARA} refinery (Limbe) for hydrodesulfurisation of petroleum fractions.
\end{itemize}
\end{savaistu}

\courssub{Consolidation Exercise}

\begin{exercice}[Practice Problems]
\begin{enumerate}
\item Write the ground-state electron configuration (spdf and electron-in-boxes) for \ce{Mn} and \ce{Mn^{2+}}.
\item Explain why the first ionisation energy of \ce{Zn} is significantly higher than that of \ce{Cu}.
\item Predict which of the following ions will be coloured in aqueous solution: \ce{Sc^{3+}}, \ce{Ti^{3+}}, \ce{V^{2+}}, \ce{Zn^{2+}}.  Justify your answer.
\item The melting point of \ce{Cr} is 1907 \ensuremath{{}^{\circ}} C while that of \ce{Zn} is 419 \ensuremath{{}^{\circ}} C.  Use metallic bonding theory to account for this difference.
\end{enumerate}
\end{exercice}

\begin{corrige}[Exercise 6.1]
\begin{enumerate}
\item \ce{Mn} (Z=25): \ce{[Ar]} 3d$^{5}$ 4s$^{2}$.  Boxes: 3d — five singly occupied ($\uparrow\ \uparrow\ \uparrow\ \uparrow\ \uparrow$); 4s — paired ($\uparrow\downarrow$).  
      \ce{Mn^{2+}}: remove 4s$^{2}$ → \ce{[Ar]} 3d$^{5}$.  Boxes: five singly occupied.
\item \ce{Zn} has configuration \ce{3d^{10} 4s^2}; the electron removed comes from a stable, fully-filled 4s orbital, and the resulting \ce{Zn^{+}} has a stable \ce{3d^{10}} core.  \ce{Cu} is \ce{3d^{10} 4s^1}; removing the single 4s electron yields stable \ce{Cu^{+}} (\ce{3d^{10}}), so less energy is required.
\item Coloured ions have partially filled d-subshells: \ce{Ti^{3+}} (3d$^{1}$), \ce{V^{2+}} (3d$^{3}$) → coloured.  \ce{Sc^{3+}} (3d$^{0}$) and \ce{Zn^{2+}} (3d$^{10}$) → colourless.
\item \ce{Cr} has six unpaired valence electrons (3d$^{5}$ 4s$^{1}$) contributing to a strong metallic bond; \ce{Zn} has only two 4s electrons (3d$^{10}$ core is inert), giving a much weaker metallic bond and lower melting point.
\end{enumerate}
\end{corrige}

This completes the foundation on d-block electronic structure and physical trends.  In the next section we will explore the hallmark chemical behaviours of transition metals: \textbf{variable oxidation states} and the fascinating world of \textbf{complex formation}.

\coursec{Variable Oxidation States and Complex Formation in Transition Metals}

In the previous section we saw how the \emph{d}-block elements of the first transition series (Sc to Zn) are built up by filling the 3d and 4s orbitals. We also observed trends in atomic radius, ionisation energy and melting points. Now we turn to the two most striking chemical consequences of that electronic structure: the \textbf{variable oxidation states} that make transition metal chemistry so rich in redox reactions, and the \textbf{formation of complex ions} through coordinate bonding with ligands. These two features are the hallmark of transition metal chemistry and underlie their importance in catalysis, biological systems and industrial processes.

\courssub{Why Transition Metals Exhibit Variable Oxidation States}

Look at the electronic configurations of the first-row transition metals and their common ions. For example, manganese:
\[
\ce{Mn}: [\ce{Ar}]\,3d^5\,4s^2 \quad \ce{Mn^{2+}}: [\ce{Ar}]\,3d^5 \quad \ce{Mn^{7+}} \text{ (in \ce{MnO4-})}: [\ce{Ar}]
\]
The 4s electrons are removed first, then the 3d electrons. Because the 3d and 4s energy levels are very close, successive ionisation energies do not jump dramatically until a stable configuration (empty, half-filled or fully-filled d subshell) is reached. This allows a metal to lose different numbers of electrons under different conditions, giving a range of oxidation states.

\begin{propriete}[Origin of Variable Oxidation States]
Transition metals show variable oxidation states because the \textbf{3d and 4s energy levels are very close}, allowing successive removal of electrons from both subshells without a prohibitively large energy gap. The maximum oxidation state usually equals the number of 4s + 3d electrons (e.g. +7 for Mn, +6 for Cr), while the minimum is typically +2 (loss of the two 4s electrons).
\end{propriete}

\textbf{Electron-in-boxes representation} (required by the syllabus) helps visualise the filling and ionisation. The 3d subshell has five degenerate orbitals; electrons occupy them singly before pairing (Hund's rule).

\begin{popfigure}
\begin{tikzpicture}[scale=0.85, every node/.style={font=\small}]
% Styles
\tikzset{orbital/.style={draw=popdark, thick, minimum width=0.7cm, minimum height=0.7cm, inner sep=0pt}}
\tikzset{labelstyle/.style={anchor=north, yshift=-0.1cm}}
% Titles
\node[anchor=south] at (0,3.5) {\textbf{3d}};
\node[anchor=south] at (5.5,3.5) {\textbf{4s}};
% Elements: Cr, Mn, Fe, Co, Ni, Cu, Zn
% We'll draw for Cr, Mn, Fe, Cu as examples
% Cr: [Ar] 3d5 4s1
\begin{scope}[xshift=-8cm]
\node[anchor=east] at (-0.5,3) {\ce{Cr}:};
\foreach \i in {1,...,5} {
  \node[orbital] (d\i) at (\i*0.8,3) {};
  \draw[thick, ->] (d\i.center) -- ++(0,0.3); % up arrow
}
\node[orbital] (s1) at (5.5,3) {};
\draw[thick, ->] (s1.center) -- ++(0,0.3);
\node[labelstyle] at (2.4,2.2) {3d\textsuperscript{5}};
\node[labelstyle] at (5.5,2.2) {4s\textsuperscript{1}};
\end{scope}
% Mn: [Ar] 3d5 4s2
\begin{scope}[xshift=-8cm, yshift=-2.5cm]
\node[anchor=east] at (-0.5,0) {\ce{Mn}:};
\foreach \i in {1,...,5} {
  \node[orbital] (d\i) at (\i*0.8,0) {};
  \draw[thick, ->] (d\i.center) -- ++(0,0.3);
}
\node[orbital] (s1) at (5.5,0) {};
\draw[thick, ->] (s1.center) -- ++(0,0.3);
\draw[thick, <-] (s1.center) -- ++(0,-0.3);
\node[labelstyle] at (2.4,-0.8) {3d\textsuperscript{5}};
\node[labelstyle] at (5.5,-0.8) {4s\textsuperscript{2}};
\end{scope}
% Fe: [Ar] 3d6 4s2
\begin{scope}[xshift=-8cm, yshift=-5cm]
\node[anchor=east] at (-0.5,0) {\ce{Fe}:};
\foreach \i in {1,...,5} {
  \node[orbital] (d\i) at (\i*0.8,0) {};
  if \i<6 \draw[thick, ->] (d\i.center) -- ++(0,0.3);
}
% Actually 6 electrons: 5 up, 1 down in first orbital
\node[orbital] (d1) at (0.8,0) {};
\draw[thick, ->] (d1.center) -- ++(0,0.3);
\draw[thick, <-] (d1.center) -- ++(0,-0.3);
\foreach \i in {2,...,5} {
  \node[orbital] (d\i) at (\i*0.8,0) {};
  \draw[thick, ->] (d\i.center) -- ++(0,0.3);
}
\node[orbital] (s1) at (5.5,0) {};
\draw[thick, ->] (s1.center) -- ++(0,0.3);
\draw[thick, <-] (s1.center) -- ++(0,-0.3);
\node[labelstyle] at (2.4,-0.8) {3d\textsuperscript{6}};
\node[labelstyle] at (5.5,-0.8) {4s\textsuperscript{2}};
\end{scope}
% Cu: [Ar] 3d10 4s1
\begin{scope}[xshift=-8cm, yshift=-7.5cm]
\node[anchor=east] at (-0.5,0) {\ce{Cu}:};
\foreach \i in {1,...,5} {
  \node[orbital] (d\i) at (\i*0.8,0) {};
  \draw[thick, ->] (d\i.center) -- ++(0,0.3);
  \draw[thick, <-] (d\i.center) -- ++(0,-0.3);
}
\node[orbital] (s1) at (5.5,0) {};
\draw[thick, ->] (s1.center) -- ++(0,0.3);
\node[labelstyle] at (2.4,-0.8) {3d\textsuperscript{10}};
\node[labelstyle] at (5.5,-0.8) {4s\textsuperscript{1}};
\end{scope}
% Ions: Fe2+, Fe3+, Cu+, Cu2+
\begin{scope}[xshift=4cm]
\node[anchor=east] at (-0.5,3) {\ce{Fe^{2+}}:};
\foreach \i in {1,...,5} {
  \node[orbital] (d\i) at (\i*0.8,3) {};
  \draw[thick, ->] (d\i.center) -- ++(0,0.3);
}
\node[orbital] (d1) at (0.8,3) {};
\draw[thick, ->] (d1.center) -- ++(0,0.3);
\draw[thick, <-] (d1.center) -- ++(0,-0.3); % 6th electron paired
\foreach \i in {2,...,5} {
  \node[orbital] (d\i) at (\i*0.8,3) {};
  \draw[thick, ->] (d\i.center) -- ++(0,0.3);
}
\node[orbital] (s1) at (5.5,3) {}; % empty 4s
\node[labelstyle] at (2.4,2.2) {3d\textsuperscript{6}};
\node[labelstyle] at (5.5,2.2) {4s\textsuperscript{0}};
\end{scope}
\begin{scope}[xshift=4cm, yshift=-2.5cm]
\node[anchor=east] at (-0.5,0) {\ce{Fe^{3+}}:};
\foreach \i in {1,...,5} {
  \node[orbital] (d\i) at (\i*0.8,0) {};
  \draw[thick, ->] (d\i.center) -- ++(0,0.3);
}
\node[orbital] (s1) at (5.5,0) {};
\node[labelstyle] at (2.4,-0.8) {3d\textsuperscript{5}};
\node[labelstyle] at (5.5,-0.8) {4s\textsuperscript{0}};
\end{scope}
\begin{scope}[xshift=4cm, yshift=-5cm]
\node[anchor=east] at (-0.5,0) {\ce{Cu^{+}}:};
\foreach \i in {1,...,5} {
  \node[orbital] (d\i) at (\i*0.8,0) {};
  \draw[thick, ->] (d\i.center) -- ++(0,0.3);
  \draw[thick, <-] (d\i.center) -- ++(0,-0.3);
}
\node[orbital] (s1) at (5.5,0) {};
\node[labelstyle] at (2.4,-0.8) {3d\textsuperscript{10}};
\node[labelstyle] at (5.5,-0.8) {4s\textsuperscript{0}};
\end{scope}
\begin{scope}[xshift=4cm, yshift=-7.5cm]
\node[anchor=east] at (-0.5,0) {\ce{Cu^{2+}}:};
\foreach \i in {1,...,4} {
  \node[orbital] (d\i) at (\i*0.8,0) {};
  \draw[thick, ->] (d\i.center) -- ++(0,0.3);
  \draw[thick, <-] (d\i.center) -- ++(0,-0.3);
}
\node[orbital] (d5) at (4.0,0) {};
\draw[thick, ->] (d5.center) -- ++(0,0.3); % one unpaired
\node[orbital] (s1) at (5.5,0) {};
\node[labelstyle] at (2.4,-0.8) {3d\textsuperscript{9}};
\node[labelstyle] at (5.5,-0.8) {4s\textsuperscript{0}};
\end{scope}
% Arrows for ionisation
\draw[->, thick, poporange] (-3.5,3) -- (-2,3) node[midway, above] {lose 4s\textsuperscript{1}};
\draw[->, thick, poporange] (-3.5,0.5) -- (-2,0.5) node[midway, above] {lose 4s\textsuperscript{2}};
\draw[->, thick, poporange] (-3.5,-2) -- (-2,-2) node[midway, above] {lose 4s\textsuperscript{2}};
\draw[->, thick, poporange] (-3.5,-4.5) -- (-2,-4.5) node[midway, above] {lose 4s\textsuperscript{1}};
\end{tikzpicture}
\legende{Electron-in-boxes diagrams for selected first-row transition metals and their common ions. Arrows indicate electron spin. Note the anomalous configurations of Cr (3d\textsuperscript{5}4s\textsuperscript{1}) and Cu (3d\textsuperscript{10}4s\textsuperscript{1}) due to extra stability of half-filled and fully-filled d subshells.}
\end{popfigure}

\begin{center}
\begin{tabularx}{\linewidth}{|c|X|X|}\hline
\rowcolor{popblueL}
\textbf{Element} & \textbf{Common Oxidation States} & \textbf{Most Stable / Typical Compounds} \\ \hline
Sc & +3 & \ce{Sc2O3}, \ce{ScCl3} \\
Ti & +2, +3, \textbf{+4} & \ce{TiO2}, \ce{TiCl4} \\
V & +2, +3, +4, \textbf{+5} & \ce{VO^{2+}}, \ce{VO3-} \\
Cr & +2, \textbf{+3}, +6 & \ce{Cr^{3+}}, \ce{Cr2O7^{2-}} \\
Mn & +2, +3, +4, +6, \textbf{+7} & \ce{Mn^{2+}}, \ce{MnO2}, \ce{MnO4-} \\
Fe & \textbf{+2}, \textbf{+3} & \ce{Fe^{2+}}, \ce{Fe^{3+}} \\
Co & +2, +3 & \ce{Co^{2+}}, \ce{Co^{3+}} (in complexes) \\
Ni & \textbf{+2} & \ce{Ni^{2+}} \\
Cu & \textbf{+1}, \textbf{+2} & \ce{Cu+}, \ce{Cu^{2+}} \\
Zn & +2 & \ce{Zn^{2+}} \\ \hline
\end{tabularx}
\end{center}
\textbf{Note:} Bold oxidation states are the most stable in aqueous solution. Higher oxidation states (+5, +6, +7) are strongly oxidising and are stabilised in oxides or oxoanions (e.g. \ce{MnO4-}, \ce{Cr2O7^{2-}}, \ce{VO3-}) where the metal is bonded to highly electronegative oxygen. Lower oxidation states (+2, +3) predominate in simple ionic compounds (halides, sulfates) and aqua ions.

\courssub{Redox Interconversions of Key Oxidation States}

The variable oxidation states allow a rich redox chemistry. Three systems are fundamental for GCE Advanced Level: iron, chromium and manganese. You must know the half-equations, the colour changes and the conditions (acidic medium).

\begin{definition}[Half-Equation]
A half-equation shows either the oxidation or the reduction process separately, with electrons explicitly written. In acidic medium, balance O with \ce{H2O}, H with \ce{H+}, and charge with \ce{e-}.
\end{definition}

\begin{methode}[Balancing a Half-Equation in Acidic Medium]
\begin{enumerate}
\item Write the species containing the element being reduced/oxidised on the left, the product on the right.
\item Balance the element (other than O and H).
\item Balance O by adding \ce{H2O} to the side deficient in O.
\item Balance H by adding \ce{H+} to the side deficient in H.
\item Balance charge by adding electrons (\ce{e-}) to the more positive side.
\end{enumerate}
\end{methode}

\noindent\textbf{1. Iron: \ce{Fe^{3+}/Fe^{2+}}}
\[
\ce{Fe^{3+}(aq) + e- <=> Fe^{2+}(aq)} \qquad E^\circ = +0.77\ \mathrm{V}
\]
\emph{Colour change:} Pale yellow (\ce{Fe^{3+}}) $\rightleftharpoons$ Pale green (\ce{Fe^{2+}}). This couple is used in redox titrations (e.g. with \ce{MnO4-} or \ce{Cr2O7^{2-}}).

\noindent\textbf{2. Chromium: \ce{Cr2O7^{2-}/Cr^{3+}} (acidic medium)}
\[
\ce{Cr2O7^{2-}(aq) + 14H+(aq) + 6e- <=> 2Cr^{3+}(aq) + 7H2O(l)} \qquad E^\circ = +1.33\ \mathrm{V}
\]
\emph{Colour change:} Orange (\ce{Cr2O7^{2-}}) $\rightarrow$ Green (\ce{Cr^{3+}}). \ce{K2Cr2O7} is a primary standard in redox titrations.

\noindent\textbf{3. Manganese: \ce{MnO4-/Mn^{2+}} (acidic medium)}
\[
\ce{MnO4-(aq) + 8H+(aq) + 5e- <=> Mn^{2+}(aq) + 4H2O(l)} \qquad E^\circ = +1.51\ \mathrm{V}
\]
\emph{Colour change:} Deep purple (\ce{MnO4-}) $\rightarrow$ Colourless (\ce{Mn^{2+}}). The reaction is self-indicating; the first permanent pink colour signals the endpoint.

\begin{attention}[Common Mistakes in Half-Equations]
\begin{itemize}
\item Forgetting to balance oxygen with \ce{H2O} and hydrogen with \ce{H+} in acidic medium.
\item Writing the wrong number of electrons: count the change in oxidation state per formula unit (e.g. Cr in \ce{Cr2O7^{2-}} is +6 each, total +12; two \ce{Cr^{3+}} give +6; change = 6 electrons).
\item Using \ce{OH-} instead of \ce{H+} when the problem states "acidic medium".
\item Confusing the colour of \ce{Cr^{3+}} (green) with \ce{Cr^{2+}} (blue) or \ce{Fe^{2+}} (pale green) with \ce{Fe^{3+}} (yellow/brown).
\end{itemize}
\end{attention}

\begin{exoresolu}[Balancing a Redox Reaction: \ce{MnO4-} with \ce{Fe^{2+}}]
\textbf{Statement:} Write the balanced ionic equation for the reaction between acidified potassium manganate(VII) and iron(II) sulfate.

\textbf{Detailed solution:}
\begin{enumerate}
\item Write the two half-equations (reduction of \ce{MnO4-}, oxidation of \ce{Fe^{2+}}):
\[
\text{Reduction: } \ce{MnO4- + 8H+ + 5e- -> Mn^{2+} + 4H2O}
\]
\[
\text{Oxidation: } \ce{Fe^{2+} -> Fe^{3+} + e-}
\]
\item Equalise electrons: multiply the oxidation half-equation by 5.
\[
\ce{5Fe^{2+} -> 5Fe^{3+} + 5e-}
\]
\item Add the two half-equations, cancelling electrons:
\[
\ce{MnO4- + 8H+ + 5Fe^{2+} -> Mn^{2+} + 4H2O + 5Fe^{3+}}
\]
\item Check: atoms and charge balanced. Left: charge = $-1 + 8 + 10 = +17$; Right: $+2 + 15 = +17$. \checkmark
\end{enumerate}
\end{exoresolu}

\courssub{Complex Ions: Ligands, Coordination and Classification}

When a transition metal ion is in solution or reacts with certain molecules/ions, it does not remain as a bare ion. It attracts species that have lone pairs of electrons, forming \textbf{coordinate (dative) bonds}. The resulting assembly is a \textbf{complex ion}.

\begin{definition}[Ligand]
A \textbf{ligand} is a molecule or ion that donates a lone pair of electrons to a central metal ion to form a coordinate bond. The ligand acts as a Lewis base; the metal ion acts as a Lewis acid.
\end{definition}

\begin{definition}[Coordination Number]
The \textbf{coordination number} is the number of coordinate bonds formed between the ligands and the central metal ion. It is \emph{not} necessarily the number of ligands (a bidentate ligand forms two bonds).
\end{definition}

\noindent Ligands are classified by the number of donor atoms (\textbf{denticity}):

\begin{popfigure}
\begin{center}
\begin{tabularx}{\linewidth}{|c|X|X|c|}\hline
\rowcolor{popblueL}
\textbf{Type} & \textbf{Examples} & \textbf{Formula / Charge} & \textbf{Denticity} \\ \hline
Monodentate (neutral) & Water, Ammonia, Carbon monoxide & \ce{H2O}, \ce{NH3}, \ce{CO} (0) & 1 \\ \hline
Monodentate (anionic) & Chloride, Cyanide, Hydroxide, Thiocyanate & \ce{Cl-}, \ce{CN-}, \ce{OH-}, \ce{SCN-} (-1) & 1 \\ \hline
Bidentate & Ethylenediamine (en), Oxalate & \ce{NH2CH2CH2NH2} (0), \ce{C2O4^{2-}} (-2) & 2 \\ \hline
Polydentate & EDTA$^{4-}$ & \ensuremath{\mathrm{[EDTA]^{4-}}} & 6 \\ \hline
Ambidentate & Thiocyanate, Nitrite & \ce{SCN-} (S or N), \ce{NO2-} (N or O) & 1 (two possible donor atoms) \\ \hline
\end{tabularx}
\end{center}
\legende{Common ligands, their charges and denticity.}
\end{popfigure}

\begin{aretenir}[Chelation and the Chelate Effect]
A ligand that forms more than one bond to the same metal ion is a \textbf{chelating ligand}. The resulting ring structure (chelate ring, usually 5- or 6-membered) gives extra thermodynamic stability compared to similar monodentate ligands. This is the \textbf{chelate effect}. Ethylenediamine (en) and EDTA are classic chelating agents. EDTA\textsuperscript{4-} wraps around the metal ion in an octahedral arrangement, forming five chelate rings.
\end{aretenir}

\courssub{Coordination Number and Geometry}

The coordination number (CN) determines the shape of the complex. For the first transition series, CN = 6 (octahedral), 4 (tetrahedral or square planar) and 2 (linear) are most common.

\begin{center}
\begin{tabularx}{\linewidth}{|c|X|X|}\hline
\rowcolor{popblueL}
\textbf{CN} & \textbf{Geometry} & \textbf{Examples} \\ \hline
2 & Linear & \ce{[Ag(NH3)2]+}, \ce{[CuCl2]-} (in conc. HCl) \\
4 & Tetrahedral & \ensuremath{\mathrm{[CoCl_{4}]^{2-}}}, \ensuremath{\mathrm{[CuCl_{4}]^{2-}}}, \ensuremath{\mathrm{[Zn(NH_{3})_{4}]^{2+}}} \\
4 & Square planar & \ensuremath{\mathrm{[Ni(CN)_{4}]^{2-}}}, \ce{[Pt(NH3)2Cl2]} (Pt is 2nd row) \\
6 & Octahedral & \ensuremath{\mathrm{[Fe(H_{2}O)_{6}]^{2+}}}, \ensuremath{\mathrm{[Co(NH_{3})_{6}]^{3+}}}, \ensuremath{\mathrm{[Cu(NH_{3})_{4}(H_{2}O)_{2}]^{2+}}} \\ \hline
\end{tabularx}
\end{center}

\begin{popfigure}
\begin{tikzpicture}[scale=1.3, every node/.style={font=\small}]
% Octahedral complex [Cu(NH3)4(H2O)2]2+
% Central Cu
\node[circle, draw=popdark, thick, fill=poporange!30, minimum size=8mm] (Cu) at (0,0) {\ce{Cu^{2+}}};
% Axial H2O (top and bottom)
\node[draw=popblue, thick, rounded corners=2pt, fill=popblue!10, minimum width=1.6cm, minimum height=0.7cm] (H2Oax1) at (0,2.8) {\ce{H2O}};
\node[draw=popblue, thick, rounded corners=2pt, fill=popblue!10, minimum width=1.6cm, minimum height=0.7cm] (H2Oax2) at (0,-2.8) {\ce{H2O}};
% Equatorial NH3 (four in a square)
\node[draw=popgreen, thick, rounded corners=2pt, fill=popgreen!10, minimum width=1.6cm, minimum height=0.7cm] (NH3eq1) at (2.5,0) {\ce{NH3}};
\node[draw=popgreen, thick, rounded corners=2pt, fill=popgreen!10, minimum width=1.6cm, minimum height=0.7cm] (NH3eq2) at (-2.5,0) {\ce{NH3}};
\node[draw=popgreen, thick, rounded corners=2pt, fill=popgreen!10, minimum width=1.6cm, minimum height=0.7cm] (NH3eq3) at (0,1.4) {\ce{NH3}};
\node[draw=popgreen, thick, rounded corners=2pt, fill=popgreen!10, minimum width=1.6cm, minimum height=0.7cm] (NH3eq4) at (0,-1.4) {\ce{NH3}};
% Dative bonds (arrows from ligand to metal)
\draw[->, thick, popdark] (H2Oax1.south) -- (Cu.north) node[midway, right] {$\sigma$};
\draw[->, thick, popdark] (H2Oax2.north) -- (Cu.south) node[midway, left] {$\sigma$};
\draw[->, thick, popdark] (NH3eq1.west) -- (Cu.east) node[midway, above] {$\sigma$};
\draw[->, thick, popdark] (NH3eq2.east) -- (Cu.west) node[midway, above] {$\sigma$};
\draw[->, thick, popdark] (NH3eq3.south) -- (Cu.north) node[midway, right, yshift=0.2cm] {$\sigma$};
\draw[->, thick, popdark] (NH3eq4.north) -- (Cu.south) node[midway, left, yshift=-0.2cm] {$\sigma$};
% Legend
\node[anchor=north west, text width=5.5cm] at (-4.5, -3.2) {
\textbf{Octahedral \ensuremath{\mathrm{[Cu(NH_{3})_{4}(H_{2}O)_{2}]^{2+}}}}\\
\ce{NH3} (green) and \ce{H2O} (blue) donate lone pairs (arrows) to \ce{Cu^{2+}}.\\
Coordination number = 6.\\
\emph{Jahn-Teller distortion:} axial bonds longer than equatorial.
};
\end{tikzpicture}
\legende{Structure of the tetraamminediaquacopper(II) ion showing dative bonds from ligands to the central metal ion.}
\end{popfigure}

\courssub{Classification of Complexes by Overall Charge}

The overall charge of a complex ion is the sum of the oxidation state of the metal and the charges of the ligands.

\begin{propriete}[Overall Charge of a Complex]
For a complex \ensuremath{\mathrm{[M_x L_y]^{z}}}, the overall charge $z$ is given by:
\[
z = \text{Oxidation state of M} + \sum (\text{charges of ligands})
\]
\end{propriete}

\begin{center}
\begin{tabularx}{\linewidth}{|c|X|X|}\hline
\rowcolor{popblueL}
\textbf{Type} & \textbf{Overall Charge} & \textbf{Examples} \\ \hline
Cationic & Positive & \ensuremath{\mathrm{[Cu(NH_{3})_{4}(H_{2}O)_{2}]^{2+}}}, \ensuremath{\mathrm{[Co(NH_{3})_{6}]^{3+}}}, \ensuremath{\mathrm{[Fe(H_{2}O)_{6}]^{2+}}} \\
Anionic & Negative & \ensuremath{\mathrm{[Fe(CN)_{6}]^{4-}}}, \ensuremath{\mathrm{[CuCl_{4}]^{2-}}}, \ensuremath{\mathrm{[MnO_{4}]^{-}}} \\
Neutral & Zero & \ce{[Ni(CO)4]}, \ce{[Pt(NH3)2Cl2]}, \ce{Fe(acac)3} \\ \hline
\end{tabularx}
\end{center}

\begin{methode}[Finding the Oxidation State of the Metal in a Complex]
\begin{enumerate}
\item Identify the overall charge of the complex ion (from the formula of the salt, e.g. \ce{K4[Fe(CN)6]} $\rightarrow$ complex is \ensuremath{\mathrm{[Fe(CN)_{6}]^{4-}}}).
\item List the ligands and their charges (e.g. \ce{CN-} is $-1$ each).
\item Let $x$ be the oxidation state of the metal. Write: $x + \sum(\text{ligand charges}) = \text{overall charge}$.
\item Solve for $x$.
\end{enumerate}
\end{methode}

\begin{exoresolu}[Oxidation State in a Complex]
\textbf{Statement:} Determine the oxidation state of chromium in \ce{K3[Cr(C2O4)3]} and cobalt in \ce{[Co(NH3)5Cl]Cl2}.

\textbf{Detailed solution:}
\begin{enumerate}
\item \textbf{\ce{K3[Cr(C2O4)3]}}: The salt has 3 \ce{K+} ions, so the complex anion is \ensuremath{\mathrm{[Cr(C_{2}O_{4})_{3}]^{3-}}}. Ligand: oxalate \ce{C2O4^{2-}} (charge $-2$), three of them $\rightarrow$ total ligand charge $= -6$. Let $x$ = oxidation state of Cr.
\[
x + (-6) = -3 \quad \Rightarrow \quad x = +3.
\]
\item \textbf{\ce{[Co(NH3)5Cl]Cl2}}: Two chloride ions outside the bracket are counter-ions, so the complex cation is \ensuremath{\mathrm{[Co(NH_{3})_{5}Cl]^{2+}}}. Ligands: 5 \ce{NH3} (neutral, charge 0) + 1 \ce{Cl-} (charge $-1$). Total ligand charge $= -1$. Let $x$ = oxidation state of Co.
\[
x + (-1) = +2 \quad \Rightarrow \quad x = +3.
\]
\end{enumerate}
\end{exoresolu}

\begin{attention}[Charge of Complex vs. Oxidation State of Metal]
Do not confuse the \textbf{charge of the complex ion} with the \textbf{oxidation state of the metal}. They are equal \emph{only when all ligands are neutral}. If any ligand carries a charge, you must do the calculation.
\end{attention}

\courssub{Ligand Exchange Reactions: The Chemistry of \ce{Cu^{2+}}}

Aqua ions like \ensuremath{\mathrm{[Cu(H_{2}O)_{6}]^{2+}}} (blue) undergo ligand exchange when a stronger ligand (or a higher concentration of a ligand) is added. Two classic examples are tested in GCE practical and theory papers.

\begin{experience}[Ligand Exchange with \ce{Cu^{2+}}]
\textbf{Reaction 1: Addition of aqueous ammonia (limited then excess).}
\begin{enumerate}
\item \textbf{Few drops \ce{NH3(aq)}:} \ce{Cu^{2+}} acts as an acid; \ce{NH3} acts as a base.
\[
\ensuremath{\mathrm{[Cu(H_{2}O)_{6}]^{2+}(aq) + 2NH_{3}(aq) \rightarrow  Cu(OH)_{2}(s) + 2NH_{4}+(aq) + 4H_{2}O(l)}}
\]
\textit{Observation:} Pale blue precipitate of \ce{Cu(OH)2}.
\item \textbf{Excess \ce{NH3(aq)}:} \ce{NH3} acts as a ligand, dissolving the precipitate via ligand exchange.
\[
\ensuremath{\mathrm{Cu(OH)_{2}(s) + 4NH_{3}(aq) + 2H_{2}O(l) \rightarrow  [Cu(NH_{3})_{4}(H_{2}O)_{2}]^{2+}(aq) + 2OH^{-}(aq)}}
\]
\textit{Observation:} Precipitate dissolves giving a \textbf{deep blue} solution (tetraamminediaquacopper(II) ion).
\end{enumerate}

\textbf{Reaction 2: Addition of concentrated hydrochloric acid.}
\[
\ensuremath{\mathrm{[Cu(H_{2}O)_{6}]^{2+}(aq) + 4Cl^{-}(aq) \rightleftharpoons  [CuCl_{4}]^{2-}(aq) + 6H_{2}O(l)}}
\]
\textit{Observation:} Blue solution turns \textbf{yellow-green} (tetrahedral \ensuremath{\mathrm{[CuCl_{4}]^{2-}}}). The equilibrium shifts right with high \ce{Cl-} concentration (conc. HCl) and left on dilution with water.
\end{experience}

\begin{aretenir}[Key Points for Ligand Exchange]
\begin{itemize}
\item Ligand exchange is a \textbf{substitution reaction}: one ligand replaces another in the coordination sphere.
\item The driving force can be: (a) stronger ligand field (e.g. \ce{NH3} > \ce{H2O}), (b) chelate effect (e.g. \ce{en} > \ce{NH3}), (c) concentration (Le Chatelier's principle, e.g. conc. \ce{HCl}).
\item Colour change is a direct indicator of the change in d-orbital splitting ($\Delta$) caused by the new ligand set.
\item \ce{Cu^{2+}} (d9) in octahedral field shows Jahn-Teller distortion: two axial bonds longer than four equatorial. In \ensuremath{\mathrm{[Cu(NH_{3})_{4}(H_{2}O)_{2}]^{2+}}} the four \ce{NH3} occupy the equatorial plane.
\end{itemize}
\end{aretenir}

\courssub{Summary and Exam Focus}

\begin{aretenir}[What You Must Master for the GCE]
\begin{enumerate}
\item \textbf{Variable oxidation states:} Explain using 3d/4s energy proximity. List common states for Cr, Mn, Fe, Cu. Know that high states are oxidising and stabilised by O; low states are stable in simple salts.
\item \textbf{Redox half-equations:} Write and balance \ce{MnO4-/Mn^{2+}}, \ce{Cr2O7^{2-}/Cr^{3+}}, \ce{Fe^{3+}/Fe^{2+}} in acidic medium. Know the colour changes and $E^\circ$ values.
\item \textbf{Complex terminology:} Ligand, coordinate bond, coordination number, chelate, denticity (mono-, bi-, poly-). Classify ligands with charge.
\item \textbf{Shapes:} Linear (CN=2), tetrahedral/square planar (CN=4), octahedral (CN=6). Draw/describe \ensuremath{\mathrm{[Cu(NH_{3})_{4}(H_{2}O)_{2}]^{2+}}}.
\item \textbf{Oxidation state calculation:} Use overall charge = metal oxidation state + sum of ligand charges.
\item \textbf{Ligand exchange:} Write equations for \ensuremath{\mathrm{[Cu(H_{2}O)_{6}]^{2+}}} with \ce{NH3} (two stages) and conc. \ce{HCl}. Explain observations.
\item \textbf{Electron-in-boxes:} Draw configurations for atoms and ions of the first transition series, noting anomalies for Cr and Cu.
\end{enumerate}
\end{aretenir}

\begin{savaistu}[Transition Metals in Cameroon]
The lateritic soils of the Adamawa and East regions of Cameroon are rich in iron and manganese oxides (\ce{Fe2O3}, \ce{MnO2}) — the high oxidation states stabilised by oxygen. These ores are the basis for potential ferroalloy industries. In our rivers, the \ce{Fe^{2+}/Fe^{3+}} redox couple plays a role in natural water chemistry, while the blue \ensuremath{\mathrm{[Cu(NH_{3})_{4}]^{2+}}} complex is a classic qualitative test for copper(II) in the school laboratory.
\end{savaistu}

\begin{exercice}[Practice Problems]
\begin{itemize}
\item Write the half-equation for the reduction of \ce{Cr2O7^{2-}} to \ce{Cr^{3+}} in acidic medium. What is the colour change?
\item Determine the oxidation state of the metal in: (a) \ensuremath{\mathrm{[Fe(CN)_{6}]^{3-}}}, (b) \ensuremath{\mathrm{[CoCl_{4}]^{2-}}}, (c) \ce{[Ni(NH3)6]Cl2}.
\item A solution of \ce{CuSO4} is treated with (i) a few drops of \ce{NaOH}, (ii) excess \ce{NH3(aq)}, (iii) conc. \ce{HCl}. Describe the observations and write equations for each step.
\item Explain why \ce{Mn} shows oxidation states from +2 to +7 while \ce{Zn} shows only +2.
\item The complex \ce{[Cr(en)3]Cl3} is a green crystalline solid. (en = ethylenediamine). State the coordination number, geometry, oxidation state of Cr, and overall charge of the complex ion.
\item Draw the electron-in-boxes diagram for \ce{Cr} atom and \ce{Cr^{3+}} ion. Explain the stability of \ce{Cr^{3+}}.
\end{itemize}
\end{exercice}

\begin{corrige}[Exercise 7]
\begin{itemize}
\item \ce{Cr2O7^{2-} + 14H+ + 6e- -> 2Cr^{3+} + 7H2O}. Colour: orange $\rightarrow$ green.
\item (a) \ce{CN-} is $-1$, six of them $= -6$. Complex charge $-3$. $x - 6 = -3 \Rightarrow x = +3$. (b) \ce{Cl-} is $-1$, four of them $= -4$. Complex charge $-2$. $x - 4 = -2 \Rightarrow x = +2$. (c) Two \ce{Cl-} counter-ions $\rightarrow$ complex is \ensuremath{\mathrm{[Ni(NH_{3})_{6}]^{2+}}}. \ce{NH3} neutral. $x = +2$.
\item (i) \ce{Cu^{2+} + 2OH- -> Cu(OH)2(s)} (pale blue ppt). (ii) \ensuremath{\mathrm{Cu(OH)_{2} + 4NH_{3} + 2H_{2}O \rightarrow  [Cu(NH_{3})_{4}(H_{2}O)_{2}]^{2+} + 2OH^{-}}} (deep blue solution). (iii) \ensuremath{\mathrm{[Cu(H_{2}O)_{6}]^{2+} + 4Cl^{-} \rightleftharpoons  [CuCl_{4}]^{2-} + 6H_{2}O}} (yellow-green solution).
\item Mn has configuration \ce{[Ar] 3d^5 4s^2}; all seven electrons can be lost because 3d and 4s are close in energy. Zn has \ce{[Ar] 3d^{10} 4s^2}; after losing 4s electrons, the 3d subshell is full and stable, so further ionisation requires much higher energy.
\item en is bidentate $\rightarrow$ 3 en ligands give 6 bonds $\rightarrow$ CN = 6, octahedral. Complex ion is \ensuremath{\mathrm{[Cr(en)_{3}]^{3+}}} (3 \ce{Cl-} outside). en is neutral. Oxidation state of Cr = +3. Overall charge of complex ion = +3.
\item \ce{Cr}: 3d\textsuperscript{5} 4s\textsuperscript{1} (five unpaired electrons in 3d, one in 4s). \ce{Cr^{3+}}: loses 4s\textsuperscript{1} and two 3d electrons $\rightarrow$ 3d\textsuperscript{3} (three unpaired electrons). The \ce{Cr^{3+}} ion has a half-filled t\textsubscript{2g} set (t\textsubscript{2g}\textsuperscript{3} e\textsubscript{g}\textsuperscript{0}) in an octahedral field, giving extra ligand field stabilisation energy (LFSE), making it particularly stable.
\end{itemize}
\end{corrige}

\coursec{Naming, Shapes and Isomerism of Complexes}

In the previous section we discovered that transition metals form \cle{complex ions} by accepting electron pairs from ligands. We also saw that a single metal ion can exist in several oxidation states and coordinate with different numbers of ligands. This richness gives rise to an enormous variety of compounds. To communicate about them unambiguously, chemists use a systematic nomenclature. Moreover, the spatial arrangement of ligands around the metal centre — the \cle{geometry} of the complex — determines not only its name but also its physical properties, chemical reactivity and even its biological activity. In this section we master the IUPAC naming rules, explore the common shapes of complexes, and study the fascinating phenomenon of \cle{isomerism} where compounds with the same formula behave differently because their atoms are arranged differently in space.

\courssub{IUPAC Nomenclature of Coordination Compounds}

Naming a coordination compound follows a precise sequence. The guiding principle is: \textbf{name the ligands first (in alphabetical order), then the metal centre with its oxidation state}. If the complex is an anion, the metal name takes an \cle{-ate} ending.

\begin{methode}[Naming a Coordination Compound — Step by Step]
\begin{enumerate}
\item \textbf{Identify the complex ion and the counter-ions.} Write the formula with the coordination sphere in square brackets, e.g. \ce{[Cu(NH3)4]SO4}.
\item \textbf{Name the ligands in alphabetical order} (ignoring multiplicative prefixes \textit{di-}, \textit{tri-}, \textit{tetra-}, etc.). Use the standard ligand names:
      \begin{itemize}
      \item Neutral ligands: \ce{H2O} = aqua, \ce{NH3} = ammine, \ce{CO} = carbonyl, \ce{NO} = nitrosyl.
      \item Anionic ligands: \ce{Cl-} = chloro, \ce{Br-} = bromo, \ce{I-} = iodo, \ce{F-} = fluoro, \ce{CN-} = cyano, \ce{OH-} = hydroxo, \ce{SO4^2-} = sulfato, \ce{NO2-} = nitro (N-bonded) or nitrito (O-bonded), \ce{SCN-} = thiocyanato (S-bonded) or isothiocyanato (N-bonded).
      \item For organic ligands use the common name: \ce{en} = ethylenediamine, \ce{ox} = oxalate, \ce{bipy} = 2,2'-bipyridine, \ce{phen} = 1,10-phenanthroline.
      \end{itemize}
\item \textbf{Indicate the number of each ligand} with prefixes: \textit{di-} (2), \textit{tri-} (3), \textit{tetra-} (4), \textit{penta-} (5), \textit{hexa-} (6). If the ligand name itself contains a prefix (e.g. ethylenediamine), use \textit{bis-}, \textit{tris-}, \textit{tetrakis-} and enclose the ligand name in parentheses.
\item \textbf{Name the metal.} For a cationic or neutral complex, use the element name (e.g. iron, copper). For an anionic complex, use the Latin root + \textit{-ate}: \cle{ferrate} (Fe), \cle{cuprate} (Cu), \cle{chromate} (Cr), \cle{argentate} (Ag), \cle{stannate} (Sn), \cle{plumbate} (Pb).
\item \textbf{Give the oxidation state of the metal} in Roman numerals in parentheses immediately after the metal name (no space): iron(III), copper(II), ferrate(II).
\item \textbf{Name the counter-ions} (if any) after the complex cation or before the complex anion, as for simple ionic compounds.
\end{enumerate}
\end{methode}

\begin{attention}[Common Naming Traps]
\begin{itemize}
\item \textbf{Alphabetical order ignores prefixes.} \textit{ammine} (a) comes before \textit{chloro} (c) even if you have \textit{triammine} and \textit{dichloro}.
\item \textbf{Oxidation state $\neq$ charge on the complex ion.} The oxidation state is the charge the metal \emph{would have} if all ligands were removed with their donor electrons. Calculate it from the overall charge of the complex sphere.
\item \textbf{The '-ate' ending is only for anionic complexes.} \ce{K4[Fe(CN)6]} $\rightarrow$ potassium hexacyanoferrate(II); but \ensuremath{\mathrm{[Fe(CN)_{6}]^{3-}}} $\rightarrow$ hexacyanoferrate(III) ion.
\item \textbf{Ligand names:} \ce{NH3} is \textit{ammine} (two m's), not amine; \ce{H2O} is \textit{aqua}, not aquo.
\item \textbf{Ambidentate ligands:} Specify the donor atom: \ce{NO2-} = nitro (N) vs nitrito (O); \ce{SCN-} = thiocyanato-S vs isothiocyanato-N; \ce{CN-} = cyano (C) vs isocyano (N).
\end{itemize}
\end{attention}

\begin{exoresolu}[Worked Naming Examples]
\textbf{Name the following compounds:}
\begin{enumerate}
\item \ce{[Cu(H2O)6]SO4}
\item \ce{K4[Fe(CN)6]}
\item \ce{[Co(NH3)5Cl]Cl2}
\item \ce{[Cr(NH3)2Cl2(en)]Cl}  (en = ethylenediamine)
\end{enumerate}

\tcblower
\textbf{Solution:}
\begin{enumerate}
\item \textbf{Complex cation:} \ensuremath{\mathrm{[Cu(H_{2}O)_{6}]^{2+}}}. Ligands: six \textit{aqua} $\rightarrow$ hexaaqua. Metal: copper(II) (overall charge +2, ligands neutral). Counter-ion: sulfate. \textbf{Name:} \cle{hexaaquacopper(II) sulfate}.
\item \textbf{Complex anion:} \ensuremath{\mathrm{[Fe(CN)_{6}]^{4-}}}. Ligands: six \textit{cyano} $\rightarrow$ hexacyano. Metal: ferrate(II) (anionic complex, Fe oxidation state +2 because 6$\times$(-1) + x = -4 $\Rightarrow$ x = +2). Counter-ion: potassium. \textbf{Name:} \cle{potassium hexacyanoferrate(II)}.
\item \textbf{Complex cation:} \ensuremath{\mathrm{[Co(NH_{3})_{5}Cl]^{2+}}}. Ligands: five \textit{ammine}, one \textit{chloro} $\rightarrow$ alphabetical: ammine (a) before chloro (c) $\rightarrow$ pentaamminechloro. Metal: cobalt(III) (overall +2, 5$\times$0 + (-1) + x = +2 $\Rightarrow$ x = +3). Counter-ions: two chloride. \textbf{Name:} \cle{pentaamminechlorocobalt(III) chloride}.
\item \textbf{Complex cation:} \ce{[Cr(NH3)2Cl2(en)]^+}. Ligands: two \textit{ammine}, two \textit{chloro}, one \textit{ethylenediamine} $\rightarrow$ alphabetical: ammine (a), chloro (c), ethylenediamine (e) $\rightarrow$ diammine dichloro ethylenediamine. Use \textit{bis} for ethylenediamine? No, ethylenediamine doesn't contain a multiplicative prefix, so \textit{di-} is fine. But wait: ethylenediamine is a bidentate ligand, its name is "ethylenediamine" which contains "di", so we use \textit{bis(ethylenediamine)}. Actually, the rule is: if the ligand name is complex or contains a multiplicative prefix, use bis, tris, etc. Ethylenediamine contains "di", so we use \textit{bis(ethylenediamine)}. Metal: chromium(III) (overall +1, 2$\times$0 + 2$\times$(-1) + 0 + x = +1 $\Rightarrow$ x = +3). Counter-ion: chloride. \textbf{Name:} \cle{diamminedichloridobis(ethylenediamine)chromium(III) chloride}.
\end{enumerate}
\end{exoresolu}

\begin{exercice}[Naming Practice]
Name the following coordination compounds:
\begin{enumerate}
\item \ce{[Ni(NH3)6]Cl2}
\item \ce{Na2[NiCl4]}
\item \ensuremath{\mathrm{[Cr(H_{2}O)_{4}Cl_{2}]Cl\cdot_{2}H_{2}O}}
\item \ce{[Pt(NH3)2Cl2]} (neutral complex)
\item \ce{K3[Cr(C2O4)3]}
\item \ce{[Co(en)2(NO2)2]NO3}  (en = ethylenediamine)
\end{enumerate}
\end{exercice}

\courssub{Shapes of Complex Ions — Geometry and Coordination Number}

The \cle{coordination number} (CN) is the number of donor atoms attached to the metal centre. The geometry adopted minimises ligand–ligand repulsion (VSEPR-like arguments) and is influenced by the metal's d-electron count and ligand field effects. The four most common geometries in the first transition series are:

\begin{popfigure}
\begin{tikzpicture}[scale=1.1, every node/.style={font=\small}]
% Octahedral
\begin{scope}[xshift=-7cm]
\draw[popblue, thick] (0,0) -- (1.5,0) node[midway, below] {90°};
\draw[popblue, thick] (0,0) -- (-1.5,0);
\draw[popblue, thick] (0,0) -- (0,1.5);
\draw[popblue, thick] (0,0) -- (0,-1.5);
\draw[popblue, thick] (0,0) -- (1.06,1.06);
\draw[popblue, thick] (0,0) -- (-1.06,-1.06);
\fill[popblue] (0,0) circle (3pt) node[below right] {M};
\foreach \ang/\lbl in {0/L, 90/L, 180/L, 270/L, 45/L, 225/L} {
  \pgfmathsetmacro{\x}{1.5*cos(\ang)}
  \pgfmathsetmacro{\y}{1.5*sin(\ang)}
  \fill[popdark] (\x,\y) circle (2.5pt);
}
\node at (0,-2.3) {\textbf{Octahedral} (CN=6)};
\node at (0,-2.7) {e.g. \ensuremath{\mathrm{[Fe(H_{2}O)_{6}]^{2+}}}};
\end{scope}

% Tetrahedral
\begin{scope}[xshift=-2cm]
\pgfmathsetmacro{\a}{1.2}
\coordinate (M) at (0,0);
\coordinate (L1) at ({0.94*\a}, {0.33*\a});
\coordinate (L2) at ({-0.94*\a}, {0.33*\a});
\coordinate (L3) at ({0}, {-1.05*\a});
\coordinate (L4) at ({0}, {1.3*\a});
\draw[popgreen, thick] (M) -- (L1) node[midway, right] {109.5°};
\draw[popgreen, thick] (M) -- (L2);
\draw[popgreen, thick] (M) -- (L3);
\draw[popgreen, thick] (M) -- (L4);
\fill[popgreen] (M) circle (3pt) node[below right] {M};
\foreach \p in {L1,L2,L3,L4} \fill[popdark] (\p) circle (2.5pt);
\node at (0,-2.3) {\textbf{Tetrahedral} (CN=4)};
\node at (0,-2.7) {e.g. \ensuremath{\mathrm{[CoCl_{4}]^{2-}}}};
\end{scope}

% Square Planar
\begin{scope}[xshift=3cm]
\draw[poporange, thick] (-1.5,0) -- (1.5,0) node[midway, below] {90°};
\draw[poporange, thick] (0,-1.5) -- (0,1.5);
\fill[poporange] (0,0) circle (3pt) node[below right] {M};
\foreach \coord in {(-1.5,0),(1.5,0),(0,-1.5),(0,1.5)} \fill[popdark] \coord circle (2.5pt);
\node at (0,-2.3) {\textbf{Square Planar} (CN=4)};
\node at (0,-2.7) {e.g. \ce{[Pt(NH3)2Cl2]}};
\end{scope}

% Linear
\begin{scope}[xshift=8cm]
\draw[poppurple, thick] (-1.5,0) -- (1.5,0) node[midway, below] {180°};
\fill[poppurple] (0,0) circle (3pt) node[below right] {M};
\foreach \coord in {(-1.5,0),(1.5,0)} \fill[popdark] \coord circle (2.5pt);
\node at (0,-2.3) {\textbf{Linear} (CN=2)};
\node at (0,-2.7) {e.g. \ce{[Ag(NH3)2]+}};
\end{scope}
\end{tikzpicture}
\legende{Common geometries of transition metal complexes. Bond angles are indicated.}
\end{popfigure}

\begin{aretenir}[Key Geometries and Examples]
\begin{itemize}
\item \textbf{Octahedral (CN=6):} Six ligands at the vertices of an octahedron. All bond angles 90° (or 180° for trans pairs). Very common for \ce{d^6} (low spin), \ce{d^3}, \ce{d^8} (e.g. \ensuremath{\mathrm{[Co(NH_{3})_{6}]^{3+}}}, \ensuremath{\mathrm{[Cr(H_{2}O)_{6}]^{3+}}}, \ensuremath{\mathrm{[Ni(H_{2}O)_{6}]^{2+}}}).
\item \textbf{Tetrahedral (CN=4):} Four ligands at the vertices of a regular tetrahedron. Bond angles 109.5°. Favoured by \ce{d^0}, \ce{d^5} (high spin), \ce{d^10} metals and bulky ligands (e.g. \ensuremath{\mathrm{[CoCl_{4}]^{2-}}}, \ce{[FeCl4]-}, \ensuremath{\mathrm{[Zn(NH_{3})_{4}]^{2+}}}).
\item \textbf{Square Planar (CN=4):} Four ligands in a plane at the corners of a square. Bond angles 90°. Characteristic of \ce{d^8} metals with strong-field ligands: \ce{Pt^2+}, \ce{Pd^2+}, \ce{Au^3+}, \ce{Ni^2+} (e.g. \ce{[Pt(NH3)2Cl2]}, \ensuremath{\mathrm{[Ni(CN)_{4}]^{2-}}}).
\item \textbf{Linear (CN=2):} Two ligands opposite each other. Bond angle 180°. Typical for \ce{d^10} ions: \ce{Ag+}, \ce{Cu+}, \ce{Au+} (e.g. \ce{[Ag(NH3)2]+}, \ce{[CuCl2]-}).
\end{itemize}
\end{aretenir}

\courssub{Isomerism in Coordination Compounds}

\begin{definition}[Isomers]
\cle{Isomers} are compounds that have the \textbf{same molecular formula} but \textbf{different arrangements of atoms} in space or in the bonding network. In coordination chemistry, isomerism is classified into two broad categories:
\begin{enumerate}
\item \textbf{Structural (constitutional) isomerism} — different connectivity (which atoms are bonded to which).
\item \textbf{Stereoisomerism} — same connectivity but different spatial arrangement.
\end{enumerate}
\end{definition}

\courssub{Structural Isomerism}

Three important types occur in coordination compounds:

\begin{propriete}[Ionization Isomerism]
Arises when a ligand that can also act as a counter-ion exchanges places with the counter-ion outside the coordination sphere. The composition of the coordination sphere changes, so the ions precipitated from solution differ.

\textbf{Classic example:} \ce{[Co(NH3)5Br]SO4} (red-violet) and \ce{[Co(NH3)5SO4]Br} (red).
\begin{itemize}
\item In \ce{[Co(NH3)5Br]SO4}, the \cle{free ion} in solution is \ce{SO4^2-} $\rightarrow$ gives white \ce{BaSO4} precipitate with \ce{BaCl2}; no precipitate with \ce{AgNO3} (Br- is bound).
\item In \ce{[Co(NH3)5SO4]Br}, the free ion is \ce{Br-} $\rightarrow$ gives pale yellow \ce{AgBr} precipitate with \ce{AgNO3}; no precipitate with \ce{BaCl2}.
\end{itemize}
\end{propriete}

\begin{propriete}[Hydrate (Solvate) Isomerism]
A special case of ionization isomerism where water molecules act as ligands inside the coordination sphere or as water of crystallization outside. The three forms of \ce{CrCl3.6H2O} are the textbook example:

\begin{center}
\begin{tabularx}{\linewidth}{|l|X|X|X|}
\hline
\textbf{Formula (modern)} & \textbf{Traditional name} & \textbf{Free \ce{Cl-} per formula} & \textbf{Colour} \\
\hline
\ce{[Cr(H2O)6]Cl3} & \textit{hexaaquachromium(III) chloride} & 3 (all ionic) & Violet \\
\ce{[Cr(H2O)5Cl]Cl2.H2O} & \textit{pentaaquachloridochromium(III) chloride monohydrate} & 2 & Light green \\
\ce{[Cr(H2O)4Cl2]Cl.2H2O} & \textit{tetraaquadichloridochromium(III) chloride dihydrate} & 1 & Dark green \\
\hline
\end{tabularx}
\end{center}

Adding \ce{AgNO3} to aqueous solutions of each form precipitates 3, 2, or 1 mole of \ce{AgCl} per mole of compound respectively.
\end{propriete}

\begin{propriete}[Coordination (Linkage) Isomerism]
Arises from \cle{ambidentate ligands} that can bind through two different donor atoms. The connectivity differs: the ligand is attached to the metal via a different atom.

\textbf{Examples:}
\begin{itemize}
\item \ce{NO2-}: \textit{nitro} (N-bonded, \ce{M-NO2}) vs \textit{nitrito} (O-bonded, \ce{M-ONO})
\item \ce{SCN-}: \textit{thiocyanato} (S-bonded, \ce{M-SCN}) vs \textit{isothiocyanato} (N-bonded, \ce{M-NCS})
\item \ce{CN-}: \textit{cyano} (C-bonded) vs \textit{isocyano} (N-bonded)
\end{itemize}
\textbf{Example pair:} \ce{[Co(NH3)5NO2]Cl2} (nitropentaamminecobalt(III) chloride) and \ce{[Co(NH3)5ONO]Cl2} (nitritopentaamminecobalt(III) chloride).
\end{propriete}

\begin{attention}[Identifying Structural Isomers — Exam Tip]
To distinguish ionization/hydrate isomers, test the \textbf{free anions} in aqueous solution:
\begin{itemize}
\item Add \ce{AgNO3}(aq): precipitate $\rightarrow$ free \ce{Cl-}, \ce{Br-}, \ce{I-} present.
\item Add \ce{BaCl2}(aq): white precipitate $\rightarrow$ free \ce{SO4^2-} present.
\end{itemize}
The isomer that gives a precipitate with \ce{AgNO3} has the halide as counter-ion; the one that gives a precipitate with \ce{BaCl2} has sulfate as counter-ion.
For linkage isomers: IR spectroscopy distinguishes M–N vs M–O or M–S vs M–N bonds.
\end{attention}

\begin{experience}[Distinguishing the Three Hydrate Isomers of \ce{CrCl3.6H2O}]
\textbf{Procedure:} Dissolve equal molar amounts of each isomer in water. To each solution add excess \ce{AgNO3}(aq). Filter, dry and weigh the \ce{AgCl} precipitate.
\textbf{Observation:} The violet isomer gives 3 mol \ce{AgCl}, the light green gives 2 mol, the dark green gives 1 mol.
\textbf{Conclusion:} The number of \textit{free} chloride ions (outside the coordination sphere) differs, proving hydration isomerism.
\end{experience}

\courssub{Stereoisomerism — Geometric (Cis–Trans) Isomerism}

Geometric isomerism occurs when ligands can be arranged in different relative positions around the metal centre. It is possible in \textbf{square planar} (CN=4) and \textbf{octahedral} (CN=6) complexes, but \textbf{not} in tetrahedral or linear geometries (all positions equivalent).

\begin{popfigure}
\begin{tikzpicture}[scale=1.3]
% cis-Pt(NH3)2Cl2
\begin{scope}[xshift=-4.5cm]
\draw[poporange, thick] (-1,0) -- (1,0);
\draw[poporange, thick] (0,-1) -- (0,1);
\fill[poporange] (0,0) circle (3pt) node[below right] {Pt};
\fill[popblue] (-1,0) circle (2.5pt) node[left] {Cl};
\fill[popblue] (0,1) circle (2.5pt) node[above] {NH3};
\fill[popgreen] (1,0) circle (2.5pt) node[right] {NH3};
\fill[popgreen] (0,-1) circle (2.5pt) node[below] {Cl};
\node at (0,-1.9) {\textbf{cis}-[\ce{Pt(NH3)2Cl2}]};
\node at (0,-2.2) {Cl–Pt–Cl = 90°};
\node at (0,-2.5) {Dipole moment $\neq$ 0};
\end{scope}

% trans-Pt(NH3)2Cl2
\begin{scope}[xshift=4.5cm]
\draw[poporange, thick] (-1,0) -- (1,0);
\draw[poporange, thick] (0,-1) -- (0,1);
\fill[poporange] (0,0) circle (3pt) node[below right] {Pt};
\fill[popblue] (-1,0) circle (2.5pt) node[left] {Cl};
\fill[popblue] (1,0) circle (2.5pt) node[right] {Cl};
\fill[popgreen] (0,1) circle (2.5pt) node[above] {NH3};
\fill[popgreen] (0,-1) circle (2.5pt) node[below] {NH3};
\node at (0,-1.9) {\textbf{trans}-[\ce{Pt(NH3)2Cl2}]};
\node at (0,-2.2) {Cl–Pt–Cl = 180°};
\node at (0,-2.5) {Dipole moment = 0};
\end{scope}
\end{tikzpicture}
\legende{Geometric isomers of square planar \ce{[Pt(NH3)2Cl2]}. \textit{cis}-platin (left) and \textit{trans}-platin (right).}
\end{popfigure}

\begin{savaistu}[Cis-Platin: Geometry Saves Lives]
\textit{Cis}-platin (\textit{cis}-diamminedichloridoplatinum(II)) is one of the most successful \cle{anti-cancer drugs}, used to treat testicular, ovarian, bladder and lung cancers. It works by binding to DNA, forming cross-links that block replication. The \textit{trans} isomer is \textbf{inactive} — its geometry prevents the necessary bifunctional binding to adjacent guanine bases. This is a spectacular example of how \cle{stereochemistry controls biological activity}. The discovery of cis-platin's activity (Rosenberg, 1965) revolutionised chemotherapy and earned a place in every inorganic chemistry syllabus.
\end{savaistu}

\begin{propriete}[Geometric Isomerism in Octahedral Complexes]
For \ce{MA4B2} type (e.g. \ce{[Co(NH3)4Cl2]+}):
\begin{itemize}
\item \textit{cis}: two B ligands adjacent (90°).
\item \textit{trans}: two B ligands opposite (180°).
\end{itemize}
For \ce{MA3B3} type (e.g. \ce{[Cr(NH3)3Cl3]}):
\begin{itemize}
\item \textit{fac} (facial): three identical ligands occupy one face of the octahedron (all mutually cis).
\item \textit{mer} (meridional): three identical ligands lie in a plane passing through the metal (one trans pair).
\end{itemize}
For \ce{M(AA)2B2} type (AA = bidentate ligand, e.g. \ce{[Co(en)2Cl2]+}):
\begin{itemize}
\item \textit{cis}: two B adjacent. Chiral (exists as enantiomers).
\item \textit{trans}: two B opposite. Achiral (centre of symmetry).
\end{itemize}
\end{propriete}

\begin{popfigure}
\begin{tikzpicture}[scale=1.1]
% fac isomer
\begin{scope}[xshift=-4.5cm]
\draw[popblue, thick] (0,0) -- (1.5,0);
\draw[popblue, thick] (0,0) -- (0,1.5);
\draw[popblue, thick] (0,0) -- (-1.06,-1.06);
\fill[popblue] (0,0) circle (3pt) node[below right] {M};
\fill[poporange] (1.5,0) circle (2.5pt) node[right] {A};
\fill[poporange] (0,1.5) circle (2.5pt) node[above] {A};
\fill[poporange] (-1.06,-1.06) circle (2.5pt) node[below left] {A};
\fill[popgreen] (-1.5,0) circle (2.5pt) node[left] {B};
\fill[popgreen] (0,-1.5) circle (2.5pt) node[below] {B};
\fill[popgreen] (1.06,1.06) circle (2.5pt) node[above right] {B};
\node at (0,-2.3) {\textbf{fac}-[\ce{MA3B3}]};
\node at (0,-2.6) {A's on one face};
\end{scope}

% mer isomer
\begin{scope}[xshift=4.5cm]
\draw[popblue, thick] (0,0) -- (1.5,0);
\draw[popblue, thick] (0,0) -- (0,1.5);
\draw[popblue, thick] (0,0) -- (-1.06,-1.06);
\fill[popblue] (0,0) circle (3pt) node[below right] {M};
\fill[poporange] (1.5,0) circle (2.5pt) node[right] {A};
\fill[poporange] (-1.5,0) circle (2.5pt) node[left] {A};
\fill[poporange] (0,1.5) circle (2.5pt) node[above] {A};
\fill[popgreen] (0,-1.5) circle (2.5pt) node[below] {B};
\fill[popgreen] (1.06,1.06) circle (2.5pt) node[above right] {B};
\fill[popgreen] (-1.06,-1.06) circle (2.5pt) node[below left] {B};
\node at (0,-2.3) {\textbf{mer}-[\ce{MA3B3}]};
\node at (0,-2.6) {A's in a meridian};
\end{scope}
\end{tikzpicture}
\legende{Facial (\textit{fac}) and meridional (\textit{mer}) isomers of octahedral \ce{MA3B3} complexes.}
\end{popfigure}

\courssub{Stereoisomerism — Optical Isomerism (Enantiomerism)}

Optical isomers (\cle{enantiomers}) are non-superimposable mirror images of each other, like left and right hands. They rotate plane-polarised light in opposite directions. A complex is \cle{chiral} (optically active) if it lacks:
\begin{itemize}
\item a plane of symmetry ($\sigma$),
\item a centre of inversion ($i$),
\item an improper rotation axis ($S_n$).
\end{itemize}

Octahedral complexes with \textbf{bidentate ligands} are the classic sources of optical isomerism. The tris-chelate \ce{[M(AA)3]} (AA = bidentate ligand like ethylenediamine, \ce{en}) exists as a pair of enantiomers ($\Lambda$ and $\Delta$).

\begin{popfigure}
\begin{tikzpicture}[scale=1.4]
% Lambda enantiomer (left-handed propeller)
\begin{scope}[xshift=-4.5cm]
\draw[poppurple, thick] (0,0) circle (1.2);
\foreach \i in {0,1,2} {
  \pgfmathsetmacro{\ang}{90 + \i*120}
  \pgfmathsetmacro{\x}{1.2*cos(\ang)}
  \pgfmathsetmacro{\y}{1.2*sin(\ang)}
  \draw[poppurple, thick] (0,0) -- (\x,\y);
  \fill[popdark] (\x,\y) circle (2pt);
  % chelate ring suggestion
  \pgfmathsetmacro{\angb}{\ang + 30}
  \pgfmathsetmacro{\xb}{1.2*cos(\angb)}
  \pgfmathsetmacro{\yb}{1.2*sin(\angb)}
  \draw[poppurple, dashed, thin] (\x,\y) -- (\xb,\yb);
  \fill[popdark] (\xb,\yb) circle (2pt);
}
\fill[poppurple] (0,0) circle (3pt) node[below right] {M};
\node at (0,-1.9) {$\Lambda$-\textbf{[M(en)3]$^{3}$+}};
\node at (0,-2.2) {Left-handed propeller};
\end{scope}

% Delta enantiomer (right-handed propeller)
\begin{scope}[xshift=4.5cm]
\draw[poppurple, thick] (0,0) circle (1.2);
\foreach \i in {0,1,2} {
  \pgfmathsetmacro{\ang}{90 - \i*120}
  \pgfmathsetmacro{\x}{1.2*cos(\ang)}
  \pgfmathsetmacro{\y}{1.2*sin(\ang)}
  \draw[poppurple, thick] (0,0) -- (\x,\y);
  \fill[popdark] (\x,\y) circle (2pt);
  \pgfmathsetmacro{\angb}{\ang - 30}
  \pgfmathsetmacro{\xb}{1.2*cos(\angb)}
  \pgfmathsetmacro{\yb}{1.2*sin(\angb)}
  \draw[poppurple, dashed, thin] (\x,\y) -- (\xb,\yb);
  \fill[popdark] (\xb,\yb) circle (2pt);
}
\fill[poppurple] (0,0) circle (3pt) node[below right] {M};
\node at (0,-1.9) {$\Delta$-\textbf{[M(en)3]$^{3}$+}};
\node at (0,-2.2) {Right-handed propeller};
\end{scope}

% Mirror plane
\draw[popline, dashed, thick] (0,-0.5) -- (0,-2.9) node[midway, right] {Mirror};
\end{tikzpicture}
\legende{Optical isomers (enantiomers) of \ensuremath{\mathrm{[M(en)_{3}]^{3+}}} ($\Lambda$ and $\Delta$). They are non-superimposable mirror images.}
\end{popfigure}

\begin{aretenir}[Optical Isomerism Checklist]
\begin{itemize}
\item \textbf{Requires:} Octahedral geometry + at least two bidentate ligands (or three) OR specific unsymmetrical tetradentate ligands.
\item \textbf{Common chiral formulas:} \ce{[M(AA)3]}, \ce{[M(AA)2B2]} (\textit{cis} only), \ce{[M(AA)2BC]} (\textit{cis}), \ce{[M(ABCD)EF]} (rare).
\item \textbf{Never chiral:} \textit{trans} isomers (have a centre of symmetry), tetrahedral complexes with monodentate ligands (rapid racemisation), square planar complexes.
\item \textbf{Separation:} Enantiomers have identical physical properties (except optical rotation) and chemical properties towards achiral reagents. They can be separated by chromatography on a chiral stationary phase or by forming diastereomeric salts with a chiral resolving agent.
\end{itemize}
\end{aretenir}

\begin{exoresolu}[Identifying Isomer Types]
\textbf{For each pair, state the type of isomerism and suggest a simple chemical test to distinguish them.}
\begin{enumerate}
\item \ce{[Co(NH3)5Cl]SO4} and \ce{[Co(NH3)5SO4]Cl}
\item \ce{[Cr(H2O)6]Cl3} and \ce{[Cr(H2O)5Cl]Cl2.H2O}
\item \textit{cis}- and \textit{trans}-[\ce{Pt(NH3)2Cl2}]
\item $\Lambda$-[\ce{Cr(en)3}]$^{3+}$ and $\Delta$-[\ce{Cr(en)3}]$^{3+}$
\item \ce{[Co(NH3)5NO2]Cl2} and \ce{[Co(NH3)5ONO]Cl2}
\end{enumerate}

\tcblower
\textbf{Solution:}
\begin{enumerate}
\item \textbf{Ionization isomerism.} Test: Add \ce{AgNO3}(aq) $\rightarrow$ first gives no precipitate (Cl- bound), second gives \ce{AgCl}(s). Add \ce{BaCl2}(aq) $\rightarrow$ first gives \ce{BaSO4}(s), second gives no precipitate.
\item \textbf{Hydrate isomerism.} Test: Add \ce{AgNO3}(aq) to equal molar solutions. First gives 3 mol \ce{AgCl}, second gives 2 mol \ce{AgCl} per mol of compound.
\item \textbf{Geometric (cis–trans) isomerism.} Test: Measure dipole moment (\textit{cis} has $\mu > 0$, \textit{trans} has $\mu = 0$). Or test biological activity: \textit{cis}-platin binds DNA, \textit{trans} does not.
\item \textbf{Optical isomerism.} Test: Polarimetry — they rotate plane-polarised light equally but in opposite directions. Chemical properties identical towards achiral reagents.
\item \textbf{Linkage (coordination) isomerism.} Test: IR spectroscopy — \ce{M-NO2} (nitro) shows N–O stretches ~1300-1400 cm$^{-1}$; \ce{M-ONO} (nitrito) shows asymmetric O–N–O stretch ~1000-1100 cm$^{-1}$.
\end{enumerate}
\end{exoresolu}

\courssub{Summary of Isomerism in Coordination Compounds}

\begin{center}
\begin{tabularx}{\linewidth}{|l|X|X|}
\hline
\textbf{Type} & \textbf{Key Feature} & \textbf{Examples} \\
\hline
\textbf{Ionization} & Counter-ion exchanges with ligand in coordination sphere & \ce{[Co(NH3)5Br]SO4} vs \ce{[Co(NH3)5SO4]Br} \\
\textbf{Hydrate} & Water as ligand vs water of crystallization & \ce{[Cr(H2O)6]Cl3}, \ce{[Cr(H2O)5Cl]Cl2.H2O}, \ce{[Cr(H2O)4Cl2]Cl.2H2O} \\
\textbf{Coordination (Linkage)} & Ambidentate ligand binds through different donor atoms & \ensuremath{\mathrm{[Co(NH_{3})_{5}NO_{2}]^{2+}}} (nitro) vs \ensuremath{\mathrm{[Co(NH_{3})_{5}ONO]^{2+}}} (nitrito); \ce{SCN-} (thiocyanato-S vs isothiocyanato-N) \\
\textbf{Geometric (cis–trans)} & Different relative positions (square planar, octahedral) & \textit{cis}-/\textit{trans}-[\ce{Pt(NH3)2Cl2}], \textit{fac}-/\textit{mer}-[\ce{Cr(NH3)3Cl3}] \\
\textbf{Optical} & Non-superimposable mirror images (chiral) & $\Lambda$- and $\Delta$-[\ce{Cr(en)3}]$^{3+}$, \textit{cis}-[\ce{Co(en)2Cl2}]$^{+}$ \\
\hline
\end{tabularx}
\end{center}

\begin{attention}[Exam Checklist for Isomerism Questions]
\begin{itemize}
\item \textbf{Draw structures clearly} — use wedge/dash or proper 3D sketches for octahedral complexes.
\item \textbf{State the \emph{type} of isomerism} precisely (e.g. "geometric isomerism", not just "isomerism").
\item \textbf{For ionization/hydrate isomers:} specify which ions are \textit{free} in solution and give the \textbf{observation with \ce{AgNO3} and \ce{BaCl2}}.
\item \textbf{For linkage isomers:} name the ligand correctly (nitro vs nitrito, thiocyanato vs isothiocyanato) and mention IR spectroscopy.
\item \textbf{For optical isomers:} mention "non-superimposable mirror images", "chiral", "rotate plane-polarised light in opposite directions".
\item \textbf{Don't confuse} geometric and optical: \textit{cis}-[\ce{M(AA)2B2}] is chiral (optical isomers exist); \textit{trans}-[\ce{M(AA)2B2}] is achiral.
\item \textbf{Count total isomers:} e.g. \ce{[M(AA)2B2]} has 3 stereoisomers (1 \textit{trans} + 2 \textit{cis} enantiomers).
\end{itemize}
\end{attention}

\begin{exercice}[Comprehensive Isomerism Exercise]
\begin{enumerate}
\item Draw all possible isomers (geometric and optical) for \ce{[Co(en)2Cl2]+}. Label each as \textit{cis} or \textit{trans}, and indicate which are chiral.
\item A compound \ce{CrCl3.6H2O} (isomer A) gives 3 mol \ce{AgCl} per mole with \ce{AgNO3}. Isomer B gives 1 mol \ce{AgCl}. Write the modern formulas for A and B and name them.
\item \ce{[Pt(NH3)2Cl2]} exists as two geometric isomers. Which one has a net dipole moment? Explain why the other has zero dipole moment.
\item The complex \ensuremath{\mathrm{[Ni(CN)_{4}]^{2-}}} is square planar while \ensuremath{\mathrm{[NiCl_{4}]^{2-}}} is tetrahedral. Explain this difference in terms of the ligand field strength and the d-electron configuration of \ce{Ni^2+}.
\item How many stereoisomers exist for \ensuremath{\mathrm{[Cr(ox)_{3}]^{3-}}} (ox = oxalate, bidentate)? Draw them and label their absolute configurations ($\Lambda$/$\Delta$).
\end{enumerate}
\end{exercice}

\begin{corrige}[Exercise 1]
\textbf{1.} \ce{[Co(en)2Cl2]+} (octahedral, \ce{en} = bidentate):
\begin{itemize}
\item \textit{trans} isomer: two Cl opposite. Has a centre of symmetry $\rightarrow$ \textbf{achiral}. One structure.
\item \textit{cis} isomer: two Cl adjacent (90°). No plane of symmetry $\rightarrow$ \textbf{chiral}. Exists as a pair of enantiomers ($\Lambda$ and $\Delta$).
\item Total: \textbf{3 stereoisomers} (1 \textit{trans} + 2 \textit{cis} enantiomers).
\end{itemize}

\textbf{2.} Isomer A: \ce{[Cr(H2O)6]Cl3} — hexaaquachromium(III) chloride. Isomer B: \ce{[Cr(H2O)4Cl2]Cl.2H2O} — tetraaquadichloridochromium(III) chloride dihydrate.

\textbf{3.} \textit{cis}-[\ce{Pt(NH3)2Cl2}] has a net dipole moment because the two \ce{Pt-Cl} bond dipoles and two \ce{Pt-NH3} bond dipoles do not cancel (vector sum $\neq$ 0). \textit{trans}-[\ce{Pt(NH3)2Cl2}] has zero dipole moment because the two \ce{Pt-Cl} dipoles are opposite and equal (cancel), and the two \ce{Pt-NH3} dipoles are opposite and equal (cancel).

\textbf{4.} \ce{Ni^2+} is \ce{d^8}. \ce{CN-} is a \textbf{strong-field ligand} $\rightarrow$ large $\Delta$ $\rightarrow$ low-spin configuration favours square planar (all electrons paired in lower orbitals, \ensuremath{\mathrm{d_{x^{2-}y^{2}}}} empty). \ce{Cl-} is a \textbf{weak-field ligand} $\rightarrow$ small $\Delta$ $\rightarrow$ high-spin $\rightarrow$ tetrahedral geometry avoids pairing energy penalty.

\textbf{5.} \ensuremath{\mathrm{[Cr(ox)_{3}]^{3-}}} is of type \ce{[M(AA)3]}. It exists as a pair of enantiomers: $\Lambda$ (left-handed propeller) and $\Delta$ (right-handed propeller). Total = \textbf{2 optical isomers}. No geometric isomers possible.
\end{corrige}

With naming, shapes and isomerism mastered, we have the language and the spatial intuition to discuss the remaining hallmark properties of transition metal complexes: \textbf{colour}, \textbf{catalysis}, \textbf{magnetism} and their \textbf{technological applications} — the subject of our final section.

\coursec{Colour, Catalysis, Magnetism and Uses of d-Block Elements}

In the previous section we mastered the nomenclature, geometries and isomerism of coordination complexes. We now turn to the \emph{consequences} of the partially filled d‑subshell that defines a transition element. It is this incomplete d‑shell that gives rise to the vivid colours, remarkable catalytic power, magnetic behaviour and the myriad industrial uses of the d‑block metals and their compounds. Understanding these properties not only explains everyday observations — from the blue of copper(II) sulfate to the platinum in a car’s exhaust — but also provides the conceptual toolkit that GCE examiners love to test.

\courssub{Origin of Colour in Transition Metal Compounds}

\begin{popfigure}
\begin{tikzpicture}[scale=1.1, every node/.style={font=\small}]
% Axes and labels
\draw[->] (-3.5,0) -- (3.5,0) node[right] {$E$};
\draw (-3.5,0) node[left] {Free ion};
\draw[->] (-3.5,3) -- (3.5,3) node[right] {$E$};
\draw (-3.5,3) node[left] {Octahedral field};

% Free ion d orbitals (degenerate)
\foreach \y in {0.5,1.0,1.5} {
    \draw[thick, popdark] (-3, \y) -- (-1, \y);
}
\node at (-2, 2.2) {$d$ orbitals (degenerate)};

% Octahedral splitting
% t2g set (lower)
\foreach \y in {0.5,1.0,1.5} {
    \draw[thick, popblue] (1, \y) -- (3, \y);
}
\node at (2, 0.0) {$t_{2g}$};
% eg set (higher)
\foreach \y in {2.2,2.7} {
    \draw[thick, poporange] (1, \y) -- (3, \y);
}
\node at (2, 3.1) {$e_g$};

% Splitting energy label
\draw[<->, poppurple] (3.2, 1.5) -- (3.2, 2.45) node[midway, right] {$\Delta_o$};

% d-d transition arrow
\draw[->, very thick, popgreen] (2, 1.5) -- (2, 2.2) node[midway, right] {\cle{d–d transition}};

% Electron configuration example: d^1
\foreach \y in {0.5} {
    \draw[->, thick, popink] (1.2, \y+0.1) -- (1.2, \y+0.4);
}
\node at (-2, -0.5) {Example: $d^1$ ion (e.g. \ce{Ti^{3+}})};
\end{tikzpicture}
\legende{d‑Orbital splitting in an octahedral ligand field. The five degenerate d‑orbitals split into a lower‑energy $t_{2g}$ triplet ($d_{xy}, d_{yz}, d_{zx}$) and a higher‑energy $e_g$ doublet ($d_{z^2}, d_{x^2-y^2}$). A d‑electron can absorb a photon of visible light and jump from $t_{2g}$ to $e_g$ (d–d transition). The energy gap $\Delta_o$ corresponds to the wavelength absorbed; the complementary colour is observed.}
\end{popfigure}

When ligands approach a transition metal ion, the electrostatic field they create lifts the degeneracy of the five d‑orbitals. In an octahedral complex the orbitals split into a lower‑energy set ($t_{2g}$: $d_{xy}, d_{yz}, d_{zx}$) and a higher‑energy set ($e_g$: $d_{z^2}, d_{x^2-y^2}$). The energy gap is denoted $\Delta_o$ (crystal field splitting parameter). If the metal ion has a partially filled d‑subshell ($d^1$–$d^9$), an electron in a lower orbital can absorb a photon of visible light and be promoted to a higher orbital — a \cle{d–d transition}. The wavelength absorbed depends on $\Delta_o$; the colour we see is the \emph{complementary} colour of the absorbed light.

\begin{propriete}[Origin of Colour]
Transition metal compounds are coloured because d‑electrons absorb specific wavelengths of visible light when jumping between split d‑orbitals (d–d transitions). The colour observed is the complementary colour of the absorbed light.
\end{propriete}

If the d‑subshell is empty ($d^0$) or completely full ($d^{10}$), no d–d transition is possible and the compound is typically colourless (or white as a solid). This explains why \ce{Sc^{3+}} ($3d^0$) and \ce{Zn^{2+}} ($3d^{10}$) form colourless solutions and white salts. Note that \ce{Sc^{3+}} and \ce{Zn^{2+}} are d‑block elements but \emph{not} transition metals by the IUPAC definition (no partially filled d‑subshell in any stable oxidation state).

\begin{exemplebox}[{Why \ensuremath{\mathrm{[Cu(H_{2}O)_{6}]^{2+}}} is blue}]
The \ce{Cu^{2+}} ion has a $d^9$ configuration. In an octahedral aqua complex the single hole in the $e_g$ level allows a d–d transition from $t_{2g}$ to $e_g$. The absorption maximum lies in the orange‑red region ($\approx 600–800\ \mathrm{nm}$), so the transmitted light appears blue.
\end{exemplebox}

\courssub{Effect of Ligand on Colour}

The magnitude of $\Delta_o$ depends strongly on the nature of the ligand (spectrochemical series). Strong‑field ligands (e.g. \ce{CN-}, \ce{CO}, \ce{NH3}) cause a larger splitting than weak‑field ligands (e.g. \ce{H2O}, \ce{F-}, \ce{Cl-}). A larger $\Delta_o$ means higher energy (shorter wavelength) light is absorbed, shifting the observed colour.

\begin{propriete}[Spectrochemical Series (common ligands, increasing field strength)]
\ensuremath{\mathrm{I^{-} < Br^{-} < SCN^{-} (S) < Cl^{-} < F^{-} < OH^{-} < H_{2}O < NH_{3} < en < CN^{-} < CO}}
\end{propriete}

\begin{center}
\begin{tabularx}{\linewidth}{|l|X|X|X|}
\hline
\rowcolor{popblueL}
\textbf{Complex} & \textbf{Ligand field strength} & \textbf{$\Delta_o$} & \textbf{Observed colour} \\
\hline
\ensuremath{\mathrm{[Cu(H_{2}O)_{6}]^{2+}}} & Weak field (\ce{H2O}) & Moderate & Blue \\
\hline
\ensuremath{\mathrm{[Cu(NH_{3})_{4}(H_{2}O)_{2}]^{2+}}} & Stronger field (\ce{NH3}) & Larger & Deep blue (indigo) \\
\hline
\ensuremath{\mathrm{[CuCl_{4}]^{2-}}} & Weak field (\ce{Cl-}), tetrahedral & Smaller ($\Delta_t \approx \frac{4}{9}\Delta_o$) & Yellow‑green \\
\hline
\end{tabularx}
\end{center}

The tetrahedral \ensuremath{\mathrm{[CuCl_{4}]^{2-}}} ion has a much smaller splitting, so it absorbs lower‑energy (red) light and appears yellow‑green. This ligand‑dependent colour change is a classic qualitative analysis observation.

\begin{attention}[Charge‑Transfer Colour]
Some intensely coloured complexes (e.g. \ce{MnO4-}, \ce{Cr2O7^{2-}}, \ensuremath{\mathrm{[Fe(SCN)]^{2+}}}) owe their colour to \emph{ligand‑to‑metal charge transfer (LMCT)} or \emph{metal‑to‑ligand charge transfer (MLCT)}, not d–d transitions. These occur when the metal has a high oxidation state ($d^0$ or $d^{10}$) or low‑lying empty ligand orbitals. They are not forbidden by selection rules and give very intense colours.
\end{attention}

\courssub{Catalytic Activity of Transition Metals}

Transition metals and their compounds are outstanding catalysts because they can:
\begin{enumerate}
\item \textbf{Adopt multiple oxidation states} — providing alternative reaction pathways with lower activation energies via redox cycles.
\item \textbf{Adsorb reactants on their surface} — using vacant d‑orbitals to form temporary bonds, weakening reactant bonds and increasing collision efficiency (heterogeneous catalysis).
\end{enumerate}

\begin{experience}[Industrial Catalysts — Key Examples for GCE]
\begin{itemize}
\item \textbf{Haber process (ammonia synthesis):} Finely divided \cle{Fe} (with \ce{K2O}/\ce{Al2O3} promoters) catalyses \ce{N2(g) + 3H2(g) <=> 2NH3(g)}. Fe cycles between oxidation states and provides a surface for \ce{N#N} bond cleavage.
\item \textbf{Contact process (sulfuric acid):} \cle{V2O5} catalyses \ce{2SO2(g) + O2(g) <=> 2SO3(g)}. The vanadium cycles between +5 and +4 oxidation states:
      \ce{V2O5 + SO2 -> V2O4 + SO3}, \quad \ce{V2O4 + 1/2 O2 -> V2O5}.
\item \textbf{Hydrogenation of oils:} \cle{Ni} (or Pt/Pd) catalyses addition of \ce{H2} to C=C bonds in unsaturated fats. \ce{H2} dissociates on the metal surface; H atoms add to the adsorbed alkene.
\item \textbf{Catalytic converters:} \cle{Pt}/\cle{Rh} (on a ceramic honeycomb) oxidise \ce{CO} and hydrocarbons to \ce{CO2} and \ce{H2O}, and reduce \ce{NO_x} to \ce{N2} and \ce{O2}. The metals cycle between oxidation states (e.g. \ce{Pt^0}/\ce{Pt^{2+}}).
\end{itemize}
\end{experience}

\courssub{Magnetic Properties: Paramagnetism and Diamagnetism}

Substances respond differently to an applied magnetic field. The behaviour is dictated by the presence or absence of unpaired electrons.

\begin{definition}[Paramagnetism]
Paramagnetism is the weak attraction of a substance to a magnetic field due to the presence of unpaired electrons. Each unpaired electron acts as a tiny magnet (spin magnetic moment). The magnetic moment $\mu_s$ (spin‑only formula) is $\mu_s = \sqrt{n(n+2)}\ \mu_B$, where $n$ = number of unpaired electrons and $\mu_B$ = Bohr magneton.
\end{definition}

\begin{definition}[Diamagnetism]
Diamagnetism is the weak repulsion of a substance from a magnetic field. It occurs in \emph{all} matter but is only observable when there are \emph{no} unpaired electrons (all electrons paired). A diamagnetic substance is \emph{not} “non‑magnetic” — it is repelled.
\end{definition}

\begin{attention}[Diamagnetic $\neq$ Non‑magnetic]
A substance with all electrons paired is \cle{diamagnetic} — it is weakly \emph{repelled} by a magnetic field. Do not call it “non‑magnetic”; that term is ambiguous and incorrect in a physics/chemistry context.
\end{attention}

To predict magnetic behaviour we write the d‑electron configuration in \cle{electron‑in‑boxes} notation (Hund’s rule: maximise unpaired electrons before pairing). For octahedral complexes we must consider whether the ligand field is strong (low‑spin) or weak (high‑spin).

\begin{methode}[Predicting Magnetic Behaviour of Transition Metal Ions/Complexes]
\begin{enumerate}
\item Determine the oxidation state of the metal ion.
\item Write its d‑electron count ($d^n$).
\item Draw the five d‑orbitals as boxes. For an octahedral complex, split into $t_{2g}$ (3 boxes) and $e_g$ (2 boxes).
\item Decide high‑spin or low‑spin: weak‑field ligands $\rightarrow$ high‑spin (Hund’s rule across all five boxes); strong‑field ligands $\rightarrow$ low‑spin (fill $t_{2g}$ completely before $e_g$).
\item Fill the boxes: one electron per box before pairing.
\item Count unpaired electrons.
\item \textbf{Unpaired $>0$ $\rightarrow$ paramagnetic}; \textbf{Unpaired $=0$ $\rightarrow$ diamagnetic}.
\end{enumerate}
\end{methode}

\begin{popfigure}
\begin{tikzpicture}[scale=0.9, every node/.style={font=\small}]
% Fe3+ d5 high-spin
\node[align=center] at (0, 2.5) {\textbf{$Fe^{3+}$ ($d^5$) -- high spin} \\ Paramagnetic (5 unpaired)};
\foreach \i in {0,...,4} {
    \draw (0.5+\i*1.2, 1.5) rectangle (1.5+\i*1.2, 2.0);
    \draw[->, thick, popblue] (1.0+\i*1.2, 1.7) -- (1.0+\i*1.2, 1.9);
}
\node at (0.5, 1.5) {$d_{xy}$};
\node at (1.7, 1.5) {$d_{yz}$};
\node at (2.9, 1.5) {$d_{zx}$};
\node at (4.1, 1.5) {$d_{z^2}$};
\node at (5.3, 1.5) {$d_{x^2-y^2}$};

% Fe3+ d5 low-spin
\node[align=center] at (0, 0.5) {\textbf{$Fe^{3+}$ ($d^5$) -- low spin} \\ Paramagnetic (1 unpaired)};
\foreach \i in {0,...,4} {
    \draw (0.5+\i*1.2, -0.5) rectangle (1.5+\i*1.2, 0.0);
}
\draw[->, thick, popblue] (1.0, -0.3) -- (1.0, -0.1);
\draw[->, thick, popblue] (2.2, -0.3) -- (2.2, -0.1);
\draw[->, thick, popblue] (3.4, -0.3) -- (3.4, -0.1);
\draw[->, thick, popblue] (4.6, -0.3) -- (4.6, -0.1);
\draw[->, thick, poporange] (4.6, -0.45) -- (4.6, -0.25);
\node at (0.5, -0.5) {$d_{xy}$};
\node at (1.7, -0.5) {$d_{yz}$};
\node at (2.9, -0.5) {$d_{zx}$};
\node at (4.1, -0.5) {$d_{z^2}$};
\node at (5.3, -0.5) {$d_{x^2-y^2}$};

% Zn2+ d10
\node[align=center] at (0, -1.5) {\textbf{$Zn^{2+}$ ($d^{10}$)} \\ Diamagnetic (0 unpaired)};
\foreach \i in {0,...,4} {
    \draw (0.5+\i*1.2, -2.5) rectangle (1.5+\i*1.2, -2.0);
    \draw[->, thick, popblue] (1.0+\i*1.2, -2.2) -- (1.0+\i*1.2, -2.0);
    \draw[->, thick, poporange] (1.0+\i*1.2, -2.35) -- (1.0+\i*1.2, -2.15);
}
\node at (0.5, -2.5) {$d_{xy}$};
\node at (1.7, -2.5) {$d_{yz}$};
\node at (2.9, -2.5) {$d_{zx}$};
\node at (4.1, -2.5) {$d_{z^2}$};
\node at (5.3, -2.5) {$d_{x^2-y^2}$};
\end{tikzpicture}
\legende{Electron‑in‑boxes diagrams. \textbf{Fe$^{3+}$} ($d^5$, high‑spin) has five unpaired electrons $\rightarrow$ strongly paramagnetic. \textbf{Fe$^{3+}$} ($d^5$, low‑spin) has one unpaired electron $\rightarrow$ paramagnetic. \textbf{Zn$^{2+}$} ($d^{10}$) has all electrons paired $\rightarrow$ diamagnetic.}
\end{popfigure}

\begin{exoresolu}[Predicting magnetic behaviour]
Predict the magnetic behaviour (paramagnetic/diamagnetic) and number of unpaired electrons for: (a) \ce{Cr^{3+}}, (b) \ce{Cu^{2+}}, (c) \ce{Zn^{2+}}, (d) \ensuremath{\mathrm{[Fe(CN)_{6}]^{4-}}}, (e) \ensuremath{\mathrm{[Fe(H_{2}O)_{6}]^{2+}}}.

\tcblower
\textbf{Solution:}
\begin{enumerate}
\item \textbf{Cr$^{3+}$}: Cr is [Ar] $3d^5 4s^1$. \ce{Cr^{3+}} loses $4s^1$ and two $3d$ electrons $\rightarrow$ $d^3$. High‑spin (weak field): three electrons occupy $t_{2g}$ singly $\rightarrow$ \textbf{3 unpaired} $\rightarrow$ \textbf{paramagnetic}. $\mu_s = \sqrt{15}\ \mu_B \approx 3.87\ \mu_B$.
\item \textbf{Cu$^{2+}$}: Cu is [Ar] $3d^{10} 4s^1$. \ce{Cu^{2+}} loses $4s^1$ and one $3d$ electron $\rightarrow$ $d^9$. In octahedral field: $t_{2g}^6 e_g^3$ $\rightarrow$ \textbf{1 unpaired} $\rightarrow$ \textbf{paramagnetic}. $\mu_s = \sqrt{3}\ \mu_B \approx 1.73\ \mu_B$.
\item \textbf{Zn$^{2+}$}: Zn is [Ar] $3d^{10} 4s^2$. \ce{Zn^{2+}} loses $4s^2$ $\rightarrow$ $d^{10}$. All orbitals full $\rightarrow$ \textbf{0 unpaired} $\rightarrow$ \textbf{diamagnetic}.
\item \textbf{[Fe(CN)6]$^{4-}$}: Fe oxidation state +2 ($d^6$). \ce{CN-} is strong field $\rightarrow$ low‑spin. $t_{2g}^6 e_g^0$ $\rightarrow$ \textbf{0 unpaired} $\rightarrow$ \textbf{diamagnetic}.
\item \textbf{[Fe(H2O)6]$^{2+}$}: Fe oxidation state +2 ($d^6$). \ce{H2O} is weak field $\rightarrow$ high‑spin. $t_{2g}^4 e_g^2$ $\rightarrow$ \textbf{4 unpaired} $\rightarrow$ \textbf{paramagnetic}. $\mu_s = \sqrt{24}\ \mu_B \approx 4.90\ \mu_B$.
\end{enumerate}
\end{exoresolu}

\courssub{Uses of d‑Block Elements and Compounds}

The unique combination of mechanical strength, electrical conductivity, variable oxidation states, catalytic activity and colour makes d‑block elements indispensable.

\begin{center}
\begin{tabularx}{\linewidth}{|l|X|}
\hline
\rowcolor{popblueL}
\textbf{Element / Compound} & \textbf{Key Uses} \\
\hline
\cle{Fe} (steel) & Construction (reinforced concrete, girders), vehicle bodies, machinery. Alloyed with C, Cr, Ni, Mn for tailored properties. \\
\hline
\cle{Cu} & Electrical wiring (high conductivity, ductility), plumbing, alloys (brass \ce{Cu-Zn}, bronze \ce{Cu-Sn}). \\
\hline
\cle{Ti} & Lightweight high‑strength alloys (aircraft, implants); \ce{TiO2} as brilliant white pigment in paints, plastics, sunscreen. \\
\hline
\cle{Cr} & Electroplating (decorative, corrosion‑resistant); stainless steel (\ce{Fe-Cr-Ni}); \ce{K2Cr2O7} as oxidising agent. \\
\hline
\cle{Ag} & Photography (\ce{AgBr} decomposition), jewellery, electrical contacts, antimicrobial coatings. \\
\hline
\cle{KMnO4} & Powerful oxidising agent in redox titrations, water treatment, disinfectant, qualitative analysis. \\
\hline
\cle{Ni} & Hydrogenation catalyst (margarine), rechargeable batteries (Ni‑Cd, Ni‑MH), stainless steel. \\
\hline
\cle{Pt}/\cle{Rh} & Catalytic converters, laboratory crucibles, chemotherapy drugs (cisplatin). \\
\hline
\end{tabularx}
\end{center}

\courssub{Summary Table: Characteristic Properties and the Partially Filled d‑Subshell}

All the hallmark properties of transition metals stem from the same root: an incomplete d‑subshell in at least one stable oxidation state.

\begin{center}
\begin{tabularx}{\linewidth}{|l|X|X|}
\hline
\rowcolor{popblueL}
\textbf{Property} & \textbf{Origin in partially filled d‑subshell} & \textbf{Typical GCE question} \\
\hline
\cle{Variable oxidation states} & Small energy gap between $(n-1)d$ and $ns$ orbitals; electrons can be lost from both. & “Explain why Mn shows oxidation states from +2 to +7.” \\
\hline
\cle{Formation of coloured compounds} & d–d transitions between split d‑orbitals absorb visible light. & “Why is \ce{ZnSO4} white while \ce{CuSO4} is blue?” \\
\hline
\cle{Catalytic activity} & Ability to change oxidation state (redox cycles) and to adsorb reactants on surface via vacant d‑orbitals. & “Describe the role of \ce{V2O5} in the Contact process.” \\
\hline
\cle{Paramagnetism} & Unpaired d‑electrons $\rightarrow$ net spin magnetic moment. & “Predict the magnetic behaviour of \ensuremath{\mathrm{[Fe(CN)_{6}]^{4-}}}.” \\
\hline
\cle{Complex formation} & Vacant low‑energy d‑orbitals accept lone pairs from ligands; crystal field stabilisation. & “Draw the shape of \ensuremath{\mathrm{[Ni(CN)_{4}]^{2-}}} and state its magnetic property.” \\
\hline
\end{tabularx}
\end{center}

\begin{aretenir}[Key Points for Revision]
\begin{itemize}
\item Colour arises from d–d transitions in split d‑orbitals; $d^0$ and $d^{10}$ ions are colourless (white solids).
\item Ligand field strength alters $\Delta_o$, changing the colour (spectrochemical series: \ensuremath{\mathrm{I^{-} < Br^{-} < Cl^{-} < F^{-} < H_{2}O < NH_{3} < CN^{-} < CO}}).
\item Tetrahedral splitting $\Delta_t \approx \frac{4}{9}\Delta_o$; tetrahedral complexes often have different colours.
\item Charge‑transfer (LMCT/MLCT) gives intense colours in some $d^0$/$d^{10}$ species (e.g. \ce{MnO4-}).
\item Catalysis exploits variable oxidation states (homogeneous) and surface adsorption (heterogeneous).
\item Know the major industrial catalysts: Fe (Haber), \ce{V2O5} (Contact), Ni (hydrogenation), Pt/Rh (converters).
\item Paramagnetism = unpaired electrons (attracted); diamagnetism = all electrons paired (repelled).
\item Use electron‑in‑boxes to count unpaired electrons; remember high‑spin vs low‑spin for octahedral complexes.
\item Spin‑only magnetic moment $\mu_s = \sqrt{n(n+2)}\ \mu_B$.
\item Common uses: Fe/steel, Cu/wiring, Ti/\ce{TiO2}, Cr/plating, Ag/photography, \ce{KMnO4}/oxidant.
\end{itemize}
\end{aretenir}

\begin{aretenir}[GCE Favourite]
A classic question: “\emph{Explain why zinc(II) compounds are white.}” The expected answer: “\ce{Zn^{2+}} has a $3d^{10}$ configuration — the d‑subshell is completely full, so no d–d transitions are possible; no visible light is absorbed.” Always mention the $d^{10}$ configuration explicitly. Another favourite: “\emph{State the magnetic behaviour of \ensuremath{\mathrm{[Fe(CN)_{6}]^{4-}}} and explain.}” Answer: “Diamagnetic; Fe(II) is $d^6$, \ce{CN-} is strong field $\rightarrow$ low‑spin $t_{2g}^6 e_g^0$, no unpaired electrons.”
\end{aretenir}

With this section we complete the descriptive inorganic chemistry of the halogens, Group IV and the d‑block. You now have a coherent framework linking electronic structure to the rich chemical behaviour that makes these elements central to both the laboratory and the modern world. Keep practising the electron‑in‑boxes diagrams and the qualitative analysis tables — they are the keys to high marks in Paper 2 and Paper 3.

\newpage

\coursec{Integration activity}

\coursec{Integrated Practical Activity: Qualitative Analysis of Halogens, Group IV and d-Block Elements}

\courssub{Aims of the Activity}
\begin{itemize}
    \item Consolidate and integrate knowledge from the whole chapter: halogen chemistry (halide tests, conc. \ce{H2SO4} reactions), Group IV chemistry (lead(II) chloride behaviour), and d-block chemistry (transition metal identification, complex formation, naming, shapes, isomerism, colour, magnetism).
    \item Develop practical planning skills: design a logical flowchart for qualitative analysis using limited reagents.
    \item Practice writing balanced ionic equations with state symbols and drawing mechanisms where appropriate.
    \item Apply IUPAC nomenclature for coordination compounds (oxidation states, ligand names, shapes, stereochemistry).
    \item Evaluate safety hazards associated with chlorine gas, lead compounds, and cyanide complexes in a Cameroonian school and domestic context.
\end{itemize}

\courssub{Situation}
\textbf{Context:} The chemistry laboratory of \emph{Lycée Bilingue d'Etoug-Ebe} in Yaoundé is preparing for the practical session of the GCE Advanced Level mock examinations. During the inventory, the laboratory technician discovers five unlabelled containers left on the preparation bench after the previous term's practicals. 

\textbf{The Unknowns:}
\begin{itemize}
    \item Four white crystalline solids labelled \textbf{Solid A}, \textbf{Solid B}, \textbf{Solid C}, and \textbf{Solid D}.
    \item One pale green aqueous solution labelled \textbf{Solution E}.
\end{itemize}

\textbf{The Suspects (known to be in the school's chemical store):}
\begin{enumerate}
    \item Sodium chloride (\ce{NaCl})
    \item Sodium bromide (\ce{NaBr})
    \item Lead(II) chloride (\ce{PbCl2})
    \item Potassium hexacyanoferrate(II) (\ce{K4[Fe(CN)6]})
    \item \emph{Either} Iron(II) sulfate (\ce{FeSO4}) \emph{or} Chromium(III) sulfate (\ce{Cr2(SO4)3}) — the label on the stock bottle is torn.
\end{enumerate}

\textbf{Available Reagents (standard school bench reagents):}
Dilute nitric acid (\ce{HNO3(aq)}), Silver nitrate solution (\ce{AgNO3(aq)}), Ammonia solution (\ce{NH3(aq)}) — both dilute and concentrated, Concentrated sulfuric acid (\ce{H2SO4(l)}), Sodium hydroxide solution (\ce{NaOH(aq)}), Bunsen burners, Nichrome wires, Test tubes, Safety goggles, Fume cupboard (functional).

\textbf{Your Mission:} As a team of Upper Sixth chemists, you are tasked with identifying all five unknowns conclusively. You must design a systematic experimental flowchart, predict every observation, write all relevant ionic equations, name any coordination complexes formed (giving their shapes and types of isomerism where applicable), and produce a final identification report. The activity concludes with a safety briefing on the handling of chlorine-based bleach (common in Cameroonian households as \emph{eau de Javel}) and lead compounds.

\courssub{Tasks}
\begin{enumerate}
    \item \textbf{Flowchart Design (Solids):} Draw a detailed flowchart showing the sequence of tests to distinguish Solids A–D. Start with solubility in water, then use \ce{AgNO3}/\ce{NH3} for halide identification, concentrated \ce{H2SO4} test for halide vs non-halide behaviour, and a specific test for the ferrocyanide ion. Indicate expected observations at each decision node.
    \item \textbf{Flowchart Design (Solution E):} Draw a flowchart to distinguish between \ce{Fe^{2+}(aq)} and \ce{Cr^{3+}(aq)}. Use \ce{NaOH(aq)} (dropwise then excess), \ce{NH3(aq)} (dropwise then excess), and a confirmatory test for \ce{Fe^{2+}} (e.g., oxidation to \ce{Fe^{3+}} followed by \ce{K4[Fe(CN)6]} or \ce{K3[Fe(CN)6]}).
    \item \textbf{Prediction Table:} For each solid (A–D) and Solution E, tabulate the expected observations for every test in your flowcharts.
    \item \textbf{Ionic Equations:} Write balanced ionic equations (with state symbols) for every positive reaction observed in Tasks 1 and 2. Include the formation of precipitates, dissolution of precipitates in excess reagent (complex formation), and redox reactions with conc. \ce{H2SO4}.
    \item \textbf{Complex Analysis:} For every coordination complex formed in your equations (e.g., \ce{[Ag(NH3)2]+}, \ensuremath{\mathrm{[PbCl_{4}]^{2-}}}, \ensuremath{\mathrm{[Fe(CN)_{6}]^{4-}}}, \ensuremath{\mathrm{[Cr(NH_{3})_{6}]^{3+}}}, \ensuremath{\mathrm{[Cr(OH)_{4}]^{-}}}, \ensuremath{\mathrm{[Fe(CN)_{6}]^{3-}}}), provide:
        \begin{itemize}
            \item The systematic IUPAC name.
            \item The oxidation state of the central metal.
            \item The coordination number and predicted shape (geometry).
            \item The type(s) of isomerism possible for that complex formula (geometric, optical, ionization, hydrate).
        \end{itemize}
    \item \textbf{Identification Report:} Based on the hypothetical observations below, identify Solids A–D and Solution E. Justify each identification by referencing the specific observation and the underlying chemical principle.
        \begin{quote}
        \textit{Hypothetical Observations:}
        \begin{itemize}
            \item \textbf{Solid A:} Dissolves in water. Gives white ppt with \ce{AgNO3}/\ce{HNO3}, soluble in dilute \ce{NH3}. With conc. \ce{H2SO4}, evolves colourless acidic gas (fumes in air), no colour change.
            \item \textbf{Solid B:} Dissolves in water. Gives cream ppt with \ce{AgNO3}/\ce{HNO3}, soluble in conc. \ce{NH3} but not dilute. With conc. \ce{H2SO4}, evolves brown/orange fumes and a gas that bleaches damp litmus.
            \item \textbf{Solid C:} Sparingly soluble in cold water, soluble in hot water. Gives white ppt with \ce{AgNO3}/\ce{HNO3}, insoluble in \ce{NH3}. With conc. \ce{H2SO4}, no gas evolution, dissolves on heating.
            \item \textbf{Solid D:} Dissolves in water. No ppt with \ce{AgNO3}/\ce{HNO3}. With conc. \ce{H2SO4}, no visible reaction. Adding \ce{Fe^{3+}} solution gives a deep blue precipitate (Prussian blue).
            \item \textbf{Solution E:} Pale green. With \ce{NaOH} (dropwise): dirty green ppt, insoluble in excess, turns brown on standing. With \ce{NH3} (dropwise): same dirty green ppt, insoluble in excess. Acidified solution + \ce{K3[Fe(CN)6]} gives a deep blue ppt.
        \end{itemize}
        \end{quote}
    \item \textbf{Safety Discussion:} Write a short briefing (150–200 words) for a school assembly in Yaoundé on the dangers of:
        \begin{itemize}
            \item Mixing household bleach (\emph{eau de Javel}, \ce{NaOCl}) with acidic toilet cleaners (HCl-based) or ammonia.
            \item Using lead-based pigments in traditional paints or pottery glazes.
            \item Improper disposal of potassium ferrocyanide waste.
        \end{itemize}
        Include the relevant chemical equations for the hazardous reactions.
\end{enumerate}

\courssub{Model Answer}

\courssub{Task 1: Flowchart for Solids A–D}
\begin{popfigure}
\centering
\begin{tikzpicture}[
    node distance=1.8cm and 2.5cm,
    decision/.style={diamond, draw=popdark, fill=poporange!20, thick, text width=3.2cm, align=center, inner sep=2pt},
    process/.style={rectangle, draw=popblue, fill=popblueL, thick, text width=3.5cm, align=center, inner sep=2pt, rounded corners=3pt},
    start/.style={ellipse, draw=popgreen, fill=popgreenL, thick, text width=2.5cm, align=center},
    arrow/.style={-Stealth, thick, popdark},
    yesno/.style={midway, fill=white, inner sep=2pt, font=\small, text width=2cm, align=center}
]
\node[start, align=center] (start){Four White Solids\\A, B, C, D};
\node[process, below of=start, align=center] (sol){Test 1: Solubility in\\cold water};
\node[decision, below of=sol, align=center] (d1){All dissolve?\\ (Note: C sparingly)};
\node[process, below left=1.5cm and 3cm of d1, align=center] (cold){Solid C: Sparingly soluble\\$\rightarrow$ Heat to dissolve};
\node[process, below right=1.5cm and 3cm of d1, align=center] (others){Solids A, B, D:\\Soluble};
\node[process, below of=others, align=center] (agno3){Test 2: Add dilute \ce{HNO3}then \ce{AgNO3(aq)}};
\node[decision, below of=agno3] (d2) {Precipitate formed?};
\node[process, below left=1.5cm and 3cm of d2, align=center] (noppt){Solid D: No ppt\\$\rightarrow$ Not a simple halide};
\node[decision, below right=1.5cm and 3cm of d2] (d3) {Ppt Colour?};
\node[process, below left=1.5cm and 2.5cm of d3, align=center] (white){White ppt\\$\rightarrow$ \ce{Cl-} (A or C)};
\node[process, below right=1.5cm and 2.5cm of d3, align=center] (cream){Cream ppt\\$\rightarrow$ \ce{Br-} (B)};
\node[process, below of=white, align=center] (nh3cl){Test 3a: Add \ce{NH3(aq)}\\dilute then conc.};
\node[decision, below of=nh3cl] (d4) {Solubility in \ce{NH3}?};
\node[process, below left=1.5cm and 2.5cm of d4, align=center] (sol_dil){Soluble in dilute \ce{NH3}\\$\rightarrow$ \ce{NaCl} (A)};
\node[process, below right=1.5cm and 2.5cm of d4, align=center] (insol_nh3){Insoluble in \ce{NH3}\\$\rightarrow$ \ce{PbCl2} (C)};
\node[process, below of=cream, align=center] (nh3br){Test 3b: Add \ce{NH3(aq)}\\dilute then conc.};
\node[decision, below of=nh3br] (d5) {Solubility in \ce{NH3}?};
\node[process, below left=1.5cm and 2.5cm of d5, align=center] (insol_dil){Insoluble in dilute,\\soluble in conc. \ce{NH3}\\$\rightarrow$ \ce{NaBr} (B)};
\node[process, below of=noppt, align=center] (ferri){Test 4 (Solid D):\\Add \ce{Fe^{3+}(aq)}};
\node[decision, below of=ferri] (d6) {Deep blue ppt?};
\node[process, below of=d6, align=center] (pruss){Yes $\rightarrow$ \ensuremath{\mathrm{[Fe(CN)_{6}]^{4-}}}\\$\rightarrow$ \ce{K4[Fe(CN)6]} (D)};
\node[process, below of=insol_nh3, align=center] (h2so4c){Test 5 (C): Conc. \ce{H2SO4}\\+ heat};
\node[process, below of=sol_dil, align=center] (h2so4a){Test 5 (A): Conc. \ce{H2SO4}\\cold};
\node[process, below of=insol_dil, align=center] (h2so4b){Test 5 (B): Conc. \ce{H2SO4}\\cold};

\draw[arrow] (start) -- (sol);
\draw[arrow] (sol) -- (d1);
\draw[arrow] (d1) -- node[yesno, left] {C: Sparingly} (cold);
\draw[arrow] (d1) -- node[yesno, right] {A, B, D: Yes} (others);
\draw[arrow] (cold) |- (agno3);
\draw[arrow] (others) -- (agno3);
\draw[arrow] (agno3) -- (d2);
\draw[arrow] (d2) -- node[yesno, left] {No} (noppt);
\draw[arrow] (d2) -- node[yesno, right] {Yes} (d3);
\draw[arrow] (d3) -- node[yesno, left] {White} (white);
\draw[arrow] (d3) -- node[yesno, right] {Cream} (cream);
\draw[arrow] (white) -- (nh3cl);
\draw[arrow] (nh3cl) -- (d4);
\draw[arrow] (d4) -- node[yesno, left] {Dilute} (sol_dil);
\draw[arrow] (d4) -- node[yesno, right] {No} (insol_nh3);
\draw[arrow] (cream) -- (nh3br);
\draw[arrow] (nh3br) -- (d5);
\draw[arrow] (d5) -- node[yesno, left] {Conc. only} (insol_dil);
\draw[arrow] (noppt) -- (ferri);
\draw[arrow] (ferri) -- (d6);
\draw[arrow] (d6) -- node[yesno, right] {Yes} (pruss);
\draw[arrow] (insol_nh3) -- (h2so4c);
\draw[arrow] (sol_dil) -- (h2so4a);
\draw[arrow] (insol_dil) -- (h2so4b);
\end{tikzpicture}
\legende{Flowchart for identification of Solids A–D.}
\end{popfigure}

\courssub{Task 2: Flowchart for Solution E}
\begin{popfigure}
\centering
\begin{tikzpicture}[
    node distance=1.8cm and 2.5cm,
    decision/.style={diamond, draw=popdark, fill=poporange!20, thick, text width=3.5cm, align=center, inner sep=2pt},
    process/.style={rectangle, draw=popblue, fill=popblueL, thick, text width=3.5cm, align=center, inner sep=2pt, rounded corners=3pt},
    start/.style={ellipse, draw=popgreen, fill=popgreenL, thick, text width=2.5cm, align=center},
    arrow/.style={-Stealth, thick, popdark},
    yesno/.style={midway, fill=white, inner sep=2pt, font=\small, text width=2.5cm, align=center}
]
\node[start, align=center] (start){Solution E\\Pale Green};
\node[process, below of=start, align=center] (naoh){Test 1: \ce{NaOH(aq)} dropwisethen excess};
\node[decision, below of=naoh] (d1) {Ppt colour \& behaviour?};
\node[process, below left=1.5cm and 3cm of d1, align=center] (fe_obs){Dirty green ppt \ce{Fe(OH)2}\\insoluble in excess\\turns brown on standing\\(\ce{Fe(OH)3})};
\node[process, below right=1.5cm and 3cm of d1, align=center] (cr_obs){Green/Grey-green ppt \ce{Cr(OH)3}\\soluble in excess giving\\green soln \ce{[Cr(OH)4]-}};
\node[process, below of=fe_obs, align=center] (nh3_fe){Test 2: \ce{NH3(aq)} dropwisethen excess};
\node[process, below of=cr_obs, align=center] (nh3_cr){Test 2: \ce{NH3(aq)} dropwisethen excess};
\node[decision, below of=nh3_fe] (d2) {Ppt soluble in excess?};
\node[decision, below of=nh3_cr] (d3) {Ppt soluble in excess?};
\node[process, below of=d2, align=center] (fe_nh3){No: \ce{Fe(OH)2} ppt remains\\$\rightarrow$ \ce{Fe^{2+}} confirmed};
\node[process, below of=d3, align=center] (cr_nh3){No: \ce{Cr(OH)3} ppt remains\\$\rightarrow$ \ce{Cr^{3+}} confirmed\\(\text{inert, no ammine in aq})};
\node[process, below of=fe_nh3, align=center] (confirm_fe){Test 3 (Confirm \ce{Fe^{2+}}):\\Acidify + \ce{K3[Fe(CN)6]}};
\node[decision, below of=confirm_fe, align=center] (d4){Deep blue ppt\\ (Turnbull's Blue)?};
\node[process, below of=d4] (fe_final) {Yes $\rightarrow$ \ce{Fe^{2+}} confirmed};

\draw[arrow] (start) -- (naoh);
\draw[arrow] (naoh) -- (d1);
\draw[arrow] (d1) -- node[yesno, left] {Dirty green $\to$ brown} (fe_obs);
\draw[arrow] (d1) -- node[yesno, right] {Green, sol. excess} (cr_obs);
\draw[arrow] (fe_obs) -- (nh3_fe);
\draw[arrow] (cr_obs) -- (nh3_cr);
\draw[arrow] (nh3_fe) -- (d2);
\draw[arrow] (nh3_cr) -- (d3);
\draw[arrow] (d2) -- node[yesno, right] {No} (fe_nh3);
\draw[arrow] (d3) -- node[yesno, right] {No} (cr_nh3);
\draw[arrow] (fe_nh3) -- (confirm_fe);
\draw[arrow] (confirm_fe) -- (d4);
\draw[arrow] (d4) -- node[yesno, right] {Yes} (fe_final);
\end{tikzpicture}
\legende{Flowchart for identification of Solution E (\ce{Fe^{2+}} vs \ce{Cr^{3+}}).}
\end{popfigure}

\courssub{Task 3: Prediction Table of Observations}

\begin{center}
\begin{tabularx}{\linewidth}{|l|X|X|X|X|X|}\hline
\rowcolor{popblueL}
\textbf{Test} & \textbf{Solid A (\ce{NaCl})} & \textbf{Solid B (\ce{NaBr})} & \textbf{Solid C (\ce{PbCl2})} & \textbf{Solid D (\ce{K4[Fe(CN)6]})} & \textbf{Solution E (\ce{Fe^{2+}})} \\ \hline
Solubility (cold) & Soluble, colourless & Soluble, colourless & Sparingly soluble & Soluble, colourless & \textemdash\ (already solution) \\ \hline
\ce{AgNO3}/\ce{HNO3} & White ppt (\ce{AgCl}) & Cream ppt (\ce{AgBr}) & White ppt (\ce{AgCl}) & No ppt & \textemdash \\ \hline
\ce{NH3} (dilute) & \ce{AgCl} dissolves & \ce{AgBr} insoluble & \ce{AgCl} insoluble & \textemdash & \textemdash \\ \hline
\ce{NH3} (conc.) & \textemdash & \ce{AgBr} dissolves & \ce{AgCl} insoluble & \textemdash & \textemdash \\ \hline
Conc. \ce{H2SO4} (cold) & \ce{HCl} gas (white fumes),\\ no colour change & \ce{Br2} vapour (brown),\\ \ce{SO2} (colourless),\\ bleaches litmus & No reaction (cold) & No visible reaction & \textemdash \\ \hline
Conc. \ce{H2SO4} (heat) & \textemdash & \textemdash & Dissolves on heating & Decomposes (\ce{HCN} gas,\\ toxic) & \textemdash \\ \hline
\ce{Fe^{3+}(aq)} & \textemdash & \textemdash & \textemdash & Deep blue ppt\\ (Prussian Blue) & \textemdash \\ \hline
\ce{NaOH} (dropwise) & \textemdash & \textemdash & \textemdash & \textemdash & Dirty green ppt\\ \ce{Fe(OH)2} \\ \hline
\ce{NaOH} (excess) & \textemdash & \textemdash & \textemdash & \textemdash & Insoluble, turns brown\\ (\ce{Fe(OH)3}) \\ \hline
\ce{NH3} (dropwise) & \textemdash & \textemdash & \textemdash & \textemdash & Dirty green ppt\\ \ce{Fe(OH)2} \\ \hline
\ce{NH3} (excess) & \textemdash & \textemdash & \textemdash & \textemdash & Insoluble \\ \hline
Acidified +\\ \ce{K3[Fe(CN)6]} & \textemdash & \textemdash & \textemdash & \textemdash & Deep blue ppt\\ (Turnbull's Blue) \\ \hline
\end{tabularx}
\end{center}

\courssub{Task 4: Ionic Equations}

\textbf{Solid A (\ce{NaCl}):}
\begin{align*}
\ensuremath{\mathrm{Ag^{+}(aq) + Cl^{-}(aq) }}&\ensuremath{\mathrm{\rightarrow  AgCl(s) \quad (white\ ppt)}} \\
\ce{AgCl(s) + 2NH3(aq) &-> [Ag(NH3)2]+(aq) + Cl-(aq)} \\
\ensuremath{\mathrm{NaCl(s) + H_{2}SO_{4}(l) }}&\ensuremath{\mathrm{\rightarrow  NaHSO_{4}(s) + HCl(g) \uparrow}} \quad \text{(white fumes with \ce{NH3})}
\end{align*}

\textbf{Solid B (\ce{NaBr}):}
\begin{align*}
\ensuremath{\mathrm{Ag^{+}(aq) + Br^{-}(aq) }}&\ensuremath{\mathrm{\rightarrow  AgBr(s) \quad (cream\ ppt)}} \\
\ce{AgBr(s) + 2NH3(aq) &-> [Ag(NH3)2]+(aq) + Br-(aq)} \quad \text{(requires conc. \ce{NH3})} \\
\ce{2NaBr(s) + 3H2SO4(l) &-> 2NaHSO4(s) + Br2(g) + SO2(g) + 2H2O(l)} \\
\ce{Br2(g) + H2O(l) &<=> HBr(aq) + HOBr(aq)} \quad \text{(bleaching action)}
\end{align*}

\textbf{Solid C (\ce{PbCl2}):}
\begin{align*}
\ce{PbCl2(s) &<=> Pb^{2+}(aq) + 2Cl-(aq)} \quad \text{(sparingly soluble, shifts right on heating)} \\
\ce{Ag+(aq) + Cl-(aq) &-> AgCl(s)} \quad \text{(from dissolved \ce{PbCl2})} \\
\ensuremath{\mathrm{PbCl_{2}(s) + 2Cl^{-}(aq) }}&\ensuremath{\mathrm{\rightarrow  [PbCl_{4}]^{2-}(aq)}} \quad \text{(in conc. \ce{Cl-} excess)} \\
\text{With conc. \ce{H2SO4} (heat):}\quad \ce{PbCl2(s) + H2SO4(l) &-> PbSO4(s) + 2HCl(g)} \quad \text{or dissolution in hot acid}
\end{align*}

\textbf{Solid D (\ce{K4[Fe(CN)6]}):}
\begin{align*}
\ensuremath{\mathrm{2Fe^{3+}(aq) + [Fe(CN)_{6}]^{4-}(aq) }}&\ensuremath{\mathrm{\rightarrow  Fe_{2}[Fe(CN)_{6}](s) \quad (Prussian\ Blue)}} \\
\text{With conc. \ce{H2SO4} (heat):}\quad \ensuremath{\mathrm{[Fe(CN)_{6}]^{4-}(aq) + 6H^{+}(aq) + 6H_{2}O(l) }}&\ensuremath{\mathrm{\rightarrow  Fe^{2+}(aq) + 6HCN(g) + 6H_{2}O(l)}} \quad \text{(Highly toxic \ce{HCN})}
\end{align*}

\textbf{Solution E (\ce{Fe^{2+}}):}
\begin{align*}
\ensuremath{\mathrm{Fe^{2+}(aq) + 2OH^{-}(aq) }}&\ensuremath{\mathrm{\rightarrow  Fe(OH)_{2}(s) \quad (dirty\ green)}} \\
\ensuremath{\mathrm{4Fe(OH)_{2}(s) + O_{2}(g) + 2H_{2}O(l) }}&\ensuremath{\mathrm{\rightarrow  4Fe(OH)_{3}(s) \quad (brown)}} \\
\ce{Fe^{2+}(aq) + 2NH3(aq) + 2H2O(l) &-> Fe(OH)2(s) + 2NH4+(aq)} \quad \text{(no dissolution in excess)} \\
\ensuremath{\mathrm{Fe^{2+}(aq) + [Fe(CN)_{6}]^{3-}(aq) }}&\ensuremath{\mathrm{\rightarrow  Fe[Fe(CN)_{6}](s) \quad (Turnbull's\ Blue)}} \\
\text{Note:}\quad \ensuremath{\mathrm{3[Fe(CN)_{6}]^{4-}(aq) + 4Fe^{3+}(aq) }}&\ensuremath{\mathrm{\rightarrow  Fe_{4}[Fe(CN)_{6}]_{3}(s) \quad (Prussian\ Blue)}} \\
\ensuremath{\mathrm{Fe^{2+}(aq) + [Fe(CN)_{6}]^{3-}(aq) }}&\ensuremath{\mathrm{\rightarrow  Fe^{3+}(aq) + [Fe(CN)_{6}]^{4-}(aq)}} \quad \text{(electron transfer in solid lattice)}
\end{align*}

\courssub{Task 5: Complex Analysis}

\begin{center}
\begin{tabularx}{\linewidth}{|l|c|c|c|X|X|}\hline
\rowcolor{popblueL}
\textbf{Complex Formula} & \textbf{Ox. State} & \textbf{Coord. No.} & \textbf{Shape} & \textbf{IUPAC Name} & \textbf{Possible Isomerism} \\ \hline
\ce{[Ag(NH3)2]+} & +1 & 2 & Linear & Diamminesilver(I) ion & None (CN=2) \\ \hline
\ensuremath{\mathrm{[PbCl_{4}]^{2-}}} & +2 & 4 & Tetrahedral & Tetrachloridoplumbate(II) ion & None (tetrahedral, identical ligands) \\ \hline
\ensuremath{\mathrm{[Fe(CN)_{6}]^{4-}}} & +2 & 6 & Octahedral & Hexacyanoferrate(II) ion & None (\ce{Ma6} type) \\ \hline
\ensuremath{\mathrm{[Cr(NH_{3})_{6}]^{3+}}} & +3 & 6 & Octahedral & Hexaamminechromium(III) ion & None (\ce{Ma6} type) \\ \hline
\ensuremath{\mathrm{[Cr(OH)_{4}]^{-}}} & +3 & 4 & Tetrahedral & Tetrahydroxidochromate(III) ion & None (tetrahedral, identical ligands) \\ \hline
\ensuremath{\mathrm{[Fe(CN)_{6}]^{3-}}} & +3 & 6 & Octahedral & Hexacyanoferrate(III) ion & None (\ce{Ma6} type) \\ \hline
\end{tabularx}
\end{center}

\textbf{Notes on Isomerism:}
\begin{itemize}
    \item \ce{[Ag(NH3)2]+} (Linear, CN=2): No geometric or optical isomerism possible.
    \item \ensuremath{\mathrm{[PbCl_{4}]^{2-}}}, \ensuremath{\mathrm{[Cr(OH)_{4}]^{-}}} (Tetrahedral, CN=4, \ce{Ma4} type): No geometric isomerism (all vertices equivalent). Optical isomerism only possible with four different ligands (chiral tetrahedral centre), not applicable here.
    \item \ensuremath{\mathrm{[Fe(CN)_{6}]^{4-}}}, \ensuremath{\mathrm{[Fe(CN)_{6}]^{3-}}}, \ensuremath{\mathrm{[Cr(NH_{3})_{6}]^{3+}}} (Octahedral, CN=6, \ce{Ma6} type): No geometric or optical isomerism.
    \item \textbf{Ionization Isomerism:} Possible for salts where a ligand can exchange with a counter-ion, e.g., \ce{[Co(NH3)5Br]SO4} (red) vs \ce{[Co(NH3)5SO4]Br} (violet). Not exhibited by the simple complex ions above.
    \item \textbf{Hydrate Isomerism:} Arises from different numbers of water molecules as ligands vs. water of crystallization, e.g., \ce{[Cr(H2O)6]Cl3} (violet), \ce{[Cr(H2O)5Cl]Cl2.H2O} (green), \ce{[Cr(H2O)4Cl2]Cl.2H2O} (dark green). Relevant for chromium(III) salts crystallised from aqueous solution.
    \item \textbf{Geometric (cis-trans) Isomerism:} Possible for \ce{Ma4b2} or \ce{Ma3b3} octahedral complexes (e.g., \ce{[Co(NH3)4Cl2]+}). Not applicable to the homoleptic complexes listed above.
    \item \textbf{Optical Isomerism:} Possible for octahedral complexes with bidentate ligands (\ce{[M(AA)3]}, \ce{[M(AA)2b2]}) or certain \ce{Ma2b2c2} types. Not applicable to the complexes listed above.
\end{itemize}

\courssub{Task 6: Identification Report}

\begin{aretenir}[Final Identification]
\begin{itemize}
    \item \textbf{Solid A is \ce{NaCl}.} Justification: Soluble in water; gives white \ce{AgCl} precipitate with \ce{AgNO3}/\ce{HNO3} soluble in \textbf{dilute} \ce{NH3} (high solubility product \ce{Ksp}(\ce{AgCl}) $\approx$ 1.8$\times$10$^{-10}$, complex formation constant \ce{Kf}(\ce{[Ag(NH3)2]+}) $\approx$ 1.6$\times$10$^{7}$ shifts equilibrium); conc. \ce{H2SO4} gives \ce{HCl} gas (white fumes with \ce{NH3}, no redox reaction as \ce{Cl-} is a weak reducing agent).
    \item \textbf{Solid B is \ce{NaBr}.} Justification: Soluble in water; gives cream \ce{AgBr} precipitate (\ce{Ksp} $\approx$ 5.0$\times$10$^{-13}$), insoluble in dilute \ce{NH3} but soluble in \textbf{conc.} \ce{NH3} (higher \ce{[NH3]} drives complex formation \ce{AgBr + 2NH3 <=> [Ag(NH3)2]+ + Br-}); conc. \ce{H2SO4} oxidises \ce{Br-} to \ce{Br2} (brown fumes) and is reduced to \ce{SO2} (\ce{Br-} is a stronger reducing agent than \ce{Cl-}, \ensuremath{\mathrm{E^\circ(\ce{Br2/Br^{-}})=+1.07\ V}} allows oxidation by \ce{H2SO4}/\ce{SO2} couple).
    \item \textbf{Solid C is \ce{PbCl2}.} Justification: Sparingly soluble in cold water (\ce{Ksp} $\approx$ 1.7$\times$10$^{-5}$ at 25°C), soluble in hot water (characteristic property); gives white \ce{AgCl} precipitate \textbf{insoluble in \ce{NH3}} (the \ce{Pb^{2+}} ions present precipitate as \ce{Pb(OH)2} in \ce{NH3}, physically occluding the \ce{AgCl} or preventing complexation under test conditions — a standard distinguishing observation for \ce{PbCl2} in this syllabus context); conc. \ce{H2SO4} causes no gas evolution (non-reducing \ce{Cl-}, \ce{Pb^{2+}} not oxidised), solid dissolves on heating (formation of soluble species or melting in acid medium).
    \item \textbf{Solid D is \ce{K4[Fe(CN)6]}.} Justification: Soluble in water; no precipitate with \ce{AgNO3}/\ce{HNO3} (\ce{Ag+} forms soluble \ce{[Ag(CN)2]-} complex instantly due to high \ce{Kf} $\approx$ 10$^{21}$ and excess \ce{CN-} from the salt); positive Prussian Blue test with \ce{Fe^{3+}} confirms ferrocyanide ion \ensuremath{\mathrm{[Fe(CN)_{6}]^{4-}}}.
    \item \textbf{Solution E contains \ce{Fe^{2+}}.} Justification: Pale green colour characteristic of \ensuremath{\mathrm{[Fe(H_{2}O)_{6}]^{2+}}}. \ce{NaOH}/\ce{NH3} give dirty green \ce{Fe(OH)2} precipitate, insoluble in excess (\ce{Fe(OH)2} is basic, not amphoteric). Oxidation by atmospheric oxygen to brown \ce{Fe(OH)3} confirms \ce{Fe(II)}. Turnbull's Blue precipitate with \ensuremath{\mathrm{[Fe(CN)_{6}]^{3-}}} confirms \ce{Fe^{2+}}. \ce{Cr^{3+}} would give green \ce{Cr(OH)3} soluble in excess \ce{NaOH} (amphoteric, forming \ce{[Cr(OH)4]-}) and insoluble in excess \ce{NH3} (kinetically inert).
\end{itemize}
\end{aretenir}

\courssub{Task 7: Safety Briefing for School Assembly}

\begin{experience}[Safety Briefing: Hidden Dangers in Our Homes and Labs]
\textbf{Theme:} "Chemistry for Safety: Protecting Our Families in Yaoundé"

\textbf{1. Bleach (\emph{Eau de Javel}) + Acid = Chlorine Gas (Deadly)}
Many households in Yaoundé use \emph{eau de Javel} (sodium hypochlorite, \ce{NaOCl}) for disinfecting floors and toilets. Some toilet cleaners contain hydrochloric acid (\ce{HCl}). \textbf{Never mix them.}
\begin{align*}
\ce{NaOCl(aq) + 2HCl(aq) &-> Cl2(g) + NaCl(aq) + H2O(l)}
\end{align*}
Chlorine gas (\ce{Cl2}) is a pulmonary agent used in WWI. Inhalation causes coughing, pulmonary oedema, and death. \textbf{Action:} Use one product, rinse thoroughly with water, then use the other. Ventilate always.

\textbf{2. Bleach + Ammonia = Chloramine / Hydrazine (Toxic)}
Some cleaning products contain ammonia (\ce{NH3}). Mixing with bleach produces chloramines (\ce{NH2Cl}, \ce{NHCl2}, \ce{NCl3}) and potentially explosive hydrazine (\ce{N2H4}).
\begin{align*}
\ce{NaOCl(aq) + NH3(aq) &-> NH2Cl(aq) + NaOH(aq)} \\
\ce{3NH2Cl(aq) + 2NH3(aq) &-> N2(g) + 3NH4Cl(aq)} \quad \text{(or \ce{N2H4} formation)}
\end{align*}
These gases cause severe respiratory distress and eye damage. \textbf{Action:} Read labels. Never mix cleaning agents.

\textbf{3. Lead Compounds: The Silent Poison}
Traditional pottery glazes or old paints may contain lead(II) oxide (\ce{PbO}) or lead chromate (\ce{PbCrO4}). Lead dust or fumes (from burning waste) are ingested/inhaled. Children are most vulnerable: lead inhibits haem synthesis and damages the developing brain (lower IQ, behavioural issues).
\begin{align*}
\ce{2PbO(s) + C(s) &-> 2Pb(l) + CO2(g)} \quad \text{(improvised smelting - DANGEROUS fumes)} \\
\ensuremath{\mathrm{Pb^{2+} + \text{biological ligands} }}&\ensuremath{\mathrm{\rightarrow  \text{Enzyme inhibition}}}
\end{align*}
\textbf{Action:} In school labs, always wear gloves handling \ce{PbCl2}, \ce{Pb(NO3)2}. Wash hands immediately. Dispose of lead waste in designated heavy-metal containers, \textbf{never} down the sink.

\textbf{4. Potassium Ferrocyanide: Safe if Cold, Deadly if Hot/Acid}
\ce{K4[Fe(CN)6]} is low toxicity (\ce{LD50} high) because \ce{CN-} is tightly bound to \ce{Fe^{2+}}. \textbf{BUT} heating with conc. \ce{H2SO4} or strong acid releases Hydrogen Cyanide gas (\ce{HCN}) — "the smell of bitter almonds" (genetic inability to smell in 20-40\% of people!).
\begin{align*}
\ensuremath{\mathrm{[Fe(CN)_{6}]^{4-}(aq) + 6H^{+}(aq) + 6H_{2}O(l) }}&\ensuremath{\mathrm{\rightarrow  Fe^{2+}(aq) + 6HCN(g) + 6H_{2}O(l)}}
\end{align*}
\textbf{Action:} Never heat ferro/ferricyanides with conc. acids. Store away from acids. In case of acid spill on cyanide salts: EVACUATE, ventilate, neutralise with \ce{NaOH} (forming \ce{NaCN}, still toxic but non-volatile) before cleanup by trained personnel.

\textbf{Conclusion:} Chemistry gives us comfort (clean water, medicine) but demands respect. "La chimie est une science de vie, pas un jeu." (Chemistry is a life science, not a game). Know your reagents. Read the MSDS. Protect your neighbour.
\end{experience}

\courssub{Examiner's Comments (for Teacher Reference)}
\begin{attention}[Marking Points Summary]
\begin{itemize}
    \item \textbf{Flowcharts (Tasks 1\&2):} Logical sequence (4 marks), correct reagents (2 marks), decision nodes with observations (4 marks). [Total 10]
    \item \textbf{Table (Task 3):} Accuracy of colours, solubility, gas evolution (1 mark per cell $\approx$ 15 marks).
    \item \textbf{Equations (Task 4):} Balancing, state symbols, correct redox products for \ce{H2SO4}/\ce{Br-} (2 marks each $\approx$ 12 marks).
    \item \textbf{Complexes (Task 5):} Name (1), Ox state (1), Shape (1), Isomerism analysis (1) per complex (24 marks).
    \item \textbf{Report (Task 6):} Correct ID (1 each), Justification linking obs to theory (2 each) [15 marks].
    \item \textbf{Safety (Task 7):} Equations (3), Hazards identified (3), Practical advice (2), Tone/Context (2) [10 marks].
    \item \textbf{Total:} $\approx$ 86 marks (scalable to 20-30 for a single session).
\end{itemize}
\end{attention}

\newpage

\coursec{Methods and worked examples}

\courssub{Method 1: Predicting and explaining trends in the halogen group}

\begin{methode}[Explaining physical trends down Group VII]
To answer any question on trends down the halogens (\ce{F2} $\rightarrow$ \ce{I2}), follow these reflexes:
\begin{enumerate}
\item \textbf{Identify the type of bonding.} The halogens are \cle{diatomic molecules} \ce{X2} held together in the solid/liquid by \emph{van der Waals (London dispersion) forces} only.
\item \textbf{Link the force to the number of electrons.} Down the group the number of electrons per molecule increases, so the electron cloud is larger and more easily polarised and the van der Waals forces become \textbf{stronger}.
\item \textbf{Deduce the physical state and melting/boiling point.} Stronger intermolecular forces $\Rightarrow$ higher melting and boiling points: \ce{F2} (pale yellow gas), \ce{Cl2} (greenish-yellow gas), \ce{Br2} (red-brown liquid), \ce{I2} (grey-black solid subliming to a violet vapour).
\item \textbf{Colour deepens down the group} because the energy gap for electronic absorption decreases, so absorption moves from the ultraviolet into the visible region.
\item \textbf{Reactivity decreases down the group} (opposite trend!): the atom gains an electron less easily because the incoming electron enters a shell further from the nucleus and is more shielded from the nuclear charge. Oxidising power: \ensuremath{\mathrm{F_{2} > Cl_{2} > Br_{2} > I_{2}}}.
\end{enumerate}
\textbf{Key idea to remember:} \emph{physical} trends (m.p., b.p., state) are governed by \textbf{intermolecular forces}; \emph{chemical} trends (oxidising power, reactivity) are governed by the \textbf{ease of gaining an electron}. Never confuse the two in an explanation.
\end{methode}

\begin{exoresolu}[Trends in the halogens]
\textbf{Statement (GCE style).} (a) State the colour and physical state of each of the halogens \ce{F2}, \ce{Cl2}, \ce{Br2} and \ce{I2} at room temperature. (b) Explain why the boiling points of the halogens increase down the group. (c) Explain why fluorine is the most reactive halogen although it has the weakest intermolecular forces.
\tcblower
\textbf{Solution.}

\textbf{(a)} Colours and states at room temperature:
\begin{center}
\begin{tabularx}{\linewidth}{|l|X|X|}\hline
\rowcolor{popblueL} \textbf{Halogen} & \textbf{Colour} & \textbf{Physical state} \\ \hline
\ce{F2} & Pale yellow & Gas \\ \hline
\ce{Cl2} & Greenish-yellow & Gas \\ \hline
\ce{Br2} & Red-brown & Volatile liquid (orange vapour) \\ \hline
\ce{I2} & Grey-black & Solid; sublimes to a violet vapour \\ \hline
\end{tabularx}
\end{center}

\textbf{(b)} The halogen molecules \ce{X2} are non-polar; the only forces between them are \textbf{van der Waals (London dispersion) forces}. Down the group the number of electrons in each molecule increases, the electron cloud becomes larger and more easily polarised, so the instantaneous dipole--induced dipole attractions become stronger. More energy is therefore needed to separate the molecules, and the boiling point increases: \ce{F2} $<$ \ce{Cl2} $<$ \ce{Br2} $<$ \ce{I2}.

\textbf{(c)} Reactivity of a halogen means its power as an \textbf{oxidising agent}, i.e.\ the ease with which the atom gains one electron to form \ce{X-}. Fluorine is the smallest halogen atom; the incoming electron enters the $n=2$ shell, very close to the nucleus, with little shielding, so it is strongly attracted: fluorine is the strongest oxidising agent of the group. The weakness of the van der Waals forces between \ce{F2} molecules concerns only the \emph{physical} state (a gas of low boiling point) and has nothing to do with the \emph{chemical} electron-gaining process. Hence physical and chemical trends need not run in parallel.
\end{exoresolu}

\courssub{Method 2: Displacement reactions and the oxidising power ladder}

\begin{methode}[Predicting halogen--halide displacement reactions]
\begin{enumerate}
\item \textbf{Rank the oxidising agents:} \ensuremath{\mathrm{F_{2} > Cl_{2} > Br_{2} > I_{2}}}. A halogen oxidises (displaces) the halide of any halogen \emph{below} it in the group.
\item \textbf{Write the half-equations.} Oxidation: \ce{2X- -> X2 + 2e-}; Reduction: \ensuremath{\mathrm{X'_{2} + 2e^{-} \rightarrow  2X'^{-}}}.
\item \textbf{Combine and cancel the electrons} to obtain the ionic equation, e.g.\ \ce{Cl2(aq) + 2Br-(aq) -> 2Cl-(aq) + Br2(aq)}.
\item \textbf{State the observation:} appearance of the colour of the displaced halogen (orange/brown for \ce{Br2}; brown solution, or a dark precipitate if concentrated, for \ce{I2}). Confirm with an organic solvent such as cyclohexane: an orange organic layer indicates \ce{Br2}, a violet layer indicates \ce{I2}.
\item \textbf{No reaction} occurs if the added halogen is \emph{below} the halide's halogen in the group (e.g.\ \ce{Br2} + \ce{Cl-}: no reaction).
\end{enumerate}
\textbf{Reflex:} the \emph{same} order applies to the reducing power of the halide ions, but reversed: \ensuremath{\mathrm{I^{-} > Br^{-} > Cl^{-} > F^{-}}}.
\end{methode}

\begin{exoresolu}[Displacement reactions of the halogens]
\textbf{Statement.} Chlorine water is added to aqueous potassium bromide; bromine water is added to aqueous potassium iodide; iodine is added to aqueous potassium chloride. For each experiment, state what is observed and write an ionic equation where a reaction occurs. Explain the results in terms of oxidising power.
\tcblower
\textbf{Solution.}

\textbf{Experiment 1: \ce{Cl2} + \ce{KBr}.} Chlorine is above bromine, hence a stronger oxidising agent; it oxidises \ce{Br-}:
\[ \ce{Cl2(aq) + 2Br-(aq) -> 2Cl-(aq) + Br2(aq)} \]
\textbf{Observation:} the solution turns orange-brown (bromine liberated); on shaking with cyclohexane an orange upper (organic) layer forms.

\textbf{Experiment 2: \ce{Br2} + \ce{KI}.} Bromine is above iodine:
\[ \ce{Br2(aq) + 2I-(aq) -> 2Br-(aq) + I2(aq)} \]
\textbf{Observation:} a brown solution of iodine appears (a dark solid may deposit if concentrated); with cyclohexane a violet upper layer forms.

\textbf{Experiment 3: \ce{I2} + \ce{KCl}.} Iodine is \emph{below} chlorine; it is a weaker oxidising agent than \ce{Cl2} and cannot oxidise \ce{Cl-}. \textbf{No reaction, no observable change.}

\textbf{Explanation.} A halogen can only remove electrons from the halide ion of a halogen \emph{below} it, because the outer electrons are held less tightly in the larger ion and are lost more easily. The experiments confirm the oxidising power order \ensuremath{\mathrm{Cl_{2} > Br_{2} > I_{2}}} and, equivalently, the reducing power order \ensuremath{\mathrm{I^{-} > Br^{-} > Cl^{-}}}.
\end{exoresolu}

\courssub{Method 3: Testing for halide ions}

\begin{methode}[Identifying \ce{Cl-}, \ce{Br-}, \ce{I-} in solution]
\begin{enumerate}
\item \textbf{Acidify} a portion of the solution with dilute \ce{HNO3} (this destroys interfering ions such as \ce{CO3^2-} that would also precipitate with \ce{Ag+}).
\item \textbf{Add \ce{AgNO3}(aq).} A precipitate of silver halide forms: \ce{Ag+(aq) + X-(aq) -> AgX(s)}.
\item \textbf{Read the colour:} \ce{AgCl} white, \ce{AgBr} cream, \ce{AgI} pale yellow.
\item \textbf{Confirm with ammonia solution:} \ce{AgCl} dissolves in \emph{dilute} \ce{NH3(aq)}; \ce{AgBr} dissolves only in \emph{concentrated} \ce{NH3(aq)}; \ce{AgI} is \emph{insoluble} even in concentrated ammonia.
\item Dissolution is due to complex formation: \ce{AgCl(s) + 2NH3(aq) -> [Ag(NH3)2]+(aq) + Cl-(aq)}.
\end{enumerate}
\textbf{Reflex:} never use \ce{HCl} to acidify — you would be adding chloride ions yourself and would obtain a white precipitate whatever the unknown!
\end{methode}

\begin{exoresolu}[Identifying an unknown halide]
\textbf{Statement.} A student is given a solution Q containing one halide ion. On adding dilute nitric acid followed by silver nitrate solution, a cream precipitate forms. The precipitate does not dissolve in dilute ammonia but dissolves in concentrated ammonia. (a) Identify the halide in Q, giving reasons. (b) Write ionic equations for the reactions observed. (c) Explain why nitric acid is added before the silver nitrate.
\tcblower
\textbf{Solution.}

\textbf{(a)} The cream colour of the silver halide precipitate suggests \ce{AgBr}. Confirmation: \ce{AgCl} (white) dissolves in \emph{dilute} ammonia, whereas \ce{AgI} (pale yellow) does not dissolve even in concentrated ammonia. A precipitate that resists dilute ammonia but dissolves in concentrated ammonia must be \ce{AgBr}. Hence Q contains the \textbf{bromide ion \ce{Br-}}.

\textbf{(b)} Precipitation:
\[ \ce{Ag+(aq) + Br-(aq) -> AgBr(s)} \]
Dissolution in concentrated ammonia (formation of the soluble diamminesilver(I) complex):
\[ \ce{AgBr(s) + 2NH3(aq) -> [Ag(NH3)2]+(aq) + Br-(aq)} \]

\textbf{(c)} Dilute \ce{HNO3} is added first to remove interfering anions, in particular carbonate ions \ce{CO3^2-}, which would otherwise give a precipitate of \ce{Ag2CO3} with \ce{Ag+} and falsify the test:
\[ \ce{2H+(aq) + CO3^2-(aq) -> H2O(l) + CO2(g)} \]
Nitric acid is chosen because it introduces no halide ion and because its own silver salt, \ce{AgNO3}, is soluble, so it cannot produce a precipitate.
\end{exoresolu}

\courssub{Method 4: Comparing the oxoacids of the halogens}

\begin{methode}[Acid strength, oxidising strength and thermal stability of oxoacids]
For the chlorine oxoacids \ce{HOCl}, \ce{HClO2}, \ce{HClO3}, \ce{HClO4} (oxidation states of Cl: $+1, +3, +5, +7$):
\begin{enumerate}
\item \textbf{Acid strength increases with the oxidation state of the halogen:} \ensuremath{\mathrm{HOCl < HClO_{2} < HClO_{3} < HClO_{4}}}. Reason: each additional oxygen withdraws electron density from the \ce{O-H} bond, weakening it, and stabilises the conjugate anion \ce{ClO_n-} by delocalisation of the negative charge over more oxygen atoms.
\item \textbf{Oxidising strength decreases} in the same direction: \ce{HOCl} is the strongest oxidising agent (chlorine in the low $+1$ state is easily reduced); dilute \ce{HClO4} is a comparatively weak oxidising agent.
\item \textbf{Thermal stability increases} with oxidation state: \ce{HOCl} decomposes readily (\ce{2HOCl -> 2HCl + O2}, accelerated by light); \ce{HClO4} is the most stable, although the pure concentrated acid can decompose explosively.
\item \textbf{Names (Stock notation):} \ce{HOCl} chloric(I) acid (hypochlorous acid), \ce{HClO2} chloric(III) acid (chlorous acid), \ce{HClO3} chloric(V) acid (chloric acid), \ce{HClO4} chloric(VII) acid (perchloric acid). The Roman numeral is the oxidation number of Cl.
\end{enumerate}
\textbf{Key formula:} for \ce{HClO_n}, the oxidation number of Cl is given by $1 + x - 2n = 0$, i.e.\ $x = 2n - 1$.
\end{methode}

\begin{exoresolu}[Oxoacids of chlorine]
\textbf{Statement.} (a) Give the systematic (Stock) name and the oxidation number of chlorine in each of the acids \ce{HOCl}, \ce{HClO3} and \ce{HClO4}. (b) Arrange the four acids \ce{HOCl}, \ce{HClO2}, \ce{HClO3}, \ce{HClO4} in order of increasing acid strength and justify the order. (c) Which of these acids is the strongest oxidising agent? Explain. (d) Write an equation for the decomposition of \ce{HOCl} in sunlight and state one use of this acid.
\tcblower
\textbf{Solution.}

\textbf{(a)} Using $1 + x - 2n = 0$:
\begin{center}
\begin{tabularx}{\linewidth}{|l|X|X|}\hline
\rowcolor{popblueL} \textbf{Formula} & \textbf{Oxidation number of Cl} & \textbf{Stock name} \\ \hline
\ce{HOCl} & $+1$ & Chloric(I) acid \\ \hline
\ce{HClO3} & $+5$ & Chloric(V) acid \\ \hline
\ce{HClO4} & $+7$ & Chloric(VII) acid \\ \hline
\end{tabularx}
\end{center}

\textbf{(b)} Increasing acid strength:
\[ \ensuremath{\mathrm{HOCl < HClO_{2} < HClO_{3} < HClO_{4}}} \]
As the number of oxygen atoms increases, the highly electronegative oxygens pull electron density away from the \ce{O-H} bond, weakening it and making loss of \ce{H+} easier. In addition, the negative charge of the conjugate base \ce{ClO_n-} is delocalised over more oxygen atoms, so the anion is increasingly stabilised and the acid increasingly strong. \ce{HClO4} is a very strong acid; \ce{HOCl} is a weak acid.

\textbf{(c)} \ce{HOCl} is the strongest oxidising agent: chlorine in the $+1$ state is easily reduced to \ce{Cl-} (oxidation state $-1$), and the \ce{O-Cl} bond is weak. Oxidising power therefore \emph{decreases} as the oxidation state of chlorine rises — the opposite order to acid strength.

\textbf{(d)} In sunlight:
\[ \ce{2HOCl(aq) -> 2HCl(aq) + O2(g)} \]
Use: \ce{HOCl} (and the \ce{OCl-} ion present in bleach) is used as a \textbf{bleaching agent and disinfectant}, e.g.\ in the treatment of drinking water, because of its strong oxidising (germ-killing) action.
\end{exoresolu}

\courssub{Method 5: Group IV — structures, allotropes and the inert pair effect}

\begin{methode}[Handling Group IV (C to Pb) questions]
\begin{enumerate}
\item \textbf{Electronic configuration} $ns^2np^2$ for all the elements; oxidation states $+4$ and $+2$. Down the group the $+2$ state becomes \textbf{more stable} (the \cle{inert pair effect}: the $ns^2$ pair is increasingly reluctant to take part in bonding). Hence \ce{PbO} is more stable than \ce{PbO2}; Pb(IV) compounds are strong oxidising agents, while Sn(II) compounds behave as reducing agents.
\item \textbf{Allotropy:} carbon exists as \textbf{diamond} (giant covalent, tetrahedral $sp^3$, very hard, non-conductor), \textbf{graphite} (layers of $sp^2$ hexagons, delocalised electrons within the layers $\Rightarrow$ conducts electricity, soft, lubricant) and fullerenes. Tin also shows allotropy (grey tin, non-metallic, and white tin, metallic).
\item \textbf{Metallic character increases down the group:} C (non-metal), Si and Ge (metalloids), Sn and Pb (metals). Atomic radius increases and first ionisation energy decreases down the group.
\item \textbf{Oxides:} \ce{CO2} is a simple molecular gas (acidic); \ce{SiO2} is a giant covalent solid of very high melting point (acidic); \ce{SnO2} and \ce{PbO2} are amphoteric with ionic character. \textbf{Chlorides:} \ce{CCl4} and \ce{SiCl4} are covalent liquids; \ce{SiCl4} is hydrolysed by water (\ce{SiCl4 + 2H2O -> SiO2 + 4HCl}) but \ce{CCl4} is not (carbon has no accessible orbital to accept a lone pair from water).
\item \textbf{Hydrides:} \ce{CH4}, \ce{SiH4}, \ce{GeH4}, \ce{SnH4}, \ce{PbH4} are tetrahedral; thermal stability \emph{decreases} down the group as the \ce{E-H} bond becomes longer and weaker.
\end{enumerate}
\end{methode}

\begin{exoresolu}[Group IV: structure and stability]
\textbf{Statement.} (a) Describe the structure and bonding in diamond and in graphite, and explain why graphite conducts electricity while diamond does not. (b) Explain why \ce{SiCl4} is hydrolysed by water whereas \ce{CCl4} is not. Write the equation for the hydrolysis of \ce{SiCl4}. (c) Explain why \ce{PbO2} oxidises concentrated hydrochloric acid to chlorine while \ce{CO2} shows no oxidising action.
\tcblower
\textbf{Solution.}

\textbf{(a)} In \textbf{diamond} each carbon atom is $sp^3$ hybridised and covalently bonded tetrahedrally to four other carbon atoms, giving a giant three-dimensional covalent lattice. All four outer electrons of each atom are localised in bonds, so there are no mobile charge carriers: diamond is a non-conductor (and is extremely hard). In \textbf{graphite} each carbon is $sp^2$ hybridised, bonded to three neighbours in flat hexagonal layers; the fourth electron of each atom is \textbf{delocalised} over the layer. These mobile electrons carry current, so graphite conducts electricity. The layers are held together by weak van der Waals forces and slide over one another, which makes graphite soft and useful as a lubricant and in pencil ``leads''.

\textbf{(b)} Silicon, in period 3, is a larger atom than carbon and has accessible empty $3d$ orbitals: a water molecule can donate a lone pair to silicon to form an intermediate, and the polar \ce{Si-Cl} bond is then readily attacked. Carbon has no available low-energy orbital (there is no $2d$ subshell) to accept the lone pair from water, so \ce{CCl4} is kinetically inert to hydrolysis even though the reaction would be energetically feasible.
\[ \ce{SiCl4(l) + 2H2O(l) -> SiO2(s) + 4HCl(aq)} \]
(White fumes of \ce{HCl} are observed.)

\textbf{(c)} This illustrates the \textbf{inert pair effect}. For lead, the $6s^2$ pair is held strongly and Pb(IV) is unstable relative to Pb(II); \ce{PbO2} therefore readily accepts electrons (it is reduced to \ce{Pb^{2+}}) and acts as a strong oxidising agent:
\[ \ce{PbO2(s) + 4HCl(aq) -> PbCl2(s) + Cl2(g) + 2H2O(l)} \]
For carbon, at the top of the group, the $+4$ state is very stable (the inert pair effect is negligible), so \ce{CO2} has no tendency to be reduced and shows no oxidising action under these conditions.
\end{exoresolu}

\courssub{Method 6: Writing electronic configurations of d-block atoms and ions}

\begin{methode}[Configurations of the first transition series (Sc to Zn)]
\begin{enumerate}
\item \textbf{Atoms:} fill $4s$ before $3d$ (Aufbau order). General pattern: $[\mathrm{Ar}]\,4s^2 3d^{n}$. \textbf{Two exceptions to memorise:} Cr $= [\mathrm{Ar}]\,4s^1 3d^5$ and Cu $= [\mathrm{Ar}]\,4s^1 3d^{10}$ (extra stability of the half-filled and fully filled $d$ subshells).
\item \textbf{Ions:} remove electrons from the \textbf{$4s$ orbital first}, then from $3d$. Example: \ce{Fe^2+} $= [\mathrm{Ar}]\,3d^6$ (not $4s^2 3d^4$!); \ce{Fe^3+} $= [\mathrm{Ar}]\,3d^5$.
\item \textbf{Electron-in-boxes notation:} draw five boxes for $3d$ and one for $4s$; fill with single parallel spins first (Hund's rule) before pairing.
\item \textbf{Definitions:} a \cle{d-block element} has its highest-energy (differentiating) electron in a $d$ subshell; a \cle{transition element} forms \emph{at least one ion with a partially filled $d$ subshell}. Hence Sc ($\ce{Sc^3+} = 3d^0$) and Zn ($\ce{Zn^2+} = 3d^{10}$) are d-block elements but \textbf{not} transition elements.
\item \textbf{Quick check:} the total number of electrons must equal $Z$ minus the positive charge.
\end{enumerate}
\end{methode}

\begin{exoresolu}[Electronic configurations of iron and copper species]
\textbf{Statement.} (a) Write the full electronic configuration (spdf notation) of \ce{Cr} ($Z=24$) and \ce{Cu} ($Z=29$), explaining any anomaly. (b) Write the configurations of \ce{Fe^2+} and \ce{Fe^3+} ($Z=26$) and explain why \ce{Fe^3+} is the more stable ion in aqueous solution. (c) Using electron-in-boxes notation for the $3d$ and $4s$ orbitals, represent \ce{Fe^2+}. (d) Explain why zinc is not regarded as a transition element.
\tcblower
\textbf{Solution.}

\textbf{(a)} Simple filling order would predict Cr $= 1s^2 2s^2 2p^6 3s^2 3p^6 4s^2 3d^4$ and Cu $= [\mathrm{Ar}]\,4s^2 3d^9$. The observed configurations are:
\[ \text{Cr} = 1s^2 2s^2 2p^6 3s^2 3p^6 4s^1 3d^5, \qquad \text{Cu} = 1s^2 2s^2 2p^6 3s^2 3p^6 4s^1 3d^{10} \]
A half-filled ($3d^5$) or completely filled ($3d^{10}$) $d$ subshell confers extra stability (symmetrical distribution of electron density and maximum exchange energy), so one $4s$ electron is effectively promoted into $3d$.

\textbf{(b)} Fe $= [\mathrm{Ar}]\,4s^2 3d^6$. Transition metals lose the $4s$ electrons first because, once the $3d$ subshell is occupied, the $4s$ electrons are on average further from the nucleus and are removed more easily:
\[ \ce{Fe^2+} = [\mathrm{Ar}]\,3d^6, \qquad \ce{Fe^3+} = [\mathrm{Ar}]\,3d^5 \]
\ce{Fe^3+} possesses the particularly stable half-filled $3d^5$ subshell, which is why iron(III) is favoured over iron(II) in aqueous (especially aerated) solution.

\textbf{(c)} \ce{Fe^2+}: the $4s$ box is empty; the six $3d$ electrons fill the five boxes singly first, then one pair forms:
\begin{center}
\begin{tikzpicture}[scale=0.55]
\foreach \x/\n in {0/2,1.2/2,2.4/1,3.6/1,4.8/1}{
  \draw (\x,0) rectangle (\x+0.8,0.8);
  \ifnum\n=2
    \draw[->,thick] (\x+0.25,0.1) -- (\x+0.25,0.7);
    \draw[->,thick] (\x+0.55,0.7) -- (\x+0.55,0.1);
  \else
    \draw[->,thick] (\x+0.4,0.1) -- (\x+0.4,0.7);
  \fi
}
\draw (7.2,0) rectangle (8.0,0.8);
\node at (2.8,-0.5) {$3d$};
\node at (7.6,-0.5) {$4s$};
\end{tikzpicture}
\end{center}

\textbf{(d)} A transition element must form at least one ion with a \emph{partially filled} $d$ subshell. Zinc's only common ion is \ce{Zn^2+} $= [\mathrm{Ar}]\,3d^{10}$, with a completely filled $d$ subshell; the Zn atom itself is $[\mathrm{Ar}]\,3d^{10}4s^2$. Since neither the atom nor its ion has an incomplete $d$ subshell, zinc is a d-block element but \textbf{not} a transition element. This is consistent with its properties: \ce{Zn^2+} compounds are white (colourless ions) and zinc shows only one oxidation state, $+2$.
\end{exoresolu}

\courssub{Method 7: Naming, shaping and finding isomers of complexes}

\begin{methode}[Naming complexes and predicting shapes and isomerism]
\begin{enumerate}
\item \textbf{Name the ligands first} (in alphabetical order), with prefixes di-, tri-, tetra- (or bis-, tris- for composite ligand names). Anionic ligands end in \emph{-o} (\ce{Cl-} chloro, \ce{CN-} cyano, \ce{OH-} hydroxo); neutral ligands keep their name except \ce{H2O} = aqua and \ce{NH3} = ammine.
\item \textbf{Then name the metal} with its oxidation state in Roman numerals; if the complex is an anion, the metal name ends in \emph{-ate} (e.g.\ ferrate, cuprate).
\item \textbf{Oxidation state:} (oxidation state of metal) $+$ (sum of ligand charges) $=$ overall charge of the complex ion.
\item \textbf{Shape from coordination number:} C.N.\ 6 $\Rightarrow$ \textbf{octahedral}; C.N.\ 4 $\Rightarrow$ \textbf{tetrahedral} (or square planar, e.g.\ \ce{[PtCl4]^2-}); C.N.\ 2 $\Rightarrow$ linear.
\item \textbf{Structural isomerism:} \emph{ionisation} isomerism (a ligand and an external ion exchange places, e.g.\ \ce{[Co(NH3)5Br]SO4} and \ce{[Co(NH3)5SO4]Br}); \emph{hydrate} isomerism (water inside vs outside the coordination sphere, e.g.\ \ce{[Cr(H2O)6]Cl3} and \ce{[Cr(H2O)5Cl]Cl2.H2O}).
\item \textbf{Stereoisomerism:} \emph{geometric} (cis/trans) isomerism in square planar \ce{MA2B2} and octahedral \ce{MA4B2} complexes; \emph{optical} isomerism (non-superimposable mirror images), e.g.\ cis-\ce{[Co(en)2Cl2]+} or \ce{[M(AA)3]} with bidentate ligands.
\end{enumerate}
\end{methode}

\begin{exoresolu}[Complexes: names, shapes, isomers]
\textbf{Statement.} (a) Name the following complexes and give the oxidation state of the metal: (i) \ce{[Cu(NH3)4]SO4}; (ii) \ce{K4[Fe(CN)6]}; (iii) \ce{[CoCl2(NH3)4]Cl}. (b) Draw and name the geometric isomers of the complex ion in (iii). (c) A compound of formula \ce{CrCl3.6H2O} exists in two forms: one gives three moles of \ce{AgCl} per mole with \ce{AgNO3}, the other gives only two. Identify the type of isomerism and write the formula of each form.
\tcblower
\textbf{Solution.}

\textbf{(a)} (i) The complex cation is \ce{[Cu(NH3)4]^2+}: $x + 4(0) = +2 \Rightarrow x = +2$. Name: \textbf{tetraamminecopper(II) sulphate}. Shape: coordination number 4 (approximately tetrahedral/square planar).

(ii) The complex anion is \ce{[Fe(CN)6]^4-}: $x + 6(-1) = -4 \Rightarrow x = +2$. Name: \textbf{potassium hexacyanoferrate(II)}. Shape: octahedral.

(iii) The complex cation is \ce{[CoCl2(NH3)4]+}: $x + 2(-1) + 4(0) = +1 \Rightarrow x = +3$. Name: \textbf{tetraamminedichlorocobalt(III) chloride} (ligands cited alphabetically: ammine before chloro). Shape: octahedral (C.N.\ 6).

\textbf{(b)} In the octahedral ion \ce{[CoCl2(NH3)4]+} the two \ce{Cl-} ligands can be:
\begin{itemize}
\item \textbf{cis}: adjacent (90° apart) — cis-tetraamminedichlorocobalt(III) ion;
\item \textbf{trans}: opposite (180° apart) — trans-tetraamminedichlorocobalt(III) ion.
\end{itemize}
\begin{center}
\begin{tikzpicture}[scale=0.9]
% cis
\node at (0,0) {Co};
\draw (0,0) -- (0,1.1) node[above]{Cl};
\draw (0,0) -- (1.1,0) node[right]{Cl};
\draw (0,0) -- (0,-1.1) node[below]{\ce{NH3}};
\draw (0,0) -- (-1.1,0) node[left]{\ce{NH3}};
\draw[dashed] (0,0) -- (0.65,0.55) node[above right=-2pt]{\ce{NH3}};
\draw (0,0) -- (-0.65,-0.55) node[below left=-2pt]{\ce{NH3}};
\node at (0,-1.8) {\emph{cis}};
% trans
\begin{scope}[xshift=5.5cm]
\node at (0,0) {Co};
\draw (0,0) -- (0,1.1) node[above]{Cl};
\draw (0,0) -- (0,-1.1) node[below]{Cl};
\draw (0,0) -- (1.1,0) node[right]{\ce{NH3}};
\draw (0,0) -- (-1.1,0) node[left]{\ce{NH3}};
\draw[dashed] (0,0) -- (0.65,0.55) node[above right=-2pt]{\ce{NH3}};
\draw (0,0) -- (-0.65,-0.55) node[below left=-2pt]{\ce{NH3}};
\node at (0,-1.8) {\emph{trans}};
\end{scope}
\end{tikzpicture}
\end{center}

\textbf{(c)} This is \textbf{hydrate isomerism} (a form of structural isomerism): water molecules and chloride ions exchange places between the coordination sphere and the outside of the complex. Only the \ce{Cl-} ions \emph{outside} the sphere are free and are precipitated by \ce{Ag+}.
\begin{itemize}
\item Form giving 3 mol \ce{AgCl}: all three \ce{Cl-} are outside the sphere: \ce{[Cr(H2O)6]Cl3}.
\item Form giving 2 mol \ce{AgCl}: one \ce{Cl-} is bonded inside the sphere (not precipitated): \ce{[Cr(H2O)5Cl]Cl2.H2O}.
\end{itemize}
(A third isomer, \ce{[Cr(H2O)4Cl2]Cl.2H2O}, would give only one mole of \ce{AgCl}.)
\end{exoresolu}

\courssub{Method 8: Colour, magnetism and catalysis in d-block chemistry}

\begin{methode}[Explaining colour and predicting magnetic behaviour]
\begin{enumerate}
\item \textbf{Colour:} in a complex, the ligands split the five $d$ orbitals into two sets of different energy (separation $\Delta$). A $d$ electron absorbs visible light of energy $E = h\nu = \Delta$ and is promoted to the higher set; the \emph{complementary} colour is transmitted and seen. Ions with $d^0$ (\ce{Sc^3+}) or $d^{10}$ (\ce{Zn^2+}, \ce{Cu+}) configurations have no possible $d$--$d$ transition and are \textbf{colourless}.
\item \textbf{Changing the ligand changes $\Delta$}, and hence the colour: e.g.\ \ce{[Cu(H2O)6]^2+} is pale blue; adding excess \ce{NH3} gives the deep blue \ce{[Cu(NH3)4(H2O)2]^2+}.
\item \textbf{Magnetism:} count the \emph{unpaired} electrons from the $3d$ box diagram. Unpaired electrons $\Rightarrow$ \textbf{paramagnetic} (attracted into a magnetic field); all electrons paired $\Rightarrow$ \textbf{diamagnetic}. The spin-only magnetic moment is $\mu = \sqrt{n(n+2)}$ B.M., where $n$ is the number of unpaired electrons.
\item \textbf{Catalysis:} transition metals catalyse by (i) providing a surface on which reactants adsorb (heterogeneous catalysis: Fe in the Haber process, Ni in the hydrogenation of oils, Pt/Rh in car exhaust converters, \ce{V2O5} in the Contact process) and (ii) using their \textbf{variable oxidation states} to provide an alternative reaction path of lower activation energy (e.g.\ \ce{Fe^2+}/\ce{Fe^3+} catalysing the reaction between \ce{S2O8^2-} and \ce{I-}).
\end{enumerate}
\end{methode}

\begin{exoresolu}[Colour, magnetism and catalysis]
\textbf{Statement.} (a) Explain why aqueous solutions of \ce{Cu^2+} salts are blue whereas solutions of \ce{Zn^2+} salts are colourless. (b) Calculate the spin-only magnetic moment of \ce{Fe^3+} ($3d^5$) and of \ce{Ni^2+} ($3d^8$) and state whether each ion is paramagnetic or diamagnetic. (c) State the catalyst used in (i) the Haber process, (ii) the Contact process, and explain in one sentence why transition metals make good catalysts.
\tcblower
\textbf{Solution.}

\textbf{(a)} \ce{Cu^2+} has the configuration $[\mathrm{Ar}]\,3d^9$: the $d$ subshell is \emph{partially filled}. In the octahedral aqua complex \ce{[Cu(H2O)6]^2+} the ligands split the $d$ orbitals into two energy levels separated by $\Delta$. Absorption of visible light (in the orange-red region) promotes a $d$ electron across $\Delta$; the complementary colour, blue, is observed. \ce{Zn^2+} is $3d^{10}$: the subshell is full, no $d$--$d$ promotion is possible, no visible light is absorbed, and the ion is colourless.

\textbf{(b)} \ce{Fe^3+} ($3d^5$): by Hund's rule the five electrons occupy the five $d$ orbitals singly, so $n = 5$ unpaired electrons:
\[ \mu = \sqrt{n(n+2)} = \sqrt{5 \times 7} = \sqrt{35} \approx 5.92\ \text{B.M.} \quad \Rightarrow \textbf{paramagnetic}. \]
\ce{Ni^2+} ($3d^8$): the eight electrons fill the five orbitals as three pairs plus two single electrons, so $n = 2$:
\[ \mu = \sqrt{2 \times 4} = \sqrt{8} \approx 2.83\ \text{B.M.} \quad \Rightarrow \textbf{paramagnetic}. \]

\textbf{(c)} (i) Haber process (\ce{N2 + 3H2 <=> 2NH3}): \textbf{iron} (with promoters). (ii) Contact process (\ce{2SO2 + O2 <=> 2SO3}): \textbf{vanadium(V) oxide, \ce{V2O5}}. Transition metals make good catalysts because they can \emph{change oxidation state} easily — offering an alternative pathway of lower activation energy (e.g.\ \ce{V2O5} cycling between V(V) and V(IV)) — and because their surfaces adsorb reactant molecules, weakening their bonds and holding them close together in favourable orientations.
\end{exoresolu}

\newpage

\coursec{Exercises}

\courssub{I check my knowledge}

\begin{exercice}[MCQ: Halogen Trends]
Which of the following statements concerning the halogens (Group 17) is \textbf{incorrect}?
\begin{enumerate}[label=\Alph*.]
\item The oxidizing power decreases down the group.
\item The bond dissociation energy of the X--X bond decreases steadily down the group.
\item The colour intensity increases down the group.
\item The first ionization energy decreases down the group.
\end{enumerate}
\end{exercice}

\begin{exercice}[True/False: Group IV Chemistry]
State whether each statement is True or False. Justify your answer with a brief explanation or a balanced chemical equation where appropriate.
\begin{enumerate}
\item Carbon tetrachloride (\ce{CCl4}) hydrolyses readily in water at room temperature to form \ce{CO2} and \ce{HCl}.
\item The stability of the +2 oxidation state in Group IV increases down the group relative to the +4 state.
\item Silicon dioxide (\ce{SiO2}) reacts with hot concentrated sodium hydroxide solution to form sodium silicate.
\item Lead(IV) oxide (\ce{PbO2}) is a stronger oxidizing agent than tin(IV) oxide (\ce{SnO2}).
\end{enumerate}
\end{exercice}

\begin{exercice}[Definitions and Configurations: d-Block]
\begin{enumerate}
\item Define the term \emph{transition element} according to the IUPAC definition.
\item Write the full electronic configuration (spdf notation) and the electron-in-boxes diagram for the \ce{3d} orbitals of:
      \begin{itemize}
      \item A chromium atom (\ce{Cr}, Z=24).
      \item A copper(II) ion (\ce{Cu^{2+}}, Z=29).
      \end{itemize}
\item Explain why the first ionization energy of zinc (\ce{Zn}) is higher than that of copper (\ce{Cu}).
\end{enumerate}
\end{exercice}

\begin{exercice}[Direct Calculations: Redox Titrations]
\SI{25.0}{cm^3} of a solution of iron(II) sulfate (\ce{FeSO4}) was acidified with dilute sulfuric acid and titrated against \SI{0.0200}{mol\,dm^{-3}} potassium permanganate (\ce{KMnO4}). The volume of permanganate required was \SI{22.5}{cm^3}.
\begin{enumerate}
\item Write the balanced ionic half-equations for the oxidation of \ce{Fe^{2+}} and the reduction of \ce{MnO4^-} in acidic medium.
\item Deduce the overall balanced ionic equation for the reaction.
\item Calculate the concentration, in \ensuremath{\mathrm{mol\,dm^{-3}}}, of the iron(II) sulfate solution.
\item Calculate the mass concentration of \ce{FeSO4.7H2O} in \ensuremath{\mathrm{g\,dm^{-3}}} (Molar mass = \SI{278}{g\,mol^{-1}}).
\end{enumerate}
\end{exercice}

\courssub{I practise}

\begin{exercice}[Halogen Displacement and Oxoacids]
\begin{enumerate}
\item A student adds chlorine water to an aqueous solution of potassium bromide, then adds an equal volume of tetrachloromethane (\ce{CCl4}) and shakes the mixture. Describe the observations in the aqueous layer and the organic layer. Write the balanced ionic equation for the reaction occurring.
\item Write the formulas and systematic names for the oxoacids of chlorine corresponding to oxidation states +1, +3, +5, and +7.
\item Arrange the oxoacids in (ii) in order of increasing acid strength and increasing thermal stability. Explain the trend in acid strength in terms of the structure of the conjugate bases.
\end{enumerate}
\end{exercice}

\begin{exercice}[Halide Ion Identification]
Three test tubes labelled A, B, and C each contain an aqueous solution of a sodium halide (\ce{NaF}, \ce{NaCl}, \ce{NaBr}, or \ce{NaI}). The following tests were performed:
\begin{itemize}
\item Test 1: Addition of \ce{AgNO3(aq)} followed by dilute \ce{NH3(aq)}.
\item Test 2: Addition of \ce{AgNO3(aq)} followed by concentrated \ce{NH3(aq)}.
\item Test 3: Addition of concentrated \ce{H2SO4} and warming.
\end{itemize}
The observations are recorded below:
\begin{center}
\begin{tabular}{|c|c|c|c|}
\hline
\textbf{Test} & \textbf{Tube A} & \textbf{Tube B} & \textbf{Tube C} \\ \hline
1 (\ce{AgNO3} + dil. \ce{NH3}) & Cream ppt, insoluble & White ppt, soluble & Yellow ppt, insoluble \\ \hline
2 (\ce{AgNO3} + conc. \ce{NH3}) & Cream ppt, insoluble & White ppt, soluble & Yellow ppt, insoluble \\ \hline
3 (conc. \ce{H2SO4}, heat) & Brown fumes, purple vapour & White fumes (dense) & White fumes (dense) \\ \hline
\end{tabular}
\end{center}
\begin{enumerate}
\item Identify the halide ion in each tube (A, B, C).
\item Write the ionic equations for the reactions with \ce{AgNO3} for each identified halide.
\item Explain the observations in Test 3 for Tube A, identifying the brown fumes and purple vapour. Why does concentrated \ce{H2SO4} not produce hydrogen bromide or hydrogen iodide gas in a pure state?
\end{enumerate}
\end{exercice}

\begin{exercice}[Group IV: Oxides and Chlorides]
\begin{enumerate}
\item Compare the structures and bonding of \ce{CO2} and \ce{SiO2}. Use diagrams to illustrate your answer. Explain why \ce{CO2} is a gas at room temperature while \ce{SiO2} is a high-melting-point solid.
\item Write balanced equations for the hydrolysis of \ce{SiCl4} and \ce{CCl4} with water. Explain why \ce{SiCl4} hydrolyses violently at room temperature whereas \ce{CCl4} is inert to hydrolysis under the same conditions.
\item Describe the amphoteric nature of \ce{SnO2} and \ce{PbO2}. Write equations for the reaction of \ce{SnO2} with hot concentrated \ce{HCl} and with hot concentrated \ce{NaOH}.
\end{enumerate}
\end{exercice}

\begin{exercice}[Complexes: Naming, Shapes and Isomerism]
\begin{enumerate}
\item Name the following complexes using IUPAC nomenclature and state the oxidation state of the central metal ion, the coordination number, and the likely shape:
      \begin{itemize}
      \item \ce{[Co(NH3)5Cl]Cl2}
      \item \ce{K3[Fe(CN)6]}
      \item \ensuremath{\mathrm{[CuCl_{4}]^{2-}}}
      \end{itemize}
\item Draw the structures (using wedge/dash notation where necessary) and name the two geometric isomers of \ce{[Pt(NH3)2Cl2]}. State which isomer is used as an anti-cancer drug (cisplatin).
\item The complex \ce{[Cr(H2O)5Cl]SO4.H2O} and \ce{[Cr(H2O)5SO4]Cl.H2O} are isomers. Identify the type of isomerism. Describe a simple chemical test (with observations) to distinguish between these two isomers in aqueous solution.
\end{enumerate}
\end{exercice}

\courssub{I prepare for the GCE}

\begin{exercice}[GCE Structured Question: Halogen Chemistry (20 marks)]
\begin{enumerate}
\item \textbf{(a)} Describe and explain the trend in the oxidizing power of the halogens (\ce{F2}, \ce{Cl2}, \ce{Br2}, \ce{I2}) down Group 17. Your answer should refer to electrode potentials, bond dissociation energies, and electron affinities. \hfill [6 marks]
\item \textbf{(b)} A student carries out displacement reactions by adding chlorine water to separate aqueous solutions of potassium bromide and potassium iodide, followed by shaking with an organic solvent (\ce{CCl4} or \ce{CHCl3}). For each reaction:
      \begin{itemize}
      \item State the colour changes observed in the aqueous and organic layers.
      \item Write the balanced ionic equation.
      \item Explain why chlorine displaces bromine and iodine but fluorine is not used for such displacements in the laboratory. \hfill [6 marks]
      \end{itemize}
\item \textbf{(c)} Concentrated sulfuric acid is added to solid sodium iodide and the mixture is heated gently.
      \begin{itemize}
      \item List all the observation(s) made.
      \item Identify all the sulfur-containing and iodine-containing products formed.
      \item Write balanced equations for the redox reactions occurring.
      \item Explain why this reaction is not a suitable method for preparing pure hydrogen iodide. \hfill [8 marks]
      \end{itemize}
\end{enumerate}
\end{exercice}

\begin{exercice}[GCE Essay Question: Group IV and Transition Metals (20 marks)]
\begin{enumerate}
\item \textbf{(a)} The elements of Group IV show a transition from non-metallic to metallic character down the group.
      \begin{itemize}
      \item Discuss the variation in the stability of the +2 and +4 oxidation states from carbon to lead. \hfill [4 marks]
      \item Compare the structures, bonding, and acid/base nature of the dioxides \ce{CO2}, \ce{SiO2}, \ce{GeO2}, \ce{SnO2}, and \ce{PbO2}. \hfill [6 marks]
      \end{itemize}
\item \textbf{(b)} Transition metals exhibit variable oxidation states, form coloured compounds, and act as catalysts.
      \begin{itemize}
      \item Using the first transition series (Sc to Zn), explain the origin of variable oxidation states in terms of electronic configuration. Why do the maximum oxidation states occur around the middle of the series (e.g., Mn)? \hfill [4 marks]
      \item Explain why most transition metal complexes are coloured while \ce{Sc^{3+}} (\ce{d^0}) and \ce{Zn^{2+}} (\ce{d^{10}}) complexes are colourless. Use the concept of d-orbital splitting in an octahedral field. \hfill [4 marks]
      \item Give one example of a homogeneous catalyst and one example of a heterogeneous catalyst involving transition metals or their compounds. State the reaction catalysed in each case. \hfill [2 marks]
      \end{itemize}
\end{enumerate}
\end{exercice}

\begin{exercice}[GCE Data Analysis and Application: Complexes and Spectrochemical Series (20 marks)]
\begin{enumerate}
\item \textbf{(a)} The complex ion \ensuremath{\mathrm{[Cu(H_{2}O)_{6}]^{2+}}} is pale blue. When excess concentrated ammonia is added, the solution turns deep blue due to the formation of \ensuremath{\mathrm{[Cu(NH_{3})_{4}(H_{2}O)_{2}]^{2+}}}.
      \begin{itemize}
      \item Draw the shape of \ensuremath{\mathrm{[Cu(NH_{3})_{4}(H_{2}O)_{2}]^{2+}}} and name its geometry. Is it a regular octahedron? Explain. \hfill [3 marks]
      \item Using Crystal Field Theory, explain the colour change in terms of the spectrochemical series and the magnitude of $\Delta_o$ (crystal field splitting parameter). \hfill [4 marks]
      \end{itemize}
\item \textbf{(b)} Consider the complex \ensuremath{\mathrm{[Co(en)_{3}]^{3+}}} (en = ethylenediamine, \ce{NH2CH2CH2NH2}).
      \begin{itemize}
      \item Draw the two optical isomers (enantiomers) of this complex. \hfill [3 marks]
      \item Explain why this complex does not have geometric isomers. \hfill [2 marks]
      \item Describe how a polarimeter could be used to distinguish between the two enantiomers. \hfill [2 marks]
      \end{itemize}
\item \textbf{(c)} A solution contains a mixture of \ce{[Co(NH3)5SO4]Br} and \ce{[Co(NH3)5Br]SO4} (ionization isomers).
      \begin{itemize}
      \item Describe a chemical test, with expected observations, to distinguish between these two solid isomers. \hfill [3 marks]
      \item Write the ionic equations for the reactions used in your test. \hfill [3 marks]
      \end{itemize}
\end{enumerate}
\end{exercice}

\newpage

\coursec{Solutions to the exercises}

\begin{corrige}[Exercise 1 — MCQ: Halogen Trends]
\textbf{Answer: B.}

Statement B is \textbf{incorrect}. The bond dissociation energy of the X--X bond does \emph{not} decrease steadily down the group. The actual order (increasing bond energy) is:
\[
\ce{I2} < \ce{F2} < \ce{Br2} < \ce{Cl2}
\]
Fluorine has an anomalously \emph{low} bond dissociation energy ($\approx \SI{158}{kJ\,mol^{-1}}$) compared with chlorine ($\approx \SI{242}{kJ\,mol^{-1}}$) and bromine ($\approx \SI{193}{kJ\,mol^{-1}}$). This is because the fluorine atom is very small; the lone pairs on the two adjacent F atoms repel each other strongly, weakening the F--F bond. From \ce{Cl2} to \ce{I2}, the bond energy decreases steadily as the atoms become larger and the orbital overlap poorer (\ce{Cl2} \num{242}, \ce{Br2} \num{193}, \ce{I2} \SI{151}{kJ\,mol^{-1}}).

The other statements are correct:
\begin{itemize}
\item \textbf{A:} Oxidizing power decreases down the group (\ce{F2} is the strongest oxidizing agent, \ce{I2} the weakest).
\item \textbf{C:} Colour intensity increases down the group: \ce{F2} pale yellow, \ce{Cl2} greenish-yellow, \ce{Br2} red-brown, \ce{I2} grey-black solid (violet vapour).
\item \textbf{D:} First ionization energy decreases down the group as the atomic radius increases and the outer electron is further from the nucleus and more shielded.
\end{itemize}
\end{corrige}

\begin{corrige}[Exercise 2 — True/False: Group IV Chemistry]
\begin{enumerate}
\item \textbf{False.} \ce{CCl4} is inert to hydrolysis at room temperature. Although the hydrolysis is thermodynamically favourable, the carbon atom has no available low-energy d-orbitals (it is in period 2) and is sterically shielded by the four chlorine atoms, so water cannot attack it. The activation energy is too high.

\item \textbf{True.} This is the \cle{inert pair effect}. Down the group, the $ns^2$ electrons become increasingly reluctant to participate in bonding, so the stability of the lower (+2) oxidation state increases relative to +4. Thus \ce{Pb^{2+}} compounds are stable while \ce{PbO2} is a strong oxidizing agent; conversely \ce{CO2} and \ce{SiO2} (+4) are very stable.

\item \textbf{True.} \ce{SiO2} is acidic and reacts with hot concentrated alkali:
\[
\ce{SiO2(s) + 2NaOH(aq) -> Na2SiO3(aq) + H2O(l)}
\]

\item \textbf{True.} Because of the inert pair effect, \ce{Pb(IV)} is unstable relative to \ce{Pb(II)}, so \ce{PbO2} readily accepts electrons (it is reduced to \ce{Pb^{2+}}) and is a strong oxidizing agent, e.g.:
\[
\ce{PbO2(s) + 4HCl(aq) -> PbCl2(aq) + Cl2(g) + 2H2O(l)}
\]
\ce{Sn(IV)} is much more stable, so \ce{SnO2} is a weaker oxidizing agent.
\end{enumerate}
\end{corrige}

\begin{corrige}[Exercise 3 — Definitions and Configurations: d-Block]
\begin{enumerate}
\item A \cle{transition element} is an element which forms at least one ion with a \textbf{partially filled d-subshell}. (Note: by this definition Sc and Zn are d-block elements but \emph{not} transition elements, since \ce{Sc^{3+}} is $d^0$ and \ce{Zn^{2+}} is $d^{10}$.)

\item \textbf{Chromium, \ce{Cr} ($Z = 24$):}
\[
\ce{Cr}: \quad 1s^2\,2s^2\,2p^6\,3s^2\,3p^6\,3d^5\,4s^1
\]
The $3d^5\,4s^1$ arrangement (rather than $3d^4\,4s^2$) gives extra stability from the half-filled $3d$ subshell. Electron-in-boxes for the $3d$ orbitals:
\begin{center}
\begin{tabular}{|c|c|c|c|c|}
\hline
$\uparrow$ & $\uparrow$ & $\uparrow$ & $\uparrow$ & $\uparrow$ \\ \hline
\end{tabular}
\quad $3d^5$ (all five orbitals singly occupied, Hund's rule)
\end{center}

\textbf{Copper(II) ion, \ce{Cu^{2+}} ($Z = 29$):} Cu atom is $[\ce{Ar}]\,3d^{10}4s^1$; the $4s$ electron and one $3d$ electron are removed:
\[
\ce{Cu^{2+}}: \quad 1s^2\,2s^2\,2p^6\,3s^2\,3p^6\,3d^9
\]
\begin{center}
\begin{tabular}{|c|c|c|c|c|}
\hline
$\uparrow\downarrow$ & $\uparrow\downarrow$ & $\uparrow\downarrow$ & $\uparrow\downarrow$ & $\uparrow$ \\ \hline
\end{tabular}
\quad $3d^9$ (one unpaired electron)
\end{center}

\item Zinc has configuration $[\ce{Ar}]\,3d^{10}4s^2$ and copper $[\ce{Ar}]\,3d^{10}4s^1$. Zinc has one more proton in its nucleus, so its effective nuclear charge on the $4s$ electrons is greater; in addition, the filled $4s^2$ subshell in Zn is relatively stable. The $4s$ electron of Cu is therefore removed more easily than a $4s$ electron of Zn, so the first ionization energy of Zn is higher than that of Cu.
\end{enumerate}
\end{corrige}

\begin{corrige}[Exercise 4 — Direct Calculations: Redox Titrations]
\begin{enumerate}
\item Half-equations:
\[
\text{Oxidation: } \ce{Fe^{2+}(aq) -> Fe^{3+}(aq) + e-}
\]
\[
\text{Reduction: } \ce{MnO4^{-}(aq) + 8H^{+}(aq) + 5e- -> Mn^{2+}(aq) + 4H2O(l)}
\]

\item Multiplying the iron half-equation by 5 and adding:
\[
\ce{MnO4^{-}(aq) + 8H^{+}(aq) + 5Fe^{2+}(aq) -> Mn^{2+}(aq) + 5Fe^{3+}(aq) + 4H2O(l)}
\]

\item Moles of \ce{MnO4^-}:
\[
n(\ce{MnO4^-}) = \SI{0.0200}{mol\,dm^{-3}} \times \frac{22.5}{1000}\ \mathrm{dm^3} = 4.50\times 10^{-4}\ \mathrm{mol}
\]
From the equation, $n(\ce{Fe^{2+}}) = 5 \times n(\ce{MnO4^-}) = 2.25\times 10^{-3}$ mol in $\SI{25.0}{cm^3}$.
\[
[\ce{FeSO4}] = \frac{2.25\times 10^{-3}}{0.0250} = \boxed{\SI{0.0900}{mol\,dm^{-3}}}
\]

\item Mass concentration:
\[
c_m = 0.0900 \times 278 = \boxed{\SI{25.0}{g\,dm^{-3}}} \quad \text{of } \ce{FeSO4.7H2O}
\]
\end{enumerate}
\end{corrige}

\begin{corrige}[Exercise 5 — Halogen Displacement and Oxoacids]
\begin{enumerate}
\item Chlorine oxidizes bromide ions to bromine:
\[
\ce{Cl2(aq) + 2Br^{-}(aq) -> 2Cl^{-}(aq) + Br2(aq)}
\]
\textbf{Observations:} the aqueous layer turns from colourless to orange/yellow (due to \ce{Br2}); after shaking with \ce{CCl4}, the denser organic layer (lower layer) becomes orange/red-brown, while the aqueous layer becomes paler or colourless, because \ce{Br2} is much more soluble in \ce{CCl4}.

\item Oxoacids of chlorine:
\begin{center}
\begin{tabular}{|c|c|c|}
\hline
\rowcolor{popblueL} \textbf{Oxidation state} & \textbf{Formula} & \textbf{Name} \\ \hline
+1 & \ce{HOCl} (\ce{HClO}) & Chloric(I) acid (hypochlorous acid) \\ \hline
+3 & \ce{HClO2} & Chloric(III) acid (chlorous acid) \\ \hline
+5 & \ce{HClO3} & Chloric(V) acid (chloric acid) \\ \hline
+7 & \ce{HClO4} & Chloric(VII) acid (perchloric acid) \\ \hline
\end{tabular}
\end{center}

\item \textbf{Increasing acid strength:} $\ce{HOCl} < \ce{HClO2} < \ce{HClO3} < \ce{HClO4}$.

\textbf{Increasing thermal stability:} $\ce{HOCl} < \ce{HClO2} < \ce{HClO3} < \ce{HClO4}$ (same order).

\textbf{Explanation of acid strength:} as the number of oxygen atoms bonded to chlorine increases, the conjugate base (\ce{ClO^-}, \ce{ClO2^-}, \ce{ClO3^-}, \ce{ClO4^-}) has its negative charge delocalized over more oxygen atoms. The anion is therefore increasingly stabilized, so the acid dissociates more readily and is stronger. The electron-withdrawing oxygens also weaken the O--H bond, facilitating loss of \ce{H+}.
\end{enumerate}
\end{corrige}

\begin{corrige}[Exercise 6 — Halide Ion Identification]
\begin{enumerate}
\item \textbf{Tube A: bromide, \ce{Br^-}} — cream precipitate of \ce{AgBr}, insoluble in dilute \ce{NH3} (only sparingly soluble even in concentrated), and Test 3 gives brown fumes (\ce{Br2}) with violet/purple vapour traces; the recorded brown fumes confirm bromide. \textbf{Tube B: chloride, \ce{Cl^-}} — white precipitate of \ce{AgCl}, soluble in dilute \ce{NH3}; Test 3 gives dense white fumes of \ce{HCl} only (chloride cannot reduce \ce{H2SO4}). \textbf{Tube C: iodide, \ce{I^-}} — yellow precipitate of \ce{AgI}, insoluble even in concentrated \ce{NH3}; Test 3 gives white fumes (\ce{HI}, \ce{SO2}) plus violet vapour (\ce{I2}).

\item Precipitation equations:
\[
\ce{Ag^{+}(aq) + Br^{-}(aq) -> AgBr(s)} \quad \text{(cream)}
\]
\[
\ce{Ag^{+}(aq) + Cl^{-}(aq) -> AgCl(s)} \quad \text{(white)}
\]
\[
\ce{Ag^{+}(aq) + I^{-}(aq) -> AgI(s)} \quad \text{(yellow)}
\]

\item \textbf{Tube A (NaBr) with conc.\ \ce{H2SO4}:} initially \ce{HBr} forms (misty white fumes):
\[
\ce{NaBr(s) + H2SO4(l) -> NaHSO4(s) + HBr(g)}
\]
But \ce{HBr} is a reducing agent and reduces the sulfuric acid, so the \textbf{brown fumes are bromine, \ce{Br2}}, and \ce{SO2} is also produced:
\[
\ce{2HBr(g) + H2SO4(l) -> Br2(g) + SO2(g) + 2H2O(l)}
\]
(Any violet colouration would come from traces of \ce{I2} impurity; the characteristic observation for bromide is the brown \ce{Br2} fumes.)

Concentrated \ce{H2SO4} is both an acid \emph{and} an oxidizing agent. \ce{HBr} and \ce{HI} are strong reducing agents, so they are oxidized by the hot acid to \ce{Br2} and \ce{I2} respectively (with \ce{SO2}, and for iodide also \ce{S} and \ce{H2S}, as reduction products). Hence pure \ce{HBr} or \ce{HI} gas cannot be obtained this way — phosphoric(V) acid, \ce{H3PO4}, which is non-oxidizing, is used instead.
\end{enumerate}
\end{corrige}

\begin{corrige}[Exercise 7 — Group IV: Oxides and Chlorides]
\begin{enumerate}
\item \textbf{\ce{CO2}} is a simple \emph{molecular} compound: carbon forms two double bonds, \ce{O=C=O}, giving discrete linear molecules held together only by weak van der Waals forces. Little energy is needed to separate the molecules, so \ce{CO2} is a gas at room temperature.

\textbf{\ce{SiO2}} is a \emph{giant covalent (macromolecular)} structure: each Si is bonded tetrahedrally to four O atoms, and each O bridges two Si atoms, forming an extended 3-D network of strong Si--O single bonds (silicon does not form stable $p\pi$--$p\pi$ double bonds with oxygen because its 3p orbitals overlap poorly with oxygen's 2p orbitals). Melting requires breaking strong covalent bonds, so \ce{SiO2} is a solid of very high melting point ($\approx \SI{1700}{\ensuremath{{}^{\circ}\mathrm{C}}}$).

\begin{popfigure}
\begin{center}
\begin{tikzpicture}[scale=0.9]
% CO2 molecule
\node at (0,0) {\ce{O=C=O}};
\node[below, text=popmuted] at (0,-0.4) {discrete linear molecule};
% SiO2 fragment
\begin{scope}[xshift=5.5cm]
\node[circle,draw=popblue,fill=popblueL,inner sep=1.5pt] (Si1) at (0,0) {\scriptsize Si};
\node[circle,draw=poporange,fill=white,inner sep=1pt] (O1) at (-1,0.8) {\scriptsize O};
\node[circle,draw=poporange,fill=white,inner sep=1pt] (O2) at (1,0.8) {\scriptsize O};
\node[circle,draw=poporange,fill=white,inner sep=1pt] (O3) at (-1,-0.8) {\scriptsize O};
\node[circle,draw=poporange,fill=white,inner sep=1pt] (O4) at (1,-0.8) {\scriptsize O};
\node[circle,draw=popblue,fill=popblueL,inner sep=1.5pt] (Si2) at (-2,1.6) {\scriptsize Si};
\node[circle,draw=popblue,fill=popblueL,inner sep=1.5pt] (Si3) at (2,1.6) {\scriptsize Si};
\node[circle,draw=popblue,fill=popblueL,inner sep=1.5pt] (Si4) at (-2,-1.6) {\scriptsize Si};
\node[circle,draw=popblue,fill=popblueL,inner sep=1.5pt] (Si5) at (2,-1.6) {\scriptsize Si};
\draw (Si1)--(O1) (Si1)--(O2) (Si1)--(O3) (Si1)--(O4);
\draw (O1)--(Si2) (O2)--(Si3) (O3)--(Si4) (O4)--(Si5);
\node[below, text=popmuted] at (0,-2.3) {giant covalent network};
\end{scope}
\end{tikzpicture}
\end{center}
\legende{\ce{CO2}: simple molecular; \ce{SiO2}: giant covalent lattice.}
\end{popfigure}

\item \textbf{\ce{SiCl4}} hydrolyses violently:
\[
\ce{SiCl4(l) + 4H2O(l) -> Si(OH)4(s) + 4HCl(aq)} \quad \text{(or } \ce{SiO2.2H2O}\text{)}
\]
\textbf{\ce{CCl4}}: no reaction with water at room temperature (inert).

\textbf{Explanation:} silicon (period 3) has vacant, energetically accessible $3d$ orbitals and a larger, more exposed atom, so a water molecule can donate a lone pair to Si to form an intermediate, allowing nucleophilic attack. Carbon (period 2) has no available d-orbitals, its maximum covalency is 4, and the small C atom is sterically shielded by the four Cl atoms; the activation energy for hydrolysis is therefore prohibitively high.

\item \ce{SnO2} and \ce{PbO2} are \cle{amphoteric} (predominantly basic in character down the group): they react with both acids and alkalis.

With hot concentrated \ce{HCl}:
\[
\ce{SnO2(s) + 4HCl(aq) -> SnCl4(aq) + 2H2O(l)}
\]
With hot concentrated \ce{NaOH}:
\[
\ce{SnO2(s) + 2NaOH(aq) + 2H2O(l) -> Na2[Sn(OH)6](aq)}
\]
(sodium hexahydroxostannate(IV), a stannate). The reaction with acid shows basic character; the reaction with alkali shows acidic character.
\end{enumerate}
\end{corrige}

\begin{corrige}[Exercise 8 — Complexes: Naming, Shapes and Isomerism]
\begin{enumerate}
\item
\begin{itemize}
\item \ce{[Co(NH3)5Cl]Cl2}: \textbf{pentaamminechloridocobalt(III) chloride}. Oxidation state of Co: $+3$; coordination number: 6; shape: \textbf{octahedral}.
\item \ce{K3[Fe(CN)6]}: \textbf{potassium hexacyanidoferrate(III)}. Oxidation state of Fe: $+3$; coordination number: 6; shape: \textbf{octahedral}.
\item \ce{[CuCl4]^{2-}}: \textbf{tetrachloridocuprate(II)} ion. Oxidation state of Cu: $+2$; coordination number: 4; shape: \textbf{tetrahedral} (distorted).
\end{itemize}

\item \ce{[Pt(NH3)2Cl2]} is square planar and exists as two geometric isomers:

\begin{popfigure}
\begin{center}
\begin{tikzpicture}[scale=1.1]
% cis
\node (Pt1) at (0,0) {Pt};
\node (a) at (0,1) {\ce{NH3}};
\node (b) at (1.4,0) {\ce{NH3}};
\node (c) at (0,-1) {Cl};
\node (d) at (-1.4,0) {Cl};
\draw (Pt1)--(a) (Pt1)--(b) (Pt1)--(c) (Pt1)--(d);
\node[below] at (0,-1.6) {\emph{cis} isomer (cisplatin)};
% trans
\begin{scope}[xshift=5.5cm]
\node (Pt2) at (0,0) {Pt};
\node (a2) at (0,1) {\ce{NH3}};
\node (b2) at (1.4,0) {Cl};
\node (c2) at (0,-1) {\ce{NH3}};
\node (d2) at (-1.4,0) {Cl};
\draw (Pt2)--(a2) (Pt2)--(b2) (Pt2)--(c2) (Pt2)--(d2);
\node[below] at (0,-1.6) {\emph{trans} isomer};
\end{scope}
\end{tikzpicture}
\end{center}
\legende{\emph{cis}- and \emph{trans}- \ce{[Pt(NH3)2Cl2]} (square planar).}
\end{popfigure}

In the \emph{cis} isomer the two identical ligands are adjacent ($90°$ apart); in the \emph{trans} isomer they are opposite ($180°$). The \textbf{\emph{cis} isomer is cisplatin}, the anti-cancer drug.

\item These are \textbf{ionization isomers}: they differ in which ion is inside the coordination sphere and which is the counter-ion. \ce{[Cr(H2O)5Cl]SO4.H2O} releases \ce{SO4^{2-}} in solution; \ce{[Cr(H2O)5SO4]Cl.H2O} releases \ce{Cl^-}.

\textbf{Test:} dissolve each compound in water and divide into two portions.
\begin{itemize}
\item Add \ce{BaCl2(aq)} (or \ce{Ba(NO3)2}): the first isomer gives a white precipitate of \ce{BaSO4} ($\ce{Ba^{2+} + SO4^{2-} -> BaSO4(s)}$); the second gives no precipitate.
\item Add \ce{AgNO3(aq)}: the second isomer gives a white precipitate of \ce{AgCl} ($\ce{Ag^{+} + Cl^{-} -> AgCl(s)}$); the first gives no precipitate.
\end{itemize}
\end{enumerate}
\end{corrige}

\begin{corrige}[Exercise 9 — GCE Structured Question: Halogen Chemistry (20 marks)]
\textbf{(a) Trend in oxidizing power [6 marks]}

Oxidizing power \emph{decreases} down the group: $\ce{F2} > \ce{Cl2} > \ce{Br2} > \ce{I2}$. \hfill [1]

Oxidizing ability corresponds to the ease of the process $\ce{X2 + 2e- -> 2X^-}$, measured by the standard electrode potential $E^\circ$, which becomes less positive down the group (e.g.\ $\ce{F2/F^-} \approx +2.87$ V, $\ce{Cl2/Cl^-} \approx +1.36$ V, $\ce{Br2/Br^-} \approx +1.07$ V, $\ce{I2/I^-} \approx +0.54$ V). \hfill [2]

Down the group the atomic radius increases, so the incoming electron is added further from the nucleus and is more shielded; the attraction for the extra electron (effective electron affinity in solution, including hydration) therefore decreases. \hfill [2]

Although the X--X bond dissociation energy also generally decreases down the group (which would favour reaction), the dominant factor is the decreasing tendency to attract and accommodate the extra electron, so oxidizing power decreases. \hfill [1]

\textbf{(b) Displacement reactions [6 marks]}

\emph{Chlorine + potassium bromide:} aqueous layer turns orange/yellow; organic layer (lower layer with \ce{CCl4}) turns orange/red-brown. \hfill [1]
\[
\ce{Cl2(aq) + 2Br^{-}(aq) -> 2Cl^{-}(aq) + Br2(aq)} \tag{1}
\]

\emph{Chlorine + potassium iodide:} aqueous layer turns brown (dark precipitate of \ce{I2} if concentrated); organic layer turns violet/purple. \hfill [1]
\[
\ce{Cl2(aq) + 2I^{-}(aq) -> 2Cl^{-}(aq) + I2(aq)} \tag{1}
\]

Chlorine displaces \ce{Br2} and \ce{I2} because it is a stronger oxidizing agent (higher $E^\circ$) and accepts electrons from \ce{Br^-} and \ce{I^-}. \hfill [1]

Fluorine is not used because it is \emph{too} reactive: it oxidizes water itself, $\ce{2F2 + 2H2O -> 4HF + O2}$, so it cannot exist in aqueous solution, and it is extremely hazardous to handle. \hfill [2]

\textbf{(c) Concentrated \ce{H2SO4} + solid NaI, heated [8 marks]}

\emph{Observations:} misty white fumes (\ce{HI}), violet/purple vapour (\ce{I2}) and a grey/black solid, smell of rotten eggs (\ce{H2S}), yellow solid (\ce{S}), effervescence of a colourless gas that turns acidified dichromate paper green (\ce{SO2}). \hfill [2]

\emph{Products:} sulfur-containing: \ce{SO2}, \ce{S}, \ce{H2S}; iodine-containing: \ce{HI} (little), \ce{I2}. \hfill [2]

\emph{Equations:}
\[
\ce{NaI(s) + H2SO4(l) -> NaHSO4(s) + HI(g)} \tag{1}
\]
\[
\ce{2HI + H2SO4 -> I2 + SO2 + 2H2O}
\]
\[
\ce{6HI + H2SO4 -> 3I2 + S + 4H2O}
\]
\[
\ce{8HI + H2SO4 -> 4I2 + H2S + 4H2O} \tag{2}
\]

\emph{Why unsuitable:} \ce{HI} is a strong reducing agent and is oxidized by the oxidizing acid \ce{H2SO4} to \ce{I2}; the gas is therefore contaminated and the yield of \ce{HI} is very low. Pure \ce{HI} is prepared using non-oxidizing phosphoric(V) acid, \ce{H3PO4}. \hfill [2]
\end{corrige}

\begin{corrige}[Exercise 10 — GCE Essay Question: Group IV and Transition Metals (20 marks)]
\textbf{(a)(i) Stability of +2 and +4 oxidation states [4 marks]}

At the top of the group (C, Si) the +4 state is by far the more stable (e.g.\ \ce{CO2}, \ce{SiO2} very stable; \ce{CO} less so). \hfill [1]

Going down the group the +2 state becomes progressively more stable relative to +4: for Ge the two are comparable; for Sn, \ce{Sn^{2+}} is a mild reducing agent; for Pb, +2 is the dominant stable state and \ce{Pb(IV)} compounds (e.g.\ \ce{PbO2}) are strong oxidizing agents. \hfill [2]

This is the \cle{inert pair effect}: the $ns^2$ electron pair becomes increasingly reluctant to participate in bonding down the group, so only the two $np$ electrons are used, giving the +2 state. \hfill [1]

\textbf{(a)(ii) The dioxides [6 marks]}

\begin{center}
\begin{tabular}{|l|l|l|l|}
\hline
\rowcolor{popblueL} \textbf{Oxide} & \textbf{Structure} & \textbf{Bonding} & \textbf{Acid/base nature} \\ \hline
\ce{CO2} & Simple linear molecules & Covalent (\ce{O=C=O}) & Acidic \\ \hline
\ce{SiO2} & Giant covalent 3-D lattice & Covalent (Si--O) & Acidic (weakly) \\ \hline
\ce{GeO2} & Giant/network, some ionic & Intermediate & Amphoteric (mainly acidic) \\ \hline
\ce{SnO2} & Giant ionic lattice & Mainly ionic & Amphoteric \\ \hline
\ce{PbO2} & Giant ionic lattice & Mainly ionic & Amphoteric (mainly basic) \\ \hline
\end{tabular}
\end{center}

Award: structure trend molecular $\to$ giant covalent $\to$ ionic [2]; bonding trend covalent $\to$ ionic as metallic character increases [2]; acid/base trend acidic $\to$ amphoteric with increasingly basic character down the group [2].

\textbf{(b)(i) Variable oxidation states [4 marks]}

In the first transition series the $3d$ and $4s$ orbitals are very close in energy, so both the $4s$ electrons and varying numbers of $3d$ electrons can be used in bonding/ion formation with little extra energy cost. \hfill [2]

The number of oxidation states increases from Sc to Mn, reaching a maximum at Mn ($+2$ to $+7$), because up to Mn all five $3d$ electrons plus the two $4s$ electrons can be removed/involved. Beyond Mn, the $3d$ electrons are held more tightly (increasing effective nuclear charge, paired d-electrons), so fewer are available and the maximum oxidation state decreases (Fe $+6$ max, commonly $+2/+3$; Zn only $+2$). \hfill [2]

\textbf{(b)(ii) Colour [4 marks]}

In an octahedral ligand field the five $3d$ orbitals split into a lower set ($t_{2g}$) and an upper set ($e_g$) separated by $\Delta_o$. \hfill [1]

In ions with a partially filled d-subshell ($d^1$ to $d^9$), an electron can absorb visible light of energy $h\nu = \Delta_o$ and be promoted from $t_{2g}$ to $e_g$ (a d--d transition); the transmitted/reflected light is the complementary colour, so the complex appears coloured. \hfill [2]

\ce{Sc^{3+}} ($d^0$) has no d-electron to excite, and \ce{Zn^{2+}} ($d^{10}$) has a full subshell with no vacancy in $e_g$; no d--d transition is possible, so their complexes are colourless. \hfill [1]

\textbf{(b)(iii) Catalysts [2 marks]}

\emph{Homogeneous:} \ce{Fe^{2+}/Fe^{3+}} (or \ce{Mn^{2+}}) catalysing the reaction of persulfate with iodide, $\ce{S2O8^{2-} + 2I^- -> 2SO4^{2-} + I2}$ — catalyst and reactants in the same (aqueous) phase. \hfill [1]

\emph{Heterogeneous:} iron in the Haber process ($\ce{N2 + 3H2 <=> 2NH3}$), or \ce{V2O5} in the Contact process, or Ni in the hydrogenation of alkenes — catalyst in a different phase from the reactants. \hfill [1]
\end{corrige}

\begin{corrige}[Exercise 11 — GCE Data Analysis: Complexes and Spectrochemical Series (20 marks)]
\textbf{(a)(i) Shape of \ce{[Cu(NH3)4(H2O)2]^{2+}} [3 marks]}

\begin{popfigure}
\begin{center}
\begin{tikzpicture}[scale=1.1]
\node (Cu) at (0,0) {Cu};
\node (n1) at (-1.5,0) {\ce{NH3}};
\node (n2) at (1.5,0) {\ce{NH3}};
\node (n3) at (-0.6,-1) {\ce{NH3}};
\node (n4) at (0.6,-1) {\ce{NH3}};
\node (w1) at (0,1.4) {\ce{H2O}};
\node (w2) at (0,-1.9) {\ce{H2O}};
\draw (Cu)--(n1) (Cu)--(n2) (Cu)--(n3) (Cu)--(n4);
\draw (Cu)--(w1);
\draw[dashed] (Cu)--(w2);
\end{tikzpicture}
\end{center}
\legende{Distorted octahedron: four \ce{NH3} in a plane, two \ce{H2O} axial (elongated).}
\end{popfigure}

The geometry is \textbf{octahedral but distorted} (tetragonal distortion): it is \emph{not} a regular octahedron because of the \cle{Jahn--Teller effect} for the $d^9$ configuration of \ce{Cu^{2+}} — the two axial bonds (to \ce{H2O}) are longer than the four equatorial bonds. \hfill [3]

\textbf{(a)(ii) Colour change and CFT [4 marks]}

In the spectrochemical series $\ce{I^-} < \ce{Br^-} < \ce{Cl^-} < \ce{H2O} < \ce{NH3} < \text{en} < \ce{CN^-}$, \ce{NH3} is a \emph{stronger-field} ligand than \ce{H2O}. \hfill [1]

Replacing \ce{H2O} by \ce{NH3} therefore \emph{increases} the splitting $\Delta_o$. \hfill [1]

A larger $\Delta_o$ means the d--d transition absorbs light of \emph{higher energy} (shorter wavelength) — absorption shifts from the orange/red region towards the yellow/green. \hfill [1]

The colour observed is the complement of the light absorbed, so the solution changes from pale blue to deep (intense) blue/violet-blue. \hfill [1]

\textbf{(b)(i) Optical isomers of \ce{[Co(en)3]^{3+}} [3 marks]}

\begin{popfigure}
\begin{center}
\begin{tikzpicture}[scale=1.0]
% Lambda
\node (Co1) at (0,0) {Co};
\draw[thick,popblue] (0,0.9) .. controls (-0.9,0.5) .. (-0.5,-0.7);
\draw[thick,popblue] (0,0.9) .. controls (0.9,0.5) .. (0.5,-0.7);
\draw[thick,popblue] (-0.5,-0.7) .. controls (0,-1.1) .. (0.5,-0.7);
\draw[thick,popblue] (-0.85,0.15) .. controls (-1.1,-0.4) .. (-0.5,-0.7);
\draw[thick,popblue] (0.85,0.15) .. controls (1.1,-0.4) .. (0.5,-0.7);
\draw[thick,popblue] (0,0.9) .. controls (-0.5,1.2) .. (-0.85,0.15);
\draw[thick,popblue] (0,0.9) .. controls (0.5,1.2) .. (0.85,0.15);
\node[below] at (0,-1.5) {$\Delta$ (delta)};
% mirror
\draw[dashed,popmuted] (2,-1.5) -- (2,1.5);
\node[above, text=popmuted] at (2,1.5) {\scriptsize mirror};
% Lambda mirror image
\begin{scope}[xshift=4cm]
\node (Co2) at (0,0) {Co};
\draw[thick,poporange] (0,0.9) .. controls (0.9,0.5) .. (0.5,-0.7);
\draw[thick,poporange] (0,0.9) .. controls (-0.9,0.5) .. (-0.5,-0.7);
\draw[thick,poporange] (0.5,-0.7) .. controls (0,-1.1) .. (-0.5,-0.7);
\draw[thick,poporange] (0.85,0.15) .. controls (1.1,-0.4) .. (0.5,-0.7);
\draw[thick,poporange] (-0.85,0.15) .. controls (-1.1,-0.4) .. (-0.5,-0.7);
\draw[thick,poporange] (0,0.9) .. controls (0.5,1.2) .. (0.85,0.15);
\draw[thick,poporange] (0,0.9) .. controls (-0.5,1.2) .. (-0.85,0.15);
\node[below] at (0,-1.5) {$\Lambda$ (lambda)};
\end{scope}
\end{tikzpicture}
\end{center}
\legende{The two non-superimposable mirror images (enantiomers) of \ce{[Co(en)3]^{3+}}; each curved line represents one en ligand.}
\end{popfigure}

The two forms are non-superimposable mirror images, labelled $\Delta$ and $\Lambda$. \hfill [3]

\textbf{(b)(ii) No geometric isomers [2 marks]}

All three ligands are identical bidentate ligands, and each en must span two \emph{adjacent} (cis) positions because its chain is too short to bridge trans positions. With all six coordination sites occupied by three identical symmetrical ligands, there is only one possible connectivity/arrangement; hence no cis/trans (geometric) isomerism is possible — only optical isomerism. \hfill [2]

\textbf{(b)(iii) Polarimetry [2 marks]}

A solution of each enantiomer is placed in turn in a polarimeter tube and plane-polarized light is passed through. One enantiomer rotates the plane of polarization \emph{clockwise} (dextrorotatory, $+$) and the other by an \emph{equal angle anticlockwise} (laevorotatory, $-$). Equal concentrations give equal magnitudes of rotation in opposite directions. \hfill [2]

\textbf{(c)(i) Test to distinguish the ionization isomers [3 marks]}

Dissolve each solid separately in distilled water and test the solution:
\begin{itemize}
\item Add \ce{BaCl2(aq)} (acidified with dilute \ce{HCl}): \ce{[Co(NH3)5Br]SO4} gives a \textbf{white precipitate} of \ce{BaSO4} (free sulfate ions in solution); \ce{[Co(NH3)5SO4]Br} gives no precipitate. \hfill [1.5]
\item Add \ce{AgNO3(aq)}: \ce{[Co(NH3)5SO4]Br} gives a \textbf{cream precipitate} of \ce{AgBr} (free bromide ions); \ce{[Co(NH3)5Br]SO4} gives no precipitate. \hfill [1.5]
\end{itemize}
(The ligand inside the coordination sphere is not released into solution, so it does not react.)

\textbf{(c)(ii) Ionic equations [3 marks]}
\[
\ce{Ba^{2+}(aq) + SO4^{2-}(aq) -> BaSO4(s)} \quad \text{(white)} \tag{1.5}
\]
\[
\ce{Ag^{+}(aq) + Br^{-}(aq) -> AgBr(s)} \quad \text{(cream)} \tag{1.5}
\]
\end{corrige}

\newpage

\coursec{Summary sheet}

\begin{popfigure}
\centering
\begin{tikzpicture}[
  root concept/.append style={concept color=popblue, minimum size=3.2cm, font=\small\bfseries},
  level 1 concept/.append style={concept color=popgreen, minimum size=2.4cm, font=\small, level distance=4.2cm, sibling angle=90},
  mindmap
]
\node[concept, concept color=popblue, minimum size=3.2cm, align=center, font=\small\bfseries] {Organic\\Chemistry}
  child[grow=0]   { node[concept, concept color=popgreen, minimum size=2.4cm, align=center, font=\small] {Alkanes} }
  child[grow=90]  { node[concept, concept color=poporange, minimum size=2.4cm, align=center, font=\small] {Alkenes} }
  child[grow=180] { node[concept, concept color=poppurple, minimum size=2.4cm, align=center, font=\small] {Alcohols} }
  child[grow=270] { node[concept, concept color=popgold, minimum size=2.4cm, align=center, font=\small] {Carboxylic\\acids} };
\end{tikzpicture}
\legende{Mind map of the main families of organic compounds studied in Upper Sixth.}
\end{popfigure}

\findecours

\end{document}
